\section{Limits of functions}
\label{sec:function-limits}
% Original schematic graphs; all coordinates are local LaTeX picture paths.
\newcommand{\rossSketchSign}{%
\begingroup
\setlength{\unitlength}{1mm}
\begin{picture}(120,65)
\linethickness{0.3pt}
\put(7,33){\vector(1,0){107}}
\put(60,5){\vector(0,1){57}}
\put(115,32){\makebox(0,0)[l]{$x$}}
\put(61.5,62){\makebox(0,0)[l]{$y$}}
\put(26.667,32.3){\line(0,1){1.4}}
\put(26.667,30.3){\makebox(0,0)[t]{\scriptsize $-2$}}
\put(93.333,32.3){\line(0,1){1.4}}
\put(93.333,30.3){\makebox(0,0)[t]{\scriptsize $2$}}
\put(59.3,17.375){\line(1,0){1.4}}
\put(58.3,17.375){\makebox(0,0)[r]{\scriptsize $-1$}}
\put(59.3,48.625){\line(1,0){1.4}}
\put(58.3,48.625){\makebox(0,0)[r]{\scriptsize $1$}}
\color{mathjockgreen}
\linethickness{0.6pt}
\qbezier(10,17.375)(35,17.375)(60,17.375)
\qbezier(60,48.625)(85,48.625)(110,48.625)
\color{black}
\put(60,17.375){{\color{white}\circle*{1.8}}}
\put(60,17.375){\circle{1.8}}
\put(60,48.625){{\color{white}\circle*{1.8}}}
\put(60,48.625){\circle{1.8}}
\end{picture}
\endgroup}
\newcommand{\rossSketchSignedSquare}{%
\begingroup
\setlength{\unitlength}{1mm}
\begin{picture}(120,65)
\linethickness{0.3pt}
\put(7,33){\vector(1,0){107}}
\put(60,5){\vector(0,1){57}}
\put(115,32){\makebox(0,0)[l]{$x$}}
\put(61.5,62){\makebox(0,0)[l]{$y$}}
\put(12.381,32.3){\line(0,1){1.4}}
\put(12.381,30.3){\makebox(0,0)[t]{\scriptsize $-2$}}
\put(107.619,32.3){\line(0,1){1.4}}
\put(107.619,30.3){\makebox(0,0)[t]{\scriptsize $2$}}
\put(59.3,10.273){\line(1,0){1.4}}
\put(58.3,10.273){\makebox(0,0)[r]{\scriptsize $-4$}}
\put(59.3,55.727){\line(1,0){1.4}}
\put(58.3,55.727){\makebox(0,0)[r]{\scriptsize $4$}}
\color{mathjockgreen}
\linethickness{0.6pt}
\qbezier(12.381,10.273)(12.976,10.834)(13.571,11.395)
\qbezier(13.571,11.395)(14.167,11.942)(14.762,12.489)
\qbezier(14.762,12.489)(15.357,13.021)(15.952,13.554)
\qbezier(15.952,13.554)(16.548,14.072)(17.143,14.591)
\qbezier(17.143,14.591)(17.738,15.095)(18.333,15.599)
\qbezier(18.333,15.599)(18.929,16.089)(19.524,16.58)
\qbezier(19.524,16.58)(20.119,17.055)(20.714,17.531)
\qbezier(20.714,17.531)(21.31,17.993)(21.905,18.455)
\qbezier(21.905,18.455)(22.5,18.902)(23.095,19.349)
\qbezier(23.095,19.349)(23.69,19.783)(24.286,20.216)
\qbezier(24.286,20.216)(24.881,20.635)(25.476,21.054)
\qbezier(25.476,21.054)(26.071,21.459)(26.667,21.864)
\qbezier(26.667,21.864)(27.262,22.254)(27.857,22.645)
\qbezier(27.857,22.645)(28.452,23.021)(29.048,23.398)
\qbezier(29.048,23.398)(29.643,23.76)(30.238,24.122)
\qbezier(30.238,24.122)(30.833,24.47)(31.429,24.818)
\qbezier(31.429,24.818)(32.024,25.152)(32.619,25.486)
\qbezier(32.619,25.486)(33.214,25.805)(33.81,26.125)
\qbezier(33.81,26.125)(34.405,26.43)(35,26.736)
\qbezier(35,26.736)(35.595,27.027)(36.19,27.318)
\qbezier(36.19,27.318)(36.786,27.595)(37.381,27.872)
\qbezier(37.381,27.872)(37.976,28.135)(38.571,28.398)
\qbezier(38.571,28.398)(39.167,28.646)(39.762,28.895)
\qbezier(39.762,28.895)(40.357,29.129)(40.952,29.364)
\qbezier(40.952,29.364)(41.548,29.584)(42.143,29.804)
\qbezier(42.143,29.804)(42.738,30.01)(43.333,30.216)
\qbezier(43.333,30.216)(43.929,30.408)(44.524,30.599)
\qbezier(44.524,30.599)(45.119,30.777)(45.714,30.955)
\qbezier(45.714,30.955)(46.31,31.118)(46.905,31.281)
\qbezier(46.905,31.281)(47.5,31.43)(48.095,31.58)
\qbezier(48.095,31.58)(48.69,31.714)(49.286,31.849)
\qbezier(49.286,31.849)(49.881,31.97)(50.476,32.091)
\qbezier(50.476,32.091)(51.071,32.197)(51.667,32.304)
\qbezier(51.667,32.304)(52.262,32.396)(52.857,32.489)
\qbezier(52.857,32.489)(53.452,32.567)(54.048,32.645)
\qbezier(54.048,32.645)(54.643,32.709)(55.238,32.773)
\qbezier(55.238,32.773)(55.833,32.822)(56.429,32.872)
\qbezier(56.429,32.872)(57.024,32.908)(57.619,32.943)
\qbezier(57.619,32.943)(58.214,32.964)(58.81,32.986)
\qbezier(58.81,32.986)(59.405,32.993)(60,33)
\qbezier(60,33)(60.595,33.007)(61.19,33.014)
\qbezier(61.19,33.014)(61.786,33.036)(62.381,33.057)
\qbezier(62.381,33.057)(62.976,33.092)(63.571,33.128)
\qbezier(63.571,33.128)(64.167,33.178)(64.762,33.227)
\qbezier(64.762,33.227)(65.357,33.291)(65.952,33.355)
\qbezier(65.952,33.355)(66.548,33.433)(67.143,33.511)
\qbezier(67.143,33.511)(67.738,33.604)(68.333,33.696)
\qbezier(68.333,33.696)(68.929,33.803)(69.524,33.909)
\qbezier(69.524,33.909)(70.119,34.03)(70.714,34.151)
\qbezier(70.714,34.151)(71.31,34.286)(71.905,34.42)
\qbezier(71.905,34.42)(72.5,34.57)(73.095,34.719)
\qbezier(73.095,34.719)(73.69,34.882)(74.286,35.045)
\qbezier(74.286,35.045)(74.881,35.223)(75.476,35.401)
\qbezier(75.476,35.401)(76.071,35.592)(76.667,35.784)
\qbezier(76.667,35.784)(77.262,35.99)(77.857,36.196)
\qbezier(77.857,36.196)(78.452,36.416)(79.048,36.636)
\qbezier(79.048,36.636)(79.643,36.871)(80.238,37.105)
\qbezier(80.238,37.105)(80.833,37.354)(81.429,37.602)
\qbezier(81.429,37.602)(82.024,37.865)(82.619,38.128)
\qbezier(82.619,38.128)(83.214,38.405)(83.81,38.682)
\qbezier(83.81,38.682)(84.405,38.973)(85,39.264)
\qbezier(85,39.264)(85.595,39.57)(86.19,39.875)
\qbezier(86.19,39.875)(86.786,40.195)(87.381,40.514)
\qbezier(87.381,40.514)(87.976,40.848)(88.571,41.182)
\qbezier(88.571,41.182)(89.167,41.53)(89.762,41.878)
\qbezier(89.762,41.878)(90.357,42.24)(90.952,42.602)
\qbezier(90.952,42.602)(91.548,42.979)(92.143,43.355)
\qbezier(92.143,43.355)(92.738,43.746)(93.333,44.136)
\qbezier(93.333,44.136)(93.929,44.541)(94.524,44.946)
\qbezier(94.524,44.946)(95.119,45.365)(95.714,45.784)
\qbezier(95.714,45.784)(96.31,46.217)(96.905,46.651)
\qbezier(96.905,46.651)(97.5,47.098)(98.095,47.545)
\qbezier(98.095,47.545)(98.69,48.007)(99.286,48.469)
\qbezier(99.286,48.469)(99.881,48.945)(100.476,49.42)
\qbezier(100.476,49.42)(101.071,49.911)(101.667,50.401)
\qbezier(101.667,50.401)(102.262,50.905)(102.857,51.409)
\qbezier(102.857,51.409)(103.452,51.928)(104.048,52.446)
\qbezier(104.048,52.446)(104.643,52.979)(105.238,53.511)
\qbezier(105.238,53.511)(105.833,54.058)(106.429,54.605)
\qbezier(106.429,54.605)(107.024,55.166)(107.619,55.727)
\color{black}
\put(60,33){{\color{white}\circle*{1.8}}}
\put(60,33){\circle{1.8}}
\end{picture}
\endgroup}
\newcommand{\rossSketchSinc}{%
\begingroup
\setlength{\unitlength}{1mm}
\begin{picture}(120,65)
\linethickness{0.3pt}
\put(7,19.29){\vector(1,0){107}}
\put(60,5){\vector(0,1){57}}
\put(115,18.29){\makebox(0,0)[l]{$x$}}
\put(61.5,62){\makebox(0,0)[l]{$y$}}
\put(44.292,18.59){\line(0,1){1.4}}
\put(44.292,16.59){\makebox(0,0)[t]{\scriptsize $-\pi$}}
\put(75.708,18.59){\line(0,1){1.4}}
\put(75.708,16.59){\makebox(0,0)[t]{\scriptsize $\pi$}}
\put(59.3,51.548){\line(1,0){1.4}}
\put(58.3,51.548){\makebox(0,0)[r]{\scriptsize $1$}}
\color{mathjockgreen}
\linethickness{0.6pt}
\qbezier(10,17.535)(10.167,17.622)(10.333,17.709)
\qbezier(10.333,17.709)(10.5,17.801)(10.667,17.892)
\qbezier(10.667,17.892)(10.833,17.988)(11,18.084)
\qbezier(11,18.084)(11.167,18.184)(11.333,18.284)
\qbezier(11.333,18.284)(11.5,18.387)(11.666,18.491)
\qbezier(11.666,18.491)(11.833,18.598)(12,18.704)
\qbezier(12,18.704)(12.166,18.814)(12.333,18.924)
\qbezier(12.333,18.924)(12.5,19.036)(12.666,19.147)
\qbezier(12.666,19.147)(12.833,19.261)(13,19.375)
\qbezier(13,19.375)(13.166,19.49)(13.333,19.606)
\qbezier(13.333,19.606)(13.5,19.722)(13.666,19.838)
\qbezier(13.666,19.838)(13.833,19.955)(14,20.072)
\qbezier(14,20.072)(14.166,20.188)(14.333,20.305)
\qbezier(14.333,20.305)(14.5,20.421)(14.666,20.537)
\qbezier(14.666,20.537)(14.833,20.652)(15,20.767)
\qbezier(15,20.767)(15.166,20.881)(15.333,20.994)
\qbezier(15.333,20.994)(15.499,21.105)(15.666,21.217)
\qbezier(15.666,21.217)(15.833,21.325)(15.999,21.434)
\qbezier(15.999,21.434)(16.166,21.54)(16.333,21.645)
\qbezier(16.333,21.645)(16.499,21.747)(16.666,21.849)
\qbezier(16.666,21.849)(16.833,21.947)(16.999,22.045)
\qbezier(16.999,22.045)(17.166,22.138)(17.333,22.231)
\qbezier(17.333,22.231)(17.499,22.319)(17.666,22.407)
\qbezier(17.666,22.407)(17.833,22.489)(17.999,22.572)
\qbezier(17.999,22.572)(18.166,22.648)(18.332,22.725)
\qbezier(18.332,22.725)(18.499,22.795)(18.666,22.865)
\qbezier(18.666,22.865)(18.832,22.928)(18.999,22.991)
\qbezier(18.999,22.991)(19.166,23.047)(19.332,23.102)
\qbezier(19.332,23.102)(19.499,23.151)(19.666,23.199)
\qbezier(19.666,23.199)(19.832,23.239)(19.999,23.279)
\qbezier(19.999,23.279)(20.166,23.311)(20.332,23.343)
\qbezier(20.332,23.343)(20.499,23.367)(20.666,23.39)
\qbezier(20.666,23.39)(20.832,23.405)(20.999,23.42)
\qbezier(20.999,23.42)(21.166,23.426)(21.332,23.431)
\qbezier(21.332,23.431)(21.499,23.428)(21.666,23.424)
\qbezier(21.666,23.424)(21.832,23.412)(21.999,23.399)
\qbezier(21.999,23.399)(22.165,23.377)(22.332,23.354)
\qbezier(22.332,23.354)(22.499,23.323)(22.665,23.291)
\qbezier(22.665,23.291)(22.832,23.25)(22.999,23.208)
\qbezier(22.999,23.208)(23.165,23.158)(23.332,23.107)
\qbezier(23.332,23.107)(23.499,23.046)(23.665,22.986)
\qbezier(23.665,22.986)(23.832,22.916)(23.999,22.847)
\qbezier(23.999,22.847)(24.165,22.768)(24.332,22.689)
\qbezier(24.332,22.689)(24.499,22.601)(24.665,22.513)
\qbezier(24.665,22.513)(24.832,22.416)(24.999,22.319)
\qbezier(24.999,22.319)(25.165,22.213)(25.332,22.108)
\qbezier(25.332,22.108)(25.498,21.994)(25.665,21.88)
\qbezier(25.665,21.88)(25.832,21.758)(25.998,21.636)
\qbezier(25.998,21.636)(26.165,21.506)(26.332,21.376)
\qbezier(26.332,21.376)(26.498,21.239)(26.665,21.102)
\qbezier(26.665,21.102)(26.832,20.958)(26.998,20.815)
\qbezier(26.998,20.815)(27.165,20.664)(27.332,20.514)
\qbezier(27.332,20.514)(27.498,20.358)(27.665,20.202)
\qbezier(27.665,20.202)(27.832,20.041)(27.998,19.88)
\qbezier(27.998,19.88)(28.165,19.714)(28.331,19.547)
\qbezier(28.331,19.547)(28.498,19.377)(28.665,19.207)
\qbezier(28.665,19.207)(28.831,19.034)(28.998,18.86)
\qbezier(28.998,18.86)(29.165,18.684)(29.331,18.507)
\qbezier(29.331,18.507)(29.498,18.329)(29.665,18.15)
\qbezier(29.665,18.15)(29.831,17.97)(29.998,17.79)
\qbezier(29.998,17.79)(30.165,17.61)(30.331,17.429)
\qbezier(30.331,17.429)(30.498,17.248)(30.665,17.068)
\qbezier(30.665,17.068)(30.831,16.888)(30.998,16.709)
\qbezier(30.998,16.709)(31.165,16.531)(31.331,16.352)
\qbezier(31.331,16.352)(31.498,16.177)(31.665,16.001)
\qbezier(31.665,16.001)(31.831,15.829)(31.998,15.656)
\qbezier(31.998,15.656)(32.164,15.488)(32.331,15.319)
\qbezier(32.331,15.319)(32.498,15.156)(32.664,14.992)
\qbezier(32.664,14.992)(32.831,14.834)(32.998,14.676)
\qbezier(32.998,14.676)(33.164,14.525)(33.331,14.373)
\qbezier(33.331,14.373)(33.498,14.229)(33.664,14.084)
\qbezier(33.664,14.084)(33.831,13.948)(33.998,13.812)
\qbezier(33.998,13.812)(34.164,13.684)(34.331,13.557)
\qbezier(34.331,13.557)(34.498,13.439)(34.664,13.321)
\qbezier(34.664,13.321)(34.831,13.213)(34.998,13.105)
\qbezier(34.998,13.105)(35.164,13.009)(35.331,12.912)
\qbezier(35.331,12.912)(35.497,12.827)(35.664,12.742)
\qbezier(35.664,12.742)(35.831,12.669)(35.997,12.597)
\qbezier(35.997,12.597)(36.164,12.537)(36.331,12.478)
\qbezier(36.331,12.478)(36.497,12.432)(36.664,12.386)
\qbezier(36.664,12.386)(36.831,12.354)(36.997,12.322)
\qbezier(36.997,12.322)(37.164,12.305)(37.331,12.288)
\qbezier(37.331,12.288)(37.497,12.287)(37.664,12.285)
\qbezier(37.664,12.285)(37.831,12.299)(37.997,12.313)
\qbezier(37.997,12.313)(38.164,12.344)(38.331,12.374)
\qbezier(38.331,12.374)(38.497,12.421)(38.664,12.468)
\qbezier(38.664,12.468)(38.83,12.531)(38.997,12.595)
\qbezier(38.997,12.595)(39.164,12.676)(39.33,12.757)
\qbezier(39.33,12.757)(39.497,12.855)(39.664,12.953)
\qbezier(39.664,12.953)(39.83,13.069)(39.997,13.185)
\qbezier(39.997,13.185)(40.164,13.318)(40.33,13.452)
\qbezier(40.33,13.452)(40.497,13.603)(40.664,13.755)
\qbezier(40.664,13.755)(40.83,13.924)(40.997,14.093)
\qbezier(40.997,14.093)(41.164,14.28)(41.33,14.467)
\qbezier(41.33,14.467)(41.497,14.672)(41.664,14.876)
\qbezier(41.664,14.876)(41.83,15.098)(41.997,15.321)
\qbezier(41.997,15.321)(42.163,15.56)(42.33,15.8)
\qbezier(42.33,15.8)(42.497,16.057)(42.663,16.313)
\qbezier(42.663,16.313)(42.83,16.587)(42.997,16.86)
\qbezier(42.997,16.86)(43.163,17.15)(43.33,17.44)
\qbezier(43.33,17.44)(43.497,17.746)(43.663,18.052)
\qbezier(43.663,18.052)(43.83,18.374)(43.997,18.695)
\qbezier(43.997,18.695)(44.163,19.032)(44.33,19.368)
\qbezier(44.33,19.368)(44.497,19.719)(44.663,20.07)
\qbezier(44.663,20.07)(44.83,20.435)(44.996,20.8)
\qbezier(44.996,20.8)(45.163,21.178)(45.33,21.556)
\qbezier(45.33,21.556)(45.496,21.946)(45.663,22.337)
\qbezier(45.663,22.337)(45.83,22.739)(45.996,23.141)
\qbezier(45.996,23.141)(46.163,23.554)(46.33,23.967)
\qbezier(46.33,23.967)(46.496,24.39)(46.663,24.812)
\qbezier(46.663,24.812)(46.83,25.244)(46.996,25.676)
\qbezier(46.996,25.676)(47.163,26.117)(47.33,26.557)
\qbezier(47.33,26.557)(47.496,27.004)(47.663,27.452)
\qbezier(47.663,27.452)(47.83,27.905)(47.996,28.359)
\qbezier(47.996,28.359)(48.163,28.818)(48.329,29.276)
\qbezier(48.329,29.276)(48.496,29.739)(48.663,30.202)
\qbezier(48.663,30.202)(48.829,30.668)(48.996,31.134)
\qbezier(48.996,31.134)(49.163,31.602)(49.329,32.07)
\qbezier(49.329,32.07)(49.496,32.539)(49.663,33.008)
\qbezier(49.663,33.008)(49.829,33.477)(49.996,33.945)
\qbezier(49.996,33.945)(50.163,34.412)(50.329,34.88)
\qbezier(50.329,34.88)(50.496,35.344)(50.663,35.809)
\qbezier(50.663,35.809)(50.829,36.27)(50.996,36.731)
\qbezier(50.996,36.731)(51.163,37.188)(51.329,37.644)
\qbezier(51.329,37.644)(51.496,38.095)(51.663,38.545)
\qbezier(51.663,38.545)(51.829,38.989)(51.996,39.432)
\qbezier(51.996,39.432)(52.162,39.867)(52.329,40.302)
\qbezier(52.329,40.302)(52.496,40.729)(52.662,41.155)
\qbezier(52.662,41.155)(52.829,41.57)(52.996,41.986)
\qbezier(52.996,41.986)(53.162,42.39)(53.329,42.795)
\qbezier(53.329,42.795)(53.496,43.186)(53.662,43.578)
\qbezier(53.662,43.578)(53.829,43.957)(53.996,44.335)
\qbezier(53.996,44.335)(54.162,44.699)(54.329,45.063)
\qbezier(54.329,45.063)(54.496,45.412)(54.662,45.761)
\qbezier(54.662,45.761)(54.829,46.093)(54.996,46.426)
\qbezier(54.996,46.426)(55.162,46.741)(55.329,47.056)
\qbezier(55.329,47.056)(55.495,47.354)(55.662,47.651)
\qbezier(55.662,47.651)(55.829,47.93)(55.995,48.209)
\qbezier(55.995,48.209)(56.162,48.468)(56.329,48.727)
\qbezier(56.329,48.727)(56.495,48.966)(56.662,49.205)
\qbezier(56.662,49.205)(56.829,49.423)(56.995,49.642)
\qbezier(56.995,49.642)(57.162,49.839)(57.329,50.035)
\qbezier(57.329,50.035)(57.495,50.21)(57.662,50.386)
\qbezier(57.662,50.386)(57.829,50.538)(57.995,50.691)
\qbezier(57.995,50.691)(58.162,50.821)(58.328,50.951)
\qbezier(58.328,50.951)(58.495,51.058)(58.662,51.165)
\qbezier(58.662,51.165)(58.828,51.248)(58.995,51.332)
\qbezier(58.995,51.332)(59.162,51.392)(59.328,51.451)
\qbezier(59.328,51.451)(59.495,51.488)(59.662,51.524)
\qbezier(59.662,51.524)(59.828,51.536)(59.995,51.548)
\qbezier(60.005,51.548)(60.172,51.536)(60.338,51.524)
\qbezier(60.338,51.524)(60.505,51.488)(60.672,51.451)
\qbezier(60.672,51.451)(60.838,51.392)(61.005,51.332)
\qbezier(61.005,51.332)(61.172,51.248)(61.338,51.165)
\qbezier(61.338,51.165)(61.505,51.058)(61.672,50.951)
\qbezier(61.672,50.951)(61.838,50.821)(62.005,50.691)
\qbezier(62.005,50.691)(62.171,50.538)(62.338,50.386)
\qbezier(62.338,50.386)(62.505,50.21)(62.671,50.035)
\qbezier(62.671,50.035)(62.838,49.839)(63.005,49.642)
\qbezier(63.005,49.642)(63.171,49.423)(63.338,49.205)
\qbezier(63.338,49.205)(63.505,48.966)(63.671,48.727)
\qbezier(63.671,48.727)(63.838,48.468)(64.005,48.209)
\qbezier(64.005,48.209)(64.171,47.93)(64.338,47.651)
\qbezier(64.338,47.651)(64.505,47.354)(64.671,47.056)
\qbezier(64.671,47.056)(64.838,46.741)(65.004,46.426)
\qbezier(65.004,46.426)(65.171,46.093)(65.338,45.761)
\qbezier(65.338,45.761)(65.504,45.412)(65.671,45.063)
\qbezier(65.671,45.063)(65.838,44.699)(66.004,44.335)
\qbezier(66.004,44.335)(66.171,43.957)(66.338,43.578)
\qbezier(66.338,43.578)(66.504,43.186)(66.671,42.795)
\qbezier(66.671,42.795)(66.838,42.39)(67.004,41.986)
\qbezier(67.004,41.986)(67.171,41.57)(67.338,41.155)
\qbezier(67.338,41.155)(67.504,40.729)(67.671,40.302)
\qbezier(67.671,40.302)(67.838,39.867)(68.004,39.432)
\qbezier(68.004,39.432)(68.171,38.989)(68.338,38.545)
\qbezier(68.338,38.545)(68.504,38.095)(68.671,37.644)
\qbezier(68.671,37.644)(68.837,37.188)(69.004,36.731)
\qbezier(69.004,36.731)(69.171,36.27)(69.337,35.809)
\qbezier(69.337,35.809)(69.504,35.344)(69.671,34.88)
\qbezier(69.671,34.88)(69.837,34.412)(70.004,33.945)
\qbezier(70.004,33.945)(70.171,33.477)(70.337,33.008)
\qbezier(70.337,33.008)(70.504,32.539)(70.671,32.07)
\qbezier(70.671,32.07)(70.837,31.602)(71.004,31.134)
\qbezier(71.004,31.134)(71.171,30.668)(71.337,30.202)
\qbezier(71.337,30.202)(71.504,29.739)(71.67,29.276)
\qbezier(71.67,29.276)(71.837,28.818)(72.004,28.359)
\qbezier(72.004,28.359)(72.17,27.905)(72.337,27.452)
\qbezier(72.337,27.452)(72.504,27.004)(72.67,26.557)
\qbezier(72.67,26.557)(72.837,26.117)(73.004,25.676)
\qbezier(73.004,25.676)(73.17,25.244)(73.337,24.812)
\qbezier(73.337,24.812)(73.504,24.39)(73.67,23.967)
\qbezier(73.67,23.967)(73.837,23.554)(74.004,23.141)
\qbezier(74.004,23.141)(74.17,22.739)(74.337,22.337)
\qbezier(74.337,22.337)(74.504,21.946)(74.67,21.556)
\qbezier(74.67,21.556)(74.837,21.178)(75.004,20.8)
\qbezier(75.004,20.8)(75.17,20.435)(75.337,20.07)
\qbezier(75.337,20.07)(75.503,19.719)(75.67,19.368)
\qbezier(75.67,19.368)(75.837,19.032)(76.003,18.695)
\qbezier(76.003,18.695)(76.17,18.374)(76.337,18.052)
\qbezier(76.337,18.052)(76.503,17.746)(76.67,17.44)
\qbezier(76.67,17.44)(76.837,17.15)(77.003,16.86)
\qbezier(77.003,16.86)(77.17,16.587)(77.337,16.313)
\qbezier(77.337,16.313)(77.503,16.057)(77.67,15.8)
\qbezier(77.67,15.8)(77.837,15.56)(78.003,15.321)
\qbezier(78.003,15.321)(78.17,15.098)(78.337,14.876)
\qbezier(78.337,14.876)(78.503,14.672)(78.67,14.467)
\qbezier(78.67,14.467)(78.836,14.28)(79.003,14.093)
\qbezier(79.003,14.093)(79.17,13.924)(79.336,13.755)
\qbezier(79.336,13.755)(79.503,13.603)(79.67,13.452)
\qbezier(79.67,13.452)(79.836,13.318)(80.003,13.185)
\qbezier(80.003,13.185)(80.17,13.069)(80.336,12.953)
\qbezier(80.336,12.953)(80.503,12.855)(80.67,12.757)
\qbezier(80.67,12.757)(80.836,12.676)(81.003,12.595)
\qbezier(81.003,12.595)(81.17,12.531)(81.336,12.468)
\qbezier(81.336,12.468)(81.503,12.421)(81.669,12.374)
\qbezier(81.669,12.374)(81.836,12.344)(82.003,12.313)
\qbezier(82.003,12.313)(82.169,12.299)(82.336,12.285)
\qbezier(82.336,12.285)(82.503,12.287)(82.669,12.288)
\qbezier(82.669,12.288)(82.836,12.305)(83.003,12.322)
\qbezier(83.003,12.322)(83.169,12.354)(83.336,12.386)
\qbezier(83.336,12.386)(83.503,12.432)(83.669,12.478)
\qbezier(83.669,12.478)(83.836,12.537)(84.003,12.597)
\qbezier(84.003,12.597)(84.169,12.669)(84.336,12.742)
\qbezier(84.336,12.742)(84.503,12.827)(84.669,12.912)
\qbezier(84.669,12.912)(84.836,13.009)(85.002,13.105)
\qbezier(85.002,13.105)(85.169,13.213)(85.336,13.321)
\qbezier(85.336,13.321)(85.502,13.439)(85.669,13.557)
\qbezier(85.669,13.557)(85.836,13.684)(86.002,13.812)
\qbezier(86.002,13.812)(86.169,13.948)(86.336,14.084)
\qbezier(86.336,14.084)(86.502,14.229)(86.669,14.373)
\qbezier(86.669,14.373)(86.836,14.525)(87.002,14.676)
\qbezier(87.002,14.676)(87.169,14.834)(87.336,14.992)
\qbezier(87.336,14.992)(87.502,15.156)(87.669,15.319)
\qbezier(87.669,15.319)(87.836,15.488)(88.002,15.656)
\qbezier(88.002,15.656)(88.169,15.829)(88.335,16.001)
\qbezier(88.335,16.001)(88.502,16.177)(88.669,16.352)
\qbezier(88.669,16.352)(88.835,16.531)(89.002,16.709)
\qbezier(89.002,16.709)(89.169,16.888)(89.335,17.068)
\qbezier(89.335,17.068)(89.502,17.248)(89.669,17.429)
\qbezier(89.669,17.429)(89.835,17.61)(90.002,17.79)
\qbezier(90.002,17.79)(90.169,17.97)(90.335,18.15)
\qbezier(90.335,18.15)(90.502,18.329)(90.669,18.507)
\qbezier(90.669,18.507)(90.835,18.684)(91.002,18.86)
\qbezier(91.002,18.86)(91.169,19.034)(91.335,19.207)
\qbezier(91.335,19.207)(91.502,19.377)(91.669,19.547)
\qbezier(91.669,19.547)(91.835,19.714)(92.002,19.88)
\qbezier(92.002,19.88)(92.168,20.041)(92.335,20.202)
\qbezier(92.335,20.202)(92.502,20.358)(92.668,20.514)
\qbezier(92.668,20.514)(92.835,20.664)(93.002,20.815)
\qbezier(93.002,20.815)(93.168,20.958)(93.335,21.102)
\qbezier(93.335,21.102)(93.502,21.239)(93.668,21.376)
\qbezier(93.668,21.376)(93.835,21.506)(94.002,21.636)
\qbezier(94.002,21.636)(94.168,21.758)(94.335,21.88)
\qbezier(94.335,21.88)(94.502,21.994)(94.668,22.108)
\qbezier(94.668,22.108)(94.835,22.213)(95.001,22.319)
\qbezier(95.001,22.319)(95.168,22.416)(95.335,22.513)
\qbezier(95.335,22.513)(95.501,22.601)(95.668,22.689)
\qbezier(95.668,22.689)(95.835,22.768)(96.001,22.847)
\qbezier(96.001,22.847)(96.168,22.916)(96.335,22.986)
\qbezier(96.335,22.986)(96.501,23.046)(96.668,23.107)
\qbezier(96.668,23.107)(96.835,23.158)(97.001,23.208)
\qbezier(97.001,23.208)(97.168,23.25)(97.335,23.291)
\qbezier(97.335,23.291)(97.501,23.323)(97.668,23.354)
\qbezier(97.668,23.354)(97.835,23.377)(98.001,23.399)
\qbezier(98.001,23.399)(98.168,23.412)(98.334,23.424)
\qbezier(98.334,23.424)(98.501,23.428)(98.668,23.431)
\qbezier(98.668,23.431)(98.834,23.426)(99.001,23.42)
\qbezier(99.001,23.42)(99.168,23.405)(99.334,23.39)
\qbezier(99.334,23.39)(99.501,23.367)(99.668,23.343)
\qbezier(99.668,23.343)(99.834,23.311)(100.001,23.279)
\qbezier(100.001,23.279)(100.168,23.239)(100.334,23.199)
\qbezier(100.334,23.199)(100.501,23.151)(100.668,23.102)
\qbezier(100.668,23.102)(100.834,23.047)(101.001,22.991)
\qbezier(101.001,22.991)(101.168,22.928)(101.334,22.865)
\qbezier(101.334,22.865)(101.501,22.795)(101.668,22.725)
\qbezier(101.668,22.725)(101.834,22.648)(102.001,22.572)
\qbezier(102.001,22.572)(102.167,22.489)(102.334,22.407)
\qbezier(102.334,22.407)(102.501,22.319)(102.667,22.231)
\qbezier(102.667,22.231)(102.834,22.138)(103.001,22.045)
\qbezier(103.001,22.045)(103.167,21.947)(103.334,21.849)
\qbezier(103.334,21.849)(103.501,21.747)(103.667,21.645)
\qbezier(103.667,21.645)(103.834,21.54)(104.001,21.434)
\qbezier(104.001,21.434)(104.167,21.325)(104.334,21.217)
\qbezier(104.334,21.217)(104.501,21.105)(104.667,20.994)
\qbezier(104.667,20.994)(104.834,20.881)(105.001,20.767)
\qbezier(105.001,20.767)(105.167,20.652)(105.334,20.537)
\qbezier(105.334,20.537)(105.5,20.421)(105.667,20.305)
\qbezier(105.667,20.305)(105.834,20.188)(106,20.072)
\qbezier(106,20.072)(106.167,19.955)(106.334,19.838)
\qbezier(106.334,19.838)(106.5,19.722)(106.667,19.606)
\qbezier(106.667,19.606)(106.834,19.49)(107,19.375)
\qbezier(107,19.375)(107.167,19.261)(107.334,19.147)
\qbezier(107.334,19.147)(107.5,19.036)(107.667,18.924)
\qbezier(107.667,18.924)(107.834,18.814)(108,18.704)
\qbezier(108,18.704)(108.167,18.598)(108.333,18.491)
\qbezier(108.333,18.491)(108.5,18.387)(108.667,18.284)
\qbezier(108.667,18.284)(108.833,18.184)(109,18.084)
\qbezier(109,18.084)(109.167,17.988)(109.333,17.892)
\qbezier(109.333,17.892)(109.5,17.801)(109.667,17.709)
\qbezier(109.667,17.709)(109.833,17.622)(110,17.535)
\color{black}
\put(60,51.548){{\color{white}\circle*{1.8}}}
\put(60,51.548){\circle{1.8}}
\end{picture}
\endgroup}
\newcommand{\rossSketchOscillatingProduct}{%
\begingroup
\setlength{\unitlength}{1mm}
\begin{picture}(120,65)
\linethickness{0.3pt}
\put(7,24.667){\vector(1,0){107}}
\put(60,5){\vector(0,1){57}}
\put(115,23.667){\makebox(0,0)[l]{$x$}}
\put(61.5,62){\makebox(0,0)[l]{$y$}}
\put(35,23.967){\line(0,1){1.4}}
\put(35,21.967){\makebox(0,0)[t]{\scriptsize $-2$}}
\put(85,23.967){\line(0,1){1.4}}
\put(85,21.967){\makebox(0,0)[t]{\scriptsize $2$}}
\put(59.3,50.308){\line(1,0){1.4}}
\put(58.3,50.308){\makebox(0,0)[r]{\scriptsize $1$}}
\multiput(10,50.308)(4,0){26}{\line(1,0){2}}
\color{mathjockgreen}
\linethickness{0.6pt}
\qbezier(10,50.041)(10.188,50.039)(10.375,50.037)
\qbezier(10.375,50.037)(10.563,50.035)(10.75,50.033)
\qbezier(10.75,50.033)(10.938,50.031)(11.125,50.029)
\qbezier(11.125,50.029)(11.312,50.027)(11.5,50.025)
\qbezier(11.5,50.025)(11.688,50.023)(11.875,50.02)
\qbezier(11.875,50.02)(12.062,50.018)(12.25,50.016)
\qbezier(12.25,50.016)(12.438,50.014)(12.625,50.011)
\qbezier(12.625,50.011)(12.812,50.009)(13,50.006)
\qbezier(13,50.006)(13.187,50.004)(13.375,50.002)
\qbezier(13.375,50.002)(13.562,49.999)(13.75,49.997)
\qbezier(13.75,49.997)(13.938,49.994)(14.125,49.992)
\qbezier(14.125,49.992)(14.312,49.989)(14.5,49.986)
\qbezier(14.5,49.986)(14.687,49.984)(14.875,49.981)
\qbezier(14.875,49.981)(15.062,49.978)(15.25,49.976)
\qbezier(15.25,49.976)(15.437,49.973)(15.625,49.97)
\qbezier(15.625,49.97)(15.812,49.967)(16,49.964)
\qbezier(16,49.964)(16.188,49.961)(16.375,49.958)
\qbezier(16.375,49.958)(16.562,49.955)(16.75,49.952)
\qbezier(16.75,49.952)(16.938,49.949)(17.125,49.946)
\qbezier(17.125,49.946)(17.312,49.943)(17.5,49.94)
\qbezier(17.5,49.94)(17.688,49.936)(17.875,49.933)
\qbezier(17.875,49.933)(18.062,49.93)(18.25,49.926)
\qbezier(18.25,49.926)(18.438,49.923)(18.625,49.919)
\qbezier(18.625,49.919)(18.812,49.916)(19,49.912)
\qbezier(19,49.912)(19.188,49.909)(19.375,49.905)
\qbezier(19.375,49.905)(19.562,49.901)(19.75,49.898)
\qbezier(19.75,49.898)(19.938,49.894)(20.125,49.89)
\qbezier(20.125,49.89)(20.312,49.886)(20.5,49.882)
\qbezier(20.5,49.882)(20.688,49.878)(20.875,49.874)
\qbezier(20.875,49.874)(21.062,49.87)(21.25,49.865)
\qbezier(21.25,49.865)(21.438,49.861)(21.625,49.857)
\qbezier(21.625,49.857)(21.812,49.852)(22,49.848)
\qbezier(22,49.848)(22.188,49.843)(22.375,49.839)
\qbezier(22.375,49.839)(22.562,49.834)(22.75,49.829)
\qbezier(22.75,49.829)(22.938,49.824)(23.125,49.819)
\qbezier(23.125,49.819)(23.312,49.814)(23.5,49.809)
\qbezier(23.5,49.809)(23.688,49.804)(23.875,49.799)
\qbezier(23.875,49.799)(24.062,49.794)(24.25,49.788)
\qbezier(24.25,49.788)(24.438,49.783)(24.625,49.777)
\qbezier(24.625,49.777)(24.812,49.772)(25,49.766)
\qbezier(25,49.766)(25.188,49.76)(25.375,49.754)
\qbezier(25.375,49.754)(25.562,49.748)(25.75,49.742)
\qbezier(25.75,49.742)(25.938,49.736)(26.125,49.73)
\qbezier(26.125,49.73)(26.312,49.723)(26.5,49.717)
\qbezier(26.5,49.717)(26.688,49.71)(26.875,49.703)
\qbezier(26.875,49.703)(27.062,49.697)(27.25,49.69)
\qbezier(27.25,49.69)(27.438,49.683)(27.625,49.675)
\qbezier(27.625,49.675)(27.812,49.668)(28,49.661)
\qbezier(28,49.661)(28.188,49.653)(28.375,49.645)
\qbezier(28.375,49.645)(28.562,49.637)(28.75,49.629)
\qbezier(28.75,49.629)(28.938,49.621)(29.125,49.613)
\qbezier(29.125,49.613)(29.312,49.604)(29.5,49.596)
\qbezier(29.5,49.596)(29.688,49.587)(29.875,49.578)
\qbezier(29.875,49.578)(30.062,49.569)(30.25,49.56)
\qbezier(30.25,49.56)(30.438,49.55)(30.625,49.541)
\qbezier(30.625,49.541)(30.812,49.531)(31,49.521)
\qbezier(31,49.521)(31.188,49.511)(31.375,49.501)
\qbezier(31.375,49.501)(31.562,49.49)(31.75,49.479)
\qbezier(31.75,49.479)(31.938,49.468)(32.125,49.457)
\qbezier(32.125,49.457)(32.312,49.445)(32.5,49.434)
\qbezier(32.5,49.434)(32.688,49.422)(32.875,49.41)
\qbezier(32.875,49.41)(33.062,49.397)(33.25,49.385)
\qbezier(33.25,49.385)(33.438,49.372)(33.625,49.359)
\qbezier(33.625,49.359)(33.812,49.345)(34,49.331)
\qbezier(34,49.331)(34.188,49.317)(34.375,49.303)
\qbezier(34.375,49.303)(34.562,49.288)(34.75,49.273)
\qbezier(34.75,49.273)(34.938,49.258)(35.125,49.242)
\qbezier(35.125,49.242)(35.312,49.226)(35.5,49.21)
\qbezier(35.5,49.21)(35.688,49.193)(35.875,49.176)
\qbezier(35.875,49.176)(36.062,49.158)(36.25,49.14)
\qbezier(36.25,49.14)(36.438,49.122)(36.625,49.103)
\qbezier(36.625,49.103)(36.812,49.083)(37,49.064)
\qbezier(37,49.064)(37.188,49.043)(37.375,49.023)
\qbezier(37.375,49.023)(37.562,49.002)(37.75,48.98)
\qbezier(37.75,48.98)(37.938,48.957)(38.125,48.935)
\qbezier(38.125,48.935)(38.312,48.911)(38.5,48.887)
\qbezier(38.5,48.887)(38.688,48.862)(38.875,48.837)
\qbezier(38.875,48.837)(39.062,48.811)(39.25,48.785)
\qbezier(39.25,48.785)(39.438,48.757)(39.625,48.729)
\qbezier(39.625,48.729)(39.812,48.7)(40,48.671)
\qbezier(40,48.671)(40.188,48.64)(40.375,48.609)
\qbezier(40.375,48.609)(40.562,48.576)(40.75,48.543)
\qbezier(40.75,48.543)(40.938,48.509)(41.125,48.474)
\qbezier(41.125,48.474)(41.312,48.437)(41.5,48.401)
\qbezier(41.5,48.401)(41.688,48.362)(41.875,48.323)
\qbezier(41.875,48.323)(42.062,48.282)(42.25,48.24)
\qbezier(42.25,48.24)(42.438,48.196)(42.625,48.152)
\qbezier(42.625,48.152)(42.812,48.106)(43,48.059)
\qbezier(43,48.059)(43.188,48.009)(43.375,47.959)
\qbezier(43.375,47.959)(43.562,47.906)(43.75,47.853)
\qbezier(43.75,47.853)(43.938,47.796)(44.125,47.739)
\qbezier(44.125,47.739)(44.312,47.678)(44.5,47.617)
\qbezier(44.5,47.617)(44.688,47.552)(44.875,47.487)
\qbezier(44.875,47.487)(45.062,47.417)(45.25,47.347)
\qbezier(45.25,47.347)(45.438,47.272)(45.625,47.196)
\qbezier(45.625,47.196)(45.812,47.115)(46,47.034)
\qbezier(46,47.034)(46.188,46.947)(46.375,46.859)
\qbezier(46.375,46.859)(46.562,46.765)(46.75,46.67)
\qbezier(46.75,46.67)(46.938,46.568)(47.125,46.465)
\qbezier(47.125,46.465)(47.312,46.354)(47.5,46.243)
\qbezier(47.5,46.243)(47.688,46.122)(47.875,46.001)
\qbezier(47.875,46.001)(48.062,45.869)(48.25,45.738)
\qbezier(48.25,45.738)(48.438,45.594)(48.625,45.45)
\qbezier(48.625,45.45)(48.812,45.292)(49,45.135)
\qbezier(49,45.135)(49.188,44.962)(49.375,44.789)
\qbezier(49.375,44.789)(49.562,44.599)(49.75,44.408)
\qbezier(49.75,44.408)(49.938,44.198)(50.125,43.988)
\qbezier(50.125,43.988)(50.312,43.756)(50.5,43.524)
\qbezier(50.5,43.524)(50.688,43.266)(50.875,43.008)
\qbezier(50.875,43.008)(51.062,42.721)(51.25,42.434)
\qbezier(51.25,42.434)(51.438,42.114)(51.625,41.794)
\qbezier(51.625,41.794)(51.812,41.435)(52,41.076)
\qbezier(52,41.076)(52.188,40.674)(52.375,40.271)
\qbezier(52.375,40.271)(52.562,39.817)(52.75,39.364)
\qbezier(52.75,39.364)(52.938,38.852)(53.125,38.34)
\qbezier(53.125,38.34)(53.312,37.761)(53.5,37.181)
\qbezier(53.5,37.181)(53.688,36.525)(53.875,35.868)
\qbezier(53.875,35.868)(54.062,35.124)(54.25,34.381)
\qbezier(54.25,34.381)(54.438,33.539)(54.625,32.698)
\qbezier(54.625,32.698)(54.812,31.751)(55,30.805)
\qbezier(55,30.805)(55.188,29.751)(55.375,28.698)
\qbezier(55.375,28.698)(55.562,27.55)(55.75,26.402)
\qbezier(55.75,26.402)(55.938,25.2)(56.125,23.998)
\qbezier(56.125,23.998)(56.312,22.836)(56.5,21.675)
\qbezier(56.5,21.675)(56.688,20.745)(56.875,19.815)
\qbezier(56.875,19.815)(57.062,19.46)(57.25,19.104)
\qbezier(57.25,19.104)(57.438,19.81)(57.625,20.515)
\qbezier(57.625,20.515)(57.812,22.523)(58,24.531)
\qbezier(58,24.531)(58.188,26.244)(58.375,27.957)
\qbezier(58.375,27.957)(58.562,25.614)(58.75,23.272)
\qbezier(58.75,23.272)(58.75,23.272)(58.75,23.272)
\qbezier(58.75,23.272)(58.773,22.964)(58.795,22.656)
\qbezier(58.795,22.656)(58.816,22.505)(58.837,22.353)
\qbezier(58.837,22.353)(58.857,22.367)(58.876,22.381)
\qbezier(58.876,22.381)(58.895,22.548)(58.913,22.715)
\qbezier(58.913,22.715)(58.93,23.002)(58.947,23.29)
\qbezier(58.947,23.29)(58.963,23.653)(58.98,24.015)
\qbezier(58.98,24.015)(58.995,24.401)(59.01,24.786)
\qbezier(59.01,24.786)(59.024,25.141)(59.038,25.495)
\qbezier(59.038,25.495)(59.052,25.774)(59.065,26.053)
\qbezier(59.065,26.053)(59.078,26.223)(59.091,26.393)
\qbezier(59.091,26.393)(59.103,26.438)(59.115,26.482)
\qbezier(59.115,26.482)(59.126,26.401)(59.138,26.32)
\qbezier(59.138,26.32)(59.149,26.131)(59.16,25.942)
\qbezier(59.16,25.942)(59.17,25.676)(59.18,25.41)
\qbezier(59.18,25.41)(59.19,25.106)(59.2,24.803)
\qbezier(59.2,24.803)(59.209,24.504)(59.219,24.205)
\qbezier(59.219,24.205)(59.228,23.952)(59.237,23.698)
\qbezier(59.237,23.698)(59.245,23.521)(59.254,23.345)
\qbezier(59.254,23.345)(59.262,23.266)(59.27,23.187)
\qbezier(59.27,23.187)(59.278,23.212)(59.286,23.237)
\qbezier(59.286,23.237)(59.293,23.358)(59.301,23.48)
\qbezier(59.301,23.48)(59.308,23.677)(59.315,23.874)
\qbezier(59.315,23.874)(59.322,24.117)(59.329,24.36)
\qbezier(59.329,24.36)(59.335,24.615)(59.342,24.869)
\qbezier(59.342,24.869)(59.348,25.1)(59.355,25.33)
\qbezier(59.355,25.33)(59.361,25.507)(59.367,25.684)
\qbezier(59.367,25.684)(59.373,25.785)(59.379,25.886)
\qbezier(59.379,25.886)(59.385,25.9)(59.39,25.913)
\qbezier(59.39,25.913)(59.396,25.842)(59.401,25.77)
\qbezier(59.401,25.77)(59.406,25.626)(59.412,25.481)
\qbezier(59.412,25.481)(59.417,25.286)(59.422,25.091)
\qbezier(59.422,25.091)(59.427,24.874)(59.432,24.656)
\qbezier(59.432,24.656)(59.437,24.447)(59.441,24.238)
\qbezier(59.441,24.238)(59.446,24.064)(59.451,23.891)
\qbezier(59.451,23.891)(59.455,23.777)(59.459,23.662)
\qbezier(59.459,23.662)(59.464,23.62)(59.468,23.578)
\qbezier(59.468,23.578)(59.472,23.611)(59.476,23.645)
\qbezier(59.476,23.645)(59.48,23.747)(59.485,23.85)
\qbezier(59.485,23.85)(59.488,24.005)(59.492,24.16)
\qbezier(59.492,24.16)(59.496,24.346)(59.5,24.531)
\qbezier(59.5,24.531)(59.504,24.72)(59.507,24.909)
\qbezier(59.507,24.909)(59.511,25.076)(59.515,25.243)
\qbezier(59.515,25.243)(59.518,25.366)(59.522,25.488)
\qbezier(59.522,25.488)(59.525,25.551)(59.528,25.614)
\qbezier(59.528,25.614)(59.532,25.611)(59.535,25.607)
\qbezier(59.535,25.607)(59.538,25.539)(59.541,25.471)
\qbezier(59.541,25.471)(59.544,25.349)(59.548,25.228)
\qbezier(59.548,25.228)(59.551,25.071)(59.554,24.915)
\qbezier(59.554,24.915)(59.557,24.745)(59.559,24.576)
\qbezier(59.559,24.576)(59.562,24.417)(59.565,24.258)
\qbezier(59.565,24.258)(59.568,24.132)(59.571,24.005)
\qbezier(59.571,24.005)(59.574,23.927)(59.576,23.849)
\qbezier(59.576,23.849)(59.579,23.829)(59.582,23.809)
\qbezier(59.582,23.809)(59.584,23.848)(59.587,23.888)
\qbezier(59.587,23.888)(59.589,23.979)(59.592,24.071)
\qbezier(59.592,24.071)(59.594,24.202)(59.597,24.332)
\qbezier(59.597,24.332)(59.599,24.483)(59.602,24.633)
\qbezier(59.602,24.633)(59.604,24.782)(59.606,24.931)
\qbezier(59.606,24.931)(59.609,25.059)(59.611,25.186)
\qbezier(59.611,25.186)(59.613,25.275)(59.615,25.364)
\qbezier(59.615,25.364)(59.618,25.403)(59.62,25.442)
\qbezier(59.62,25.442)(59.622,25.427)(59.624,25.411)
\qbezier(59.624,25.411)(59.626,25.345)(59.628,25.279)
\qbezier(59.628,25.279)(59.63,25.172)(59.632,25.066)
\qbezier(59.632,25.066)(59.634,24.934)(59.636,24.802)
\qbezier(59.636,24.802)(59.638,24.664)(59.64,24.526)
\qbezier(59.64,24.526)(59.642,24.4)(59.644,24.274)
\qbezier(59.644,24.274)(59.646,24.178)(59.648,24.082)
\qbezier(59.648,24.082)(59.65,24.029)(59.652,23.975)
\qbezier(59.652,23.975)(59.653,23.97)(59.655,23.965)
\qbezier(59.655,23.965)(59.657,24.008)(59.659,24.051)
\qbezier(59.659,24.051)(59.66,24.136)(59.662,24.221)
\qbezier(59.662,24.221)(59.664,24.334)(59.666,24.448)
\qbezier(59.666,24.448)(59.667,24.575)(59.669,24.701)
\qbezier(59.669,24.701)(59.671,24.823)(59.672,24.945)
\qbezier(59.672,24.945)(59.674,25.045)(59.675,25.145)
\qbezier(59.675,25.145)(59.677,25.21)(59.678,25.275)
\qbezier(59.678,25.275)(59.68,25.298)(59.682,25.32)
\qbezier(59.682,25.32)(59.683,25.297)(59.685,25.273)
\qbezier(59.685,25.273)(59.686,25.209)(59.688,25.144)
\qbezier(59.688,25.144)(59.689,25.048)(59.69,24.952)
\qbezier(59.69,24.952)(59.692,24.838)(59.693,24.724)
\qbezier(59.693,24.724)(59.695,24.608)(59.696,24.492)
\qbezier(59.696,24.492)(59.697,24.39)(59.699,24.288)
\qbezier(59.699,24.288)(59.7,24.214)(59.701,24.14)
\qbezier(59.701,24.14)(59.703,24.104)(59.704,24.068)
\qbezier(59.704,24.068)(59.705,24.073)(59.707,24.079)
\qbezier(59.707,24.079)(59.708,24.125)(59.709,24.171)
\qbezier(59.709,24.171)(59.711,24.25)(59.712,24.329)
\qbezier(59.712,24.329)(59.713,24.431)(59.714,24.532)
\qbezier(59.714,24.532)(59.715,24.641)(59.717,24.749)
\qbezier(59.717,24.749)(59.718,24.851)(59.719,24.952)
\qbezier(59.719,24.952)(59.72,25.032)(59.721,25.112)
\qbezier(59.721,25.112)(59.723,25.16)(59.724,25.207)
\qbezier(59.724,25.207)(59.725,25.217)(59.726,25.227)
\qbezier(59.726,25.227)(59.727,25.198)(59.728,25.169)
\qbezier(59.728,25.169)(59.729,25.106)(59.73,25.043)
\qbezier(59.73,25.043)(59.732,24.955)(59.733,24.867)
\qbezier(59.733,24.867)(59.734,24.767)(59.735,24.666)
\qbezier(59.735,24.666)(59.736,24.567)(59.737,24.468)
\qbezier(59.737,24.468)(59.738,24.385)(59.739,24.301)
\qbezier(59.739,24.301)(59.74,24.244)(59.741,24.187)
\qbezier(59.741,24.187)(59.742,24.164)(59.743,24.141)
\qbezier(59.743,24.141)(59.744,24.154)(59.745,24.168)
\qbezier(59.745,24.168)(59.746,24.215)(59.747,24.263)
\qbezier(59.747,24.263)(59.748,24.338)(59.749,24.413)
\qbezier(59.749,24.413)(59.75,24.504)(59.751,24.595)
\qbezier(59.751,24.595)(59.752,24.69)(59.752,24.785)
\qbezier(59.752,24.785)(59.753,24.87)(59.754,24.955)
\qbezier(59.754,24.955)(59.755,25.019)(59.756,25.083)
\qbezier(59.756,25.083)(59.757,25.118)(59.758,25.152)
\qbezier(59.758,25.152)(59.759,25.153)(59.76,25.153)
\qbezier(59.76,25.153)(59.76,25.12)(59.761,25.087)
\qbezier(59.761,25.087)(59.762,25.025)(59.763,24.964)
\qbezier(59.763,24.964)(59.764,24.882)(59.765,24.801)
\qbezier(59.765,24.801)(59.766,24.712)(59.766,24.622)
\qbezier(59.766,24.622)(59.767,24.537)(59.768,24.452)
\qbezier(59.768,24.452)(59.769,24.383)(59.77,24.314)
\qbezier(59.77,24.314)(59.77,24.27)(59.771,24.226)
\qbezier(59.771,24.226)(59.772,24.213)(59.773,24.201)
\qbezier(59.773,24.201)(59.773,24.22)(59.774,24.24)
\qbezier(59.774,24.24)(59.775,24.288)(59.776,24.337)
\qbezier(59.776,24.337)(59.777,24.408)(59.777,24.479)
\qbezier(59.777,24.479)(59.778,24.562)(59.779,24.645)
\qbezier(59.779,24.645)(59.779,24.728)(59.78,24.811)
\qbezier(59.78,24.811)(59.781,24.883)(59.782,24.956)
\qbezier(59.782,24.956)(59.782,25.007)(59.783,25.058)
\qbezier(59.783,25.058)(59.784,25.082)(59.784,25.106)
\qbezier(59.784,25.106)(59.785,25.099)(59.786,25.092)
\qbezier(59.786,25.092)(59.787,25.055)(59.787,25.019)
\qbezier(59.787,25.019)(59.788,24.959)(59.789,24.899)
\qbezier(59.789,24.899)(59.789,24.824)(59.79,24.748)
\qbezier(59.79,24.748)(59.791,24.668)(59.791,24.588)
\qbezier(59.791,24.588)(59.792,24.514)(59.793,24.441)
\qbezier(59.793,24.441)(59.793,24.384)(59.794,24.326)
\qbezier(59.794,24.326)(59.794,24.293)(59.795,24.261)
\qbezier(59.795,24.261)(59.796,24.256)(59.796,24.252)
\qbezier(59.796,24.252)(59.797,24.276)(59.798,24.3)
\qbezier(59.798,24.3)(59.798,24.349)(59.799,24.399)
\qbezier(59.799,24.399)(59.799,24.466)(59.8,24.533)
\qbezier(59.8,24.533)(59.801,24.609)(59.801,24.684)
\qbezier(59.801,24.684)(59.802,24.758)(59.802,24.831)
\qbezier(59.802,24.831)(59.803,24.893)(59.804,24.954)
\qbezier(59.804,24.954)(59.804,24.995)(59.805,25.035)
\qbezier(59.805,25.035)(59.805,25.05)(59.806,25.065)
\qbezier(59.806,25.065)(59.806,25.052)(59.807,25.039)
\qbezier(59.807,25.039)(59.808,25.001)(59.808,24.962)
\qbezier(59.808,24.962)(59.809,24.904)(59.809,24.845)
\qbezier(59.809,24.845)(59.81,24.775)(59.81,24.705)
\qbezier(59.81,24.705)(59.811,24.633)(59.811,24.561)
\qbezier(59.811,24.561)(59.812,24.497)(59.812,24.433)
\qbezier(59.812,24.433)(59.813,24.386)(59.813,24.339)
\qbezier(59.813,24.339)(59.814,24.315)(59.814,24.292)
\qbezier(59.814,24.292)(59.815,24.294)(59.815,24.296)
\qbezier(59.815,24.296)(59.816,24.324)(59.817,24.352)
\qbezier(59.817,24.352)(59.817,24.401)(59.818,24.451)
\qbezier(59.818,24.451)(59.818,24.514)(59.819,24.578)
\qbezier(59.819,24.578)(59.819,24.647)(59.819,24.716)
\qbezier(59.819,24.716)(59.82,24.781)(59.82,24.846)
\qbezier(59.82,24.846)(59.821,24.898)(59.821,24.95)
\qbezier(59.821,24.95)(59.822,24.982)(59.822,25.014)
\qbezier(59.822,25.014)(59.823,25.021)(59.823,25.028)
\qbezier(59.823,25.028)(59.824,25.011)(59.824,24.993)
\qbezier(59.824,24.993)(59.825,24.953)(59.825,24.913)
\qbezier(59.825,24.913)(59.826,24.856)(59.826,24.8)
\qbezier(59.826,24.8)(59.827,24.734)(59.827,24.669)
\qbezier(59.827,24.669)(59.827,24.604)(59.828,24.54)
\qbezier(59.828,24.54)(59.828,24.484)(59.829,24.429)
\qbezier(59.829,24.429)(59.829,24.391)(59.83,24.352)
\qbezier(59.83,24.352)(59.83,24.336)(59.831,24.32)
\qbezier(59.831,24.32)(59.831,24.328)(59.831,24.336)
\qbezier(59.831,24.336)(59.832,24.367)(59.832,24.397)
\qbezier(59.832,24.397)(59.833,24.446)(59.833,24.496)
\qbezier(59.833,24.496)(59.833,24.556)(59.834,24.616)
\qbezier(59.834,24.616)(59.834,24.68)(59.835,24.743)
\qbezier(59.835,24.743)(59.835,24.8)(59.836,24.858)
\qbezier(59.836,24.858)(59.836,24.901)(59.836,24.945)
\qbezier(59.836,24.945)(59.837,24.969)(59.837,24.993)
\qbezier(59.837,24.993)(59.838,24.994)(59.838,24.995)
\qbezier(59.838,24.995)(59.838,24.974)(59.839,24.952)
\qbezier(59.839,24.952)(59.839,24.911)(59.839,24.87)
\qbezier(59.839,24.87)(59.84,24.815)(59.84,24.76)
\qbezier(59.84,24.76)(59.841,24.7)(59.841,24.639)
\qbezier(59.841,24.639)(59.841,24.581)(59.842,24.523)
\qbezier(59.842,24.523)(59.842,24.475)(59.843,24.427)
\qbezier(59.843,24.427)(59.843,24.396)(59.843,24.366)
\qbezier(59.843,24.366)(59.844,24.356)(59.844,24.347)
\qbezier(59.844,24.347)(59.844,24.359)(59.845,24.372)
\qbezier(59.845,24.372)(59.845,24.405)(59.845,24.438)
\qbezier(59.845,24.438)(59.846,24.486)(59.846,24.535)
\qbezier(59.846,24.535)(59.847,24.592)(59.847,24.649)
\qbezier(59.847,24.649)(59.847,24.707)(59.848,24.765)
\qbezier(59.848,24.765)(59.848,24.815)(59.848,24.866)
\qbezier(59.848,24.866)(59.849,24.902)(59.849,24.938)
\qbezier(59.849,24.938)(59.849,24.955)(59.85,24.973)
\qbezier(59.85,24.973)(59.85,24.969)(59.85,24.964)
\qbezier(59.85,24.964)(59.851,24.94)(59.851,24.915)
\qbezier(59.851,24.915)(59.851,24.873)(59.852,24.832)
\qbezier(59.852,24.832)(59.852,24.779)(59.852,24.726)
\qbezier(59.852,24.726)(59.853,24.67)(59.853,24.614)
\qbezier(59.853,24.614)(59.853,24.561)(59.854,24.509)
\qbezier(59.854,24.509)(59.854,24.468)(59.854,24.428)
\qbezier(59.854,24.428)(59.855,24.404)(59.855,24.38)
\qbezier(59.855,24.38)(59.855,24.376)(59.855,24.372)
\qbezier(59.855,24.372)(59.856,24.388)(59.856,24.405)
\qbezier(59.856,24.405)(59.856,24.44)(59.857,24.474)
\qbezier(59.857,24.474)(59.857,24.522)(59.857,24.569)
\qbezier(59.857,24.569)(59.858,24.623)(59.858,24.677)
\qbezier(59.858,24.677)(59.858,24.73)(59.859,24.782)
\qbezier(59.859,24.782)(59.859,24.827)(59.859,24.871)
\qbezier(59.859,24.871)(59.859,24.901)(59.86,24.93)
\qbezier(59.86,24.93)(59.86,24.942)(59.86,24.953)
\qbezier(59.86,24.953)(59.861,24.944)(59.861,24.936)
\qbezier(59.861,24.936)(59.861,24.908)(59.861,24.881)
\qbezier(59.861,24.881)(59.862,24.839)(59.862,24.798)
\qbezier(59.862,24.798)(59.862,24.747)(59.863,24.697)
\qbezier(59.863,24.697)(59.863,24.644)(59.863,24.592)
\qbezier(59.863,24.592)(59.863,24.546)(59.864,24.499)
\qbezier(59.864,24.499)(59.864,24.464)(59.864,24.43)
\qbezier(59.864,24.43)(59.865,24.412)(59.865,24.394)
\qbezier(59.865,24.394)(59.865,24.395)(59.865,24.396)
\qbezier(59.865,24.396)(59.866,24.416)(59.866,24.436)
\qbezier(59.866,24.436)(59.866,24.471)(59.866,24.507)
\qbezier(59.866,24.507)(59.867,24.553)(59.867,24.6)
\qbezier(59.867,24.6)(59.867,24.65)(59.868,24.701)
\qbezier(59.868,24.701)(59.868,24.749)(59.868,24.797)
\qbezier(59.868,24.797)(59.868,24.835)(59.869,24.874)
\qbezier(59.869,24.874)(59.869,24.898)(59.869,24.922)
\qbezier(59.869,24.922)(59.869,24.928)(59.87,24.934)
\qbezier(59.87,24.934)(59.87,24.921)(59.87,24.909)
\qbezier(59.87,24.909)(59.87,24.879)(59.871,24.85)
\qbezier(59.871,24.85)(59.871,24.809)(59.871,24.767)
\qbezier(59.871,24.767)(59.871,24.719)(59.872,24.67)
\qbezier(59.872,24.67)(59.872,24.622)(59.872,24.574)
\qbezier(59.872,24.574)(59.872,24.533)(59.873,24.491)
\qbezier(59.873,24.491)(59.873,24.462)(59.873,24.433)
\qbezier(59.873,24.433)(59.873,24.421)(59.874,24.408)
\qbezier(59.874,24.408)(59.874,24.414)(59.874,24.419)
\qbezier(59.874,24.419)(59.874,24.442)(59.875,24.464)
\qbezier(59.875,24.464)(59.875,24.5)(59.875,24.537)
\qbezier(59.875,24.537)(59.875,24.582)(59.875,24.627)
\qbezier(59.875,24.627)(59.876,24.674)(59.876,24.722)
\qbezier(59.876,24.722)(59.876,24.765)(59.876,24.809)
\qbezier(59.876,24.809)(59.877,24.842)(59.877,24.875)
\qbezier(59.877,24.875)(59.877,24.893)(59.877,24.912)
\qbezier(59.877,24.912)(59.878,24.913)(59.878,24.915)
\qbezier(59.878,24.915)(59.878,24.899)(59.878,24.883)
\qbezier(59.878,24.883)(59.878,24.852)(59.879,24.822)
\qbezier(59.879,24.822)(59.879,24.781)(59.879,24.739)
\qbezier(59.879,24.739)(59.879,24.693)(59.88,24.648)
\qbezier(59.88,24.648)(59.88,24.603)(59.88,24.559)
\qbezier(59.88,24.559)(59.88,24.523)(59.88,24.486)
\qbezier(59.88,24.486)(59.881,24.462)(59.881,24.439)
\qbezier(59.881,24.439)(59.881,24.431)(59.881,24.423)
\qbezier(59.881,24.423)(59.881,24.432)(59.882,24.441)
\qbezier(59.882,24.441)(59.882,24.466)(59.882,24.491)
\qbezier(59.882,24.491)(59.882,24.527)(59.882,24.564)
\qbezier(59.882,24.564)(59.883,24.608)(59.883,24.651)
\qbezier(59.883,24.651)(59.883,24.696)(59.883,24.74)
\qbezier(59.883,24.74)(59.884,24.779)(59.884,24.818)
\qbezier(59.884,24.818)(59.884,24.846)(59.884,24.874)
\qbezier(59.884,24.874)(59.884,24.888)(59.885,24.901)
\qbezier(59.885,24.901)(59.885,24.899)(59.885,24.896)
\qbezier(59.885,24.896)(59.885,24.877)(59.885,24.859)
\qbezier(59.885,24.859)(59.886,24.827)(59.886,24.795)
\qbezier(59.886,24.795)(59.886,24.755)(59.886,24.714)
\qbezier(59.886,24.714)(59.886,24.671)(59.886,24.627)
\qbezier(59.886,24.627)(59.887,24.587)(59.887,24.547)
\qbezier(59.887,24.547)(59.887,24.515)(59.887,24.483)
\qbezier(59.887,24.483)(59.887,24.464)(59.888,24.445)
\qbezier(59.888,24.445)(59.888,24.442)(59.888,24.438)
\qbezier(59.888,24.438)(59.888,24.45)(59.888,24.463)
\qbezier(59.888,24.463)(59.889,24.489)(59.889,24.516)
\qbezier(59.889,24.516)(59.889,24.553)(59.889,24.589)
\qbezier(59.889,24.589)(59.889,24.631)(59.89,24.673)
\qbezier(59.89,24.673)(59.89,24.714)(59.89,24.755)
\qbezier(59.89,24.755)(59.89,24.79)(59.89,24.825)
\qbezier(59.89,24.825)(59.89,24.848)(59.891,24.872)
\qbezier(59.891,24.872)(59.891,24.881)(59.891,24.89)
\qbezier(59.891,24.89)(59.891,24.884)(59.891,24.877)
\qbezier(59.891,24.877)(59.891,24.857)(59.892,24.836)
\qbezier(59.892,24.836)(59.892,24.803)(59.892,24.771)
\qbezier(59.892,24.771)(59.892,24.731)(59.892,24.692)
\qbezier(59.892,24.692)(59.893,24.651)(59.893,24.61)
\qbezier(59.893,24.61)(59.893,24.573)(59.893,24.536)
\qbezier(59.893,24.536)(59.893,24.509)(59.893,24.481)
\qbezier(59.893,24.481)(59.894,24.467)(59.894,24.452)
\qbezier(59.894,24.452)(59.894,24.453)(59.894,24.453)
\qbezier(59.894,24.453)(59.894,24.469)(59.894,24.484)
\qbezier(59.894,24.484)(59.895,24.512)(59.895,24.539)
\qbezier(59.895,24.539)(59.895,24.576)(59.895,24.612)
\qbezier(59.895,24.612)(59.895,24.652)(59.895,24.692)
\qbezier(59.895,24.692)(59.896,24.73)(59.896,24.768)
\qbezier(59.896,24.768)(59.896,24.799)(59.896,24.83)
\qbezier(59.896,24.83)(59.896,24.849)(59.896,24.868)
\qbezier(59.896,24.868)(59.897,24.873)(59.897,24.878)
\qbezier(59.897,24.878)(59.897,24.869)(59.897,24.859)
\qbezier(59.897,24.859)(59.897,24.836)(59.897,24.814)
\qbezier(59.897,24.814)(59.897,24.781)(59.898,24.748)
\qbezier(59.898,24.748)(59.898,24.71)(59.898,24.671)
\qbezier(59.898,24.671)(59.898,24.633)(59.898,24.595)
\qbezier(59.898,24.595)(59.898,24.561)(59.899,24.528)
\qbezier(59.899,24.528)(59.899,24.505)(59.899,24.482)
\qbezier(59.899,24.482)(59.899,24.471)(59.899,24.461)
\qbezier(59.899,24.461)(59.899,24.465)(59.899,24.469)
\qbezier(59.899,24.469)(59.9,24.486)(59.9,24.504)
\qbezier(59.9,24.504)(59.9,24.533)(59.9,24.562)
\qbezier(59.9,24.562)(59.9,24.597)(59.9,24.633)
\qbezier(59.9,24.633)(59.901,24.671)(59.901,24.709)
\qbezier(59.901,24.709)(59.901,24.744)(59.901,24.779)
\qbezier(59.901,24.779)(59.901,24.806)(59.901,24.833)
\qbezier(59.901,24.833)(59.901,24.848)(59.902,24.863)
\qbezier(59.902,24.863)(59.902,24.865)(59.902,24.866)
\qbezier(59.902,24.866)(59.902,24.854)(59.902,24.841)
\qbezier(59.902,24.841)(59.902,24.817)(59.902,24.793)
\qbezier(59.902,24.793)(59.903,24.76)(59.903,24.727)
\qbezier(59.903,24.727)(59.903,24.69)(59.903,24.653)
\qbezier(59.903,24.653)(59.903,24.617)(59.903,24.581)
\qbezier(59.903,24.581)(59.903,24.552)(59.904,24.522)
\qbezier(59.904,24.522)(59.904,24.503)(59.904,24.483)
\qbezier(59.904,24.483)(59.904,24.477)(59.904,24.47)
\qbezier(59.904,24.47)(59.904,24.477)(59.904,24.484)
\qbezier(59.904,24.484)(59.905,24.504)(59.905,24.523)
\qbezier(59.905,24.523)(59.905,24.553)(59.905,24.582)
\qbezier(59.905,24.582)(59.905,24.618)(59.905,24.653)
\qbezier(59.905,24.653)(59.905,24.689)(59.905,24.725)
\qbezier(59.905,24.725)(59.906,24.756)(59.906,24.788)
\qbezier(59.906,24.788)(59.906,24.811)(59.906,24.834)
\qbezier(59.906,24.834)(59.906,24.846)(59.906,24.857)
\qbezier(59.906,24.857)(59.906,24.855)(59.907,24.853)
\qbezier(59.907,24.853)(59.907,24.838)(59.907,24.824)
\qbezier(59.907,24.824)(59.907,24.798)(59.907,24.772)
\qbezier(59.907,24.772)(59.907,24.74)(59.907,24.707)
\qbezier(59.907,24.707)(59.907,24.672)(59.908,24.636)
\qbezier(59.908,24.636)(59.908,24.603)(59.908,24.57)
\qbezier(59.908,24.57)(59.908,24.544)(59.908,24.518)
\qbezier(59.908,24.518)(59.908,24.502)(59.908,24.486)
\qbezier(59.908,24.486)(59.908,24.483)(59.909,24.48)
\qbezier(59.909,24.48)(59.909,24.49)(59.909,24.5)
\qbezier(59.909,24.5)(59.909,24.521)(59.909,24.542)
\qbezier(59.909,24.542)(59.909,24.572)(59.909,24.602)
\qbezier(59.909,24.602)(59.909,24.636)(59.91,24.67)
\qbezier(59.91,24.67)(59.91,24.704)(59.91,24.738)
\qbezier(59.91,24.738)(59.91,24.766)(59.91,24.795)
\qbezier(59.91,24.795)(59.91,24.815)(59.91,24.834)
\qbezier(59.91,24.834)(59.91,24.842)(59.911,24.85)
\qbezier(59.911,24.85)(59.911,24.845)(59.911,24.84)
\qbezier(59.911,24.84)(59.911,24.823)(59.911,24.806)
\qbezier(59.911,24.806)(59.911,24.78)(59.911,24.753)
\qbezier(59.911,24.753)(59.911,24.721)(59.912,24.689)
\qbezier(59.912,24.689)(59.912,24.655)(59.912,24.621)
\qbezier(59.912,24.621)(59.912,24.591)(59.912,24.561)
\qbezier(59.912,24.561)(59.912,24.538)(59.912,24.515)
\qbezier(59.912,24.515)(59.912,24.503)(59.912,24.49)
\qbezier(59.912,24.49)(59.913,24.49)(59.913,24.49)
\qbezier(59.913,24.49)(59.913,24.503)(59.913,24.515)
\qbezier(59.913,24.515)(59.913,24.538)(59.913,24.56)
\qbezier(59.913,24.56)(59.913,24.59)(59.913,24.62)
\qbezier(59.913,24.62)(59.913,24.653)(59.914,24.686)
\qbezier(59.914,24.686)(59.914,24.718)(59.914,24.749)
\qbezier(59.914,24.749)(59.914,24.775)(59.914,24.801)
\qbezier(59.914,24.801)(59.914,24.817)(59.914,24.833)
\qbezier(59.914,24.833)(59.914,24.837)(59.914,24.842)
\qbezier(59.914,24.842)(59.915,24.834)(59.915,24.827)
\qbezier(59.915,24.827)(59.915,24.808)(59.915,24.789)
\qbezier(59.915,24.789)(59.915,24.762)(59.915,24.735)
\qbezier(59.915,24.735)(59.915,24.704)(59.915,24.672)
\qbezier(59.915,24.672)(59.915,24.64)(59.916,24.608)
\qbezier(59.916,24.608)(59.916,24.58)(59.916,24.553)
\qbezier(59.916,24.553)(59.916,24.533)(59.916,24.513)
\qbezier(59.916,24.513)(59.916,24.505)(59.916,24.496)
\qbezier(59.916,24.496)(59.916,24.499)(59.916,24.502)
\qbezier(59.916,24.502)(59.916,24.516)(59.917,24.53)
\qbezier(59.917,24.53)(59.917,24.554)(59.917,24.578)
\qbezier(59.917,24.578)(59.917,24.608)(59.917,24.638)
\qbezier(59.917,24.638)(59.917,24.669)(59.917,24.701)
\qbezier(59.917,24.701)(59.917,24.73)(59.917,24.759)
\qbezier(59.917,24.759)(59.918,24.782)(59.918,24.804)
\qbezier(59.918,24.804)(59.918,24.817)(59.918,24.83)
\qbezier(59.918,24.83)(59.918,24.832)(59.918,24.833)
\qbezier(59.918,24.833)(59.918,24.823)(59.918,24.813)
\qbezier(59.918,24.813)(59.918,24.793)(59.918,24.773)
\qbezier(59.918,24.773)(59.919,24.745)(59.919,24.718)
\qbezier(59.919,24.718)(59.919,24.687)(59.919,24.656)
\qbezier(59.919,24.656)(59.919,24.626)(59.919,24.596)
\qbezier(59.919,24.596)(59.919,24.571)(59.919,24.546)
\qbezier(59.919,24.546)(59.919,24.53)(59.919,24.513)
\qbezier(59.919,24.513)(59.92,24.508)(59.92,24.502)
\qbezier(59.92,24.502)(59.92,24.507)(59.92,24.513)
\qbezier(59.92,24.513)(59.92,24.529)(59.92,24.546)
\qbezier(59.92,24.546)(59.92,24.57)(59.92,24.595)
\qbezier(59.92,24.595)(59.92,24.624)(59.92,24.654)
\qbezier(59.92,24.654)(59.92,24.684)(59.921,24.714)
\qbezier(59.921,24.714)(59.921,24.741)(59.921,24.768)
\qbezier(59.921,24.768)(59.921,24.787)(59.921,24.807)
\qbezier(59.921,24.807)(59.921,24.817)(59.921,24.827)
\qbezier(59.921,24.827)(59.921,24.825)(59.921,24.824)
\qbezier(59.921,24.824)(59.921,24.812)(59.922,24.8)
\qbezier(59.922,24.8)(59.922,24.778)(59.922,24.757)
\qbezier(59.922,24.757)(59.922,24.729)(59.922,24.702)
\qbezier(60.078,24.702)(60.078,24.729)(60.078,24.757)
\qbezier(60.078,24.757)(60.078,24.778)(60.078,24.8)
\qbezier(60.078,24.8)(60.079,24.812)(60.079,24.824)
\qbezier(60.079,24.824)(60.079,24.825)(60.079,24.827)
\qbezier(60.079,24.827)(60.079,24.817)(60.079,24.807)
\qbezier(60.079,24.807)(60.079,24.787)(60.079,24.768)
\qbezier(60.079,24.768)(60.079,24.741)(60.079,24.714)
\qbezier(60.079,24.714)(60.08,24.684)(60.08,24.654)
\qbezier(60.08,24.654)(60.08,24.624)(60.08,24.595)
\qbezier(60.08,24.595)(60.08,24.57)(60.08,24.546)
\qbezier(60.08,24.546)(60.08,24.529)(60.08,24.513)
\qbezier(60.08,24.513)(60.08,24.507)(60.08,24.502)
\qbezier(60.08,24.502)(60.08,24.508)(60.081,24.513)
\qbezier(60.081,24.513)(60.081,24.53)(60.081,24.546)
\qbezier(60.081,24.546)(60.081,24.571)(60.081,24.596)
\qbezier(60.081,24.596)(60.081,24.626)(60.081,24.656)
\qbezier(60.081,24.656)(60.081,24.687)(60.081,24.718)
\qbezier(60.081,24.718)(60.081,24.745)(60.082,24.773)
\qbezier(60.082,24.773)(60.082,24.793)(60.082,24.813)
\qbezier(60.082,24.813)(60.082,24.823)(60.082,24.833)
\qbezier(60.082,24.833)(60.082,24.832)(60.082,24.83)
\qbezier(60.082,24.83)(60.082,24.817)(60.082,24.804)
\qbezier(60.082,24.804)(60.082,24.782)(60.083,24.759)
\qbezier(60.083,24.759)(60.083,24.73)(60.083,24.701)
\qbezier(60.083,24.701)(60.083,24.669)(60.083,24.638)
\qbezier(60.083,24.638)(60.083,24.608)(60.083,24.578)
\qbezier(60.083,24.578)(60.083,24.554)(60.083,24.53)
\qbezier(60.083,24.53)(60.084,24.516)(60.084,24.502)
\qbezier(60.084,24.502)(60.084,24.499)(60.084,24.496)
\qbezier(60.084,24.496)(60.084,24.505)(60.084,24.513)
\qbezier(60.084,24.513)(60.084,24.533)(60.084,24.553)
\qbezier(60.084,24.553)(60.084,24.58)(60.084,24.608)
\qbezier(60.084,24.608)(60.085,24.64)(60.085,24.672)
\qbezier(60.085,24.672)(60.085,24.704)(60.085,24.735)
\qbezier(60.085,24.735)(60.085,24.762)(60.085,24.789)
\qbezier(60.085,24.789)(60.085,24.808)(60.085,24.827)
\qbezier(60.085,24.827)(60.085,24.834)(60.086,24.842)
\qbezier(60.086,24.842)(60.086,24.837)(60.086,24.833)
\qbezier(60.086,24.833)(60.086,24.817)(60.086,24.801)
\qbezier(60.086,24.801)(60.086,24.775)(60.086,24.749)
\qbezier(60.086,24.749)(60.086,24.718)(60.086,24.686)
\qbezier(60.086,24.686)(60.087,24.653)(60.087,24.62)
\qbezier(60.087,24.62)(60.087,24.59)(60.087,24.56)
\qbezier(60.087,24.56)(60.087,24.538)(60.087,24.515)
\qbezier(60.087,24.515)(60.087,24.503)(60.087,24.49)
\qbezier(60.087,24.49)(60.087,24.49)(60.088,24.49)
\qbezier(60.088,24.49)(60.088,24.503)(60.088,24.515)
\qbezier(60.088,24.515)(60.088,24.538)(60.088,24.561)
\qbezier(60.088,24.561)(60.088,24.591)(60.088,24.621)
\qbezier(60.088,24.621)(60.088,24.655)(60.088,24.689)
\qbezier(60.088,24.689)(60.089,24.721)(60.089,24.753)
\qbezier(60.089,24.753)(60.089,24.78)(60.089,24.806)
\qbezier(60.089,24.806)(60.089,24.823)(60.089,24.84)
\qbezier(60.089,24.84)(60.089,24.845)(60.089,24.85)
\qbezier(60.089,24.85)(60.09,24.842)(60.09,24.834)
\qbezier(60.09,24.834)(60.09,24.815)(60.09,24.795)
\qbezier(60.09,24.795)(60.09,24.766)(60.09,24.738)
\qbezier(60.09,24.738)(60.09,24.704)(60.09,24.67)
\qbezier(60.09,24.67)(60.091,24.636)(60.091,24.602)
\qbezier(60.091,24.602)(60.091,24.572)(60.091,24.542)
\qbezier(60.091,24.542)(60.091,24.521)(60.091,24.5)
\qbezier(60.091,24.5)(60.091,24.49)(60.091,24.48)
\qbezier(60.091,24.48)(60.092,24.483)(60.092,24.486)
\qbezier(60.092,24.486)(60.092,24.502)(60.092,24.518)
\qbezier(60.092,24.518)(60.092,24.544)(60.092,24.57)
\qbezier(60.092,24.57)(60.092,24.603)(60.092,24.636)
\qbezier(60.092,24.636)(60.093,24.672)(60.093,24.707)
\qbezier(60.093,24.707)(60.093,24.74)(60.093,24.772)
\qbezier(60.093,24.772)(60.093,24.798)(60.093,24.824)
\qbezier(60.093,24.824)(60.093,24.838)(60.093,24.853)
\qbezier(60.093,24.853)(60.094,24.855)(60.094,24.857)
\qbezier(60.094,24.857)(60.094,24.846)(60.094,24.834)
\qbezier(60.094,24.834)(60.094,24.811)(60.094,24.788)
\qbezier(60.094,24.788)(60.094,24.756)(60.095,24.725)
\qbezier(60.095,24.725)(60.095,24.689)(60.095,24.653)
\qbezier(60.095,24.653)(60.095,24.618)(60.095,24.582)
\qbezier(60.095,24.582)(60.095,24.553)(60.095,24.523)
\qbezier(60.095,24.523)(60.095,24.504)(60.096,24.484)
\qbezier(60.096,24.484)(60.096,24.477)(60.096,24.47)
\qbezier(60.096,24.47)(60.096,24.477)(60.096,24.483)
\qbezier(60.096,24.483)(60.096,24.503)(60.096,24.522)
\qbezier(60.096,24.522)(60.097,24.552)(60.097,24.581)
\qbezier(60.097,24.581)(60.097,24.617)(60.097,24.653)
\qbezier(60.097,24.653)(60.097,24.69)(60.097,24.727)
\qbezier(60.097,24.727)(60.097,24.76)(60.098,24.793)
\qbezier(60.098,24.793)(60.098,24.817)(60.098,24.841)
\qbezier(60.098,24.841)(60.098,24.854)(60.098,24.866)
\qbezier(60.098,24.866)(60.098,24.865)(60.098,24.863)
\qbezier(60.098,24.863)(60.099,24.848)(60.099,24.833)
\qbezier(60.099,24.833)(60.099,24.806)(60.099,24.779)
\qbezier(60.099,24.779)(60.099,24.744)(60.099,24.709)
\qbezier(60.099,24.709)(60.099,24.671)(60.1,24.633)
\qbezier(60.1,24.633)(60.1,24.597)(60.1,24.562)
\qbezier(60.1,24.562)(60.1,24.533)(60.1,24.504)
\qbezier(60.1,24.504)(60.1,24.486)(60.101,24.469)
\qbezier(60.101,24.469)(60.101,24.465)(60.101,24.461)
\qbezier(60.101,24.461)(60.101,24.471)(60.101,24.482)
\qbezier(60.101,24.482)(60.101,24.505)(60.101,24.528)
\qbezier(60.101,24.528)(60.102,24.561)(60.102,24.595)
\qbezier(60.102,24.595)(60.102,24.633)(60.102,24.671)
\qbezier(60.102,24.671)(60.102,24.71)(60.102,24.748)
\qbezier(60.102,24.748)(60.103,24.781)(60.103,24.814)
\qbezier(60.103,24.814)(60.103,24.836)(60.103,24.859)
\qbezier(60.103,24.859)(60.103,24.869)(60.103,24.878)
\qbezier(60.103,24.878)(60.103,24.873)(60.104,24.868)
\qbezier(60.104,24.868)(60.104,24.849)(60.104,24.83)
\qbezier(60.104,24.83)(60.104,24.799)(60.104,24.768)
\qbezier(60.104,24.768)(60.104,24.73)(60.105,24.692)
\qbezier(60.105,24.692)(60.105,24.652)(60.105,24.612)
\qbezier(60.105,24.612)(60.105,24.576)(60.105,24.539)
\qbezier(60.105,24.539)(60.105,24.512)(60.106,24.484)
\qbezier(60.106,24.484)(60.106,24.469)(60.106,24.453)
\qbezier(60.106,24.453)(60.106,24.453)(60.106,24.452)
\qbezier(60.106,24.452)(60.106,24.467)(60.107,24.481)
\qbezier(60.107,24.481)(60.107,24.509)(60.107,24.536)
\qbezier(60.107,24.536)(60.107,24.573)(60.107,24.61)
\qbezier(60.107,24.61)(60.107,24.651)(60.108,24.692)
\qbezier(60.108,24.692)(60.108,24.731)(60.108,24.771)
\qbezier(60.108,24.771)(60.108,24.803)(60.108,24.836)
\qbezier(60.108,24.836)(60.109,24.857)(60.109,24.877)
\qbezier(60.109,24.877)(60.109,24.884)(60.109,24.89)
\qbezier(60.109,24.89)(60.109,24.881)(60.109,24.872)
\qbezier(60.109,24.872)(60.11,24.848)(60.11,24.825)
\qbezier(60.11,24.825)(60.11,24.79)(60.11,24.755)
\qbezier(60.11,24.755)(60.11,24.714)(60.11,24.673)
\qbezier(60.11,24.673)(60.111,24.631)(60.111,24.589)
\qbezier(60.111,24.589)(60.111,24.553)(60.111,24.516)
\qbezier(60.111,24.516)(60.111,24.489)(60.112,24.463)
\qbezier(60.112,24.463)(60.112,24.45)(60.112,24.438)
\qbezier(60.112,24.438)(60.112,24.442)(60.112,24.445)
\qbezier(60.112,24.445)(60.113,24.464)(60.113,24.483)
\qbezier(60.113,24.483)(60.113,24.515)(60.113,24.547)
\qbezier(60.113,24.547)(60.113,24.587)(60.114,24.627)
\qbezier(60.114,24.627)(60.114,24.671)(60.114,24.714)
\qbezier(60.114,24.714)(60.114,24.755)(60.114,24.795)
\qbezier(60.114,24.795)(60.114,24.827)(60.115,24.859)
\qbezier(60.115,24.859)(60.115,24.877)(60.115,24.896)
\qbezier(60.115,24.896)(60.115,24.899)(60.115,24.901)
\qbezier(60.115,24.901)(60.116,24.888)(60.116,24.874)
\qbezier(60.116,24.874)(60.116,24.846)(60.116,24.818)
\qbezier(60.116,24.818)(60.116,24.779)(60.117,24.74)
\qbezier(60.117,24.74)(60.117,24.696)(60.117,24.651)
\qbezier(60.117,24.651)(60.117,24.608)(60.118,24.564)
\qbezier(60.118,24.564)(60.118,24.527)(60.118,24.491)
\qbezier(60.118,24.491)(60.118,24.466)(60.118,24.441)
\qbezier(60.118,24.441)(60.119,24.432)(60.119,24.423)
\qbezier(60.119,24.423)(60.119,24.431)(60.119,24.439)
\qbezier(60.119,24.439)(60.119,24.462)(60.12,24.486)
\qbezier(60.12,24.486)(60.12,24.523)(60.12,24.559)
\qbezier(60.12,24.559)(60.12,24.603)(60.12,24.648)
\qbezier(60.12,24.648)(60.121,24.693)(60.121,24.739)
\qbezier(60.121,24.739)(60.121,24.781)(60.121,24.822)
\qbezier(60.121,24.822)(60.122,24.852)(60.122,24.883)
\qbezier(60.122,24.883)(60.122,24.899)(60.122,24.915)
\qbezier(60.122,24.915)(60.122,24.913)(60.123,24.912)
\qbezier(60.123,24.912)(60.123,24.893)(60.123,24.875)
\qbezier(60.123,24.875)(60.123,24.842)(60.124,24.809)
\qbezier(60.124,24.809)(60.124,24.765)(60.124,24.722)
\qbezier(60.124,24.722)(60.124,24.674)(60.125,24.627)
\qbezier(60.125,24.627)(60.125,24.582)(60.125,24.537)
\qbezier(60.125,24.537)(60.125,24.5)(60.125,24.464)
\qbezier(60.125,24.464)(60.126,24.442)(60.126,24.419)
\qbezier(60.126,24.419)(60.126,24.414)(60.126,24.408)
\qbezier(60.126,24.408)(60.127,24.421)(60.127,24.433)
\qbezier(60.127,24.433)(60.127,24.462)(60.127,24.491)
\qbezier(60.127,24.491)(60.128,24.533)(60.128,24.574)
\qbezier(60.128,24.574)(60.128,24.622)(60.128,24.67)
\qbezier(60.128,24.67)(60.129,24.719)(60.129,24.767)
\qbezier(60.129,24.767)(60.129,24.809)(60.129,24.85)
\qbezier(60.129,24.85)(60.13,24.879)(60.13,24.909)
\qbezier(60.13,24.909)(60.13,24.921)(60.13,24.934)
\qbezier(60.13,24.934)(60.131,24.928)(60.131,24.922)
\qbezier(60.131,24.922)(60.131,24.898)(60.131,24.874)
\qbezier(60.131,24.874)(60.132,24.835)(60.132,24.797)
\qbezier(60.132,24.797)(60.132,24.749)(60.132,24.701)
\qbezier(60.132,24.701)(60.133,24.65)(60.133,24.6)
\qbezier(60.133,24.6)(60.133,24.553)(60.134,24.507)
\qbezier(60.134,24.507)(60.134,24.471)(60.134,24.436)
\qbezier(60.134,24.436)(60.134,24.416)(60.135,24.396)
\qbezier(60.135,24.396)(60.135,24.395)(60.135,24.394)
\qbezier(60.135,24.394)(60.135,24.412)(60.136,24.43)
\qbezier(60.136,24.43)(60.136,24.464)(60.136,24.499)
\qbezier(60.136,24.499)(60.137,24.546)(60.137,24.592)
\qbezier(60.137,24.592)(60.137,24.644)(60.137,24.697)
\qbezier(60.137,24.697)(60.138,24.747)(60.138,24.798)
\qbezier(60.138,24.798)(60.138,24.839)(60.139,24.881)
\qbezier(60.139,24.881)(60.139,24.908)(60.139,24.936)
\qbezier(60.139,24.936)(60.139,24.944)(60.14,24.953)
\qbezier(60.14,24.953)(60.14,24.942)(60.14,24.93)
\qbezier(60.14,24.93)(60.141,24.901)(60.141,24.871)
\qbezier(60.141,24.871)(60.141,24.827)(60.141,24.782)
\qbezier(60.141,24.782)(60.142,24.73)(60.142,24.677)
\qbezier(60.142,24.677)(60.142,24.623)(60.143,24.569)
\qbezier(60.143,24.569)(60.143,24.522)(60.143,24.474)
\qbezier(60.143,24.474)(60.144,24.44)(60.144,24.405)
\qbezier(60.144,24.405)(60.144,24.388)(60.145,24.372)
\qbezier(60.145,24.372)(60.145,24.376)(60.145,24.38)
\qbezier(60.145,24.38)(60.145,24.404)(60.146,24.428)
\qbezier(60.146,24.428)(60.146,24.468)(60.146,24.509)
\qbezier(60.146,24.509)(60.147,24.561)(60.147,24.614)
\qbezier(60.147,24.614)(60.147,24.67)(60.148,24.726)
\qbezier(60.148,24.726)(60.148,24.779)(60.148,24.832)
\qbezier(60.148,24.832)(60.149,24.873)(60.149,24.915)
\qbezier(60.149,24.915)(60.149,24.94)(60.15,24.964)
\qbezier(60.15,24.964)(60.15,24.969)(60.15,24.973)
\qbezier(60.15,24.973)(60.151,24.955)(60.151,24.938)
\qbezier(60.151,24.938)(60.151,24.902)(60.152,24.866)
\qbezier(60.152,24.866)(60.152,24.815)(60.152,24.765)
\qbezier(60.152,24.765)(60.153,24.707)(60.153,24.649)
\qbezier(60.153,24.649)(60.153,24.592)(60.154,24.535)
\qbezier(60.154,24.535)(60.154,24.486)(60.155,24.438)
\qbezier(60.155,24.438)(60.155,24.405)(60.155,24.372)
\qbezier(60.155,24.372)(60.156,24.359)(60.156,24.347)
\qbezier(60.156,24.347)(60.156,24.356)(60.157,24.366)
\qbezier(60.157,24.366)(60.157,24.396)(60.157,24.427)
\qbezier(60.157,24.427)(60.158,24.475)(60.158,24.523)
\qbezier(60.158,24.523)(60.159,24.581)(60.159,24.639)
\qbezier(60.159,24.639)(60.159,24.7)(60.16,24.76)
\qbezier(60.16,24.76)(60.16,24.815)(60.161,24.87)
\qbezier(60.161,24.87)(60.161,24.911)(60.161,24.952)
\qbezier(60.161,24.952)(60.162,24.974)(60.162,24.995)
\qbezier(60.162,24.995)(60.162,24.994)(60.163,24.993)
\qbezier(60.163,24.993)(60.163,24.969)(60.164,24.945)
\qbezier(60.164,24.945)(60.164,24.901)(60.164,24.858)
\qbezier(60.164,24.858)(60.165,24.8)(60.165,24.743)
\qbezier(60.165,24.743)(60.166,24.68)(60.166,24.616)
\qbezier(60.166,24.616)(60.167,24.556)(60.167,24.496)
\qbezier(60.167,24.496)(60.167,24.446)(60.168,24.397)
\qbezier(60.168,24.397)(60.168,24.367)(60.169,24.336)
\qbezier(60.169,24.336)(60.169,24.328)(60.169,24.32)
\qbezier(60.169,24.32)(60.17,24.336)(60.17,24.352)
\qbezier(60.17,24.352)(60.171,24.391)(60.171,24.429)
\qbezier(60.171,24.429)(60.172,24.484)(60.172,24.54)
\qbezier(60.172,24.54)(60.173,24.604)(60.173,24.669)
\qbezier(60.173,24.669)(60.173,24.734)(60.174,24.8)
\qbezier(60.174,24.8)(60.174,24.856)(60.175,24.913)
\qbezier(60.175,24.913)(60.175,24.953)(60.176,24.993)
\qbezier(60.176,24.993)(60.176,25.011)(60.177,25.028)
\qbezier(60.177,25.028)(60.177,25.021)(60.178,25.014)
\qbezier(60.178,25.014)(60.178,24.982)(60.179,24.95)
\qbezier(60.179,24.95)(60.179,24.898)(60.18,24.846)
\qbezier(60.18,24.846)(60.18,24.781)(60.181,24.716)
\qbezier(60.181,24.716)(60.181,24.647)(60.181,24.578)
\qbezier(60.181,24.578)(60.182,24.514)(60.182,24.451)
\qbezier(60.182,24.451)(60.183,24.401)(60.183,24.352)
\qbezier(60.183,24.352)(60.184,24.324)(60.185,24.296)
\qbezier(60.185,24.296)(60.185,24.294)(60.186,24.292)
\qbezier(60.186,24.292)(60.186,24.315)(60.187,24.339)
\qbezier(60.187,24.339)(60.187,24.386)(60.188,24.433)
\qbezier(60.188,24.433)(60.188,24.497)(60.189,24.561)
\qbezier(60.189,24.561)(60.189,24.633)(60.19,24.705)
\qbezier(60.19,24.705)(60.19,24.775)(60.191,24.845)
\qbezier(60.191,24.845)(60.191,24.904)(60.192,24.962)
\qbezier(60.192,24.962)(60.192,25.001)(60.193,25.039)
\qbezier(60.193,25.039)(60.194,25.052)(60.194,25.065)
\qbezier(60.194,25.065)(60.195,25.05)(60.195,25.035)
\qbezier(60.195,25.035)(60.196,24.995)(60.196,24.954)
\qbezier(60.196,24.954)(60.197,24.893)(60.198,24.831)
\qbezier(60.198,24.831)(60.198,24.758)(60.199,24.684)
\qbezier(60.199,24.684)(60.199,24.609)(60.2,24.533)
\qbezier(60.2,24.533)(60.201,24.466)(60.201,24.399)
\qbezier(60.201,24.399)(60.202,24.349)(60.202,24.3)
\qbezier(60.202,24.3)(60.203,24.276)(60.204,24.252)
\qbezier(60.204,24.252)(60.204,24.256)(60.205,24.261)
\qbezier(60.205,24.261)(60.206,24.293)(60.206,24.326)
\qbezier(60.206,24.326)(60.207,24.384)(60.207,24.441)
\qbezier(60.207,24.441)(60.208,24.514)(60.209,24.588)
\qbezier(60.209,24.588)(60.209,24.668)(60.21,24.748)
\qbezier(60.21,24.748)(60.211,24.824)(60.211,24.899)
\qbezier(60.211,24.899)(60.212,24.959)(60.213,25.019)
\qbezier(60.213,25.019)(60.213,25.055)(60.214,25.092)
\qbezier(60.214,25.092)(60.215,25.099)(60.216,25.106)
\qbezier(60.216,25.106)(60.216,25.082)(60.217,25.058)
\qbezier(60.217,25.058)(60.218,25.007)(60.218,24.956)
\qbezier(60.218,24.956)(60.219,24.883)(60.22,24.811)
\qbezier(60.22,24.811)(60.221,24.728)(60.221,24.645)
\qbezier(60.221,24.645)(60.222,24.562)(60.223,24.479)
\qbezier(60.223,24.479)(60.223,24.408)(60.224,24.337)
\qbezier(60.224,24.337)(60.225,24.288)(60.226,24.24)
\qbezier(60.226,24.24)(60.227,24.22)(60.227,24.201)
\qbezier(60.227,24.201)(60.228,24.213)(60.229,24.226)
\qbezier(60.229,24.226)(60.23,24.27)(60.23,24.314)
\qbezier(60.23,24.314)(60.231,24.383)(60.232,24.452)
\qbezier(60.232,24.452)(60.233,24.537)(60.234,24.622)
\qbezier(60.234,24.622)(60.234,24.712)(60.235,24.801)
\qbezier(60.235,24.801)(60.236,24.882)(60.237,24.964)
\qbezier(60.237,24.964)(60.238,25.025)(60.239,25.087)
\qbezier(60.239,25.087)(60.24,25.12)(60.24,25.153)
\qbezier(60.24,25.153)(60.241,25.153)(60.242,25.152)
\qbezier(60.242,25.152)(60.243,25.118)(60.244,25.083)
\qbezier(60.244,25.083)(60.245,25.019)(60.246,24.955)
\qbezier(60.246,24.955)(60.247,24.87)(60.248,24.785)
\qbezier(60.248,24.785)(60.248,24.69)(60.249,24.595)
\qbezier(60.249,24.595)(60.25,24.504)(60.251,24.413)
\qbezier(60.251,24.413)(60.252,24.338)(60.253,24.263)
\qbezier(60.253,24.263)(60.254,24.215)(60.255,24.168)
\qbezier(60.255,24.168)(60.256,24.154)(60.257,24.141)
\qbezier(60.257,24.141)(60.258,24.164)(60.259,24.187)
\qbezier(60.259,24.187)(60.26,24.244)(60.261,24.301)
\qbezier(60.261,24.301)(60.262,24.385)(60.263,24.468)
\qbezier(60.263,24.468)(60.264,24.567)(60.265,24.666)
\qbezier(60.265,24.666)(60.266,24.767)(60.267,24.867)
\qbezier(60.267,24.867)(60.268,24.955)(60.27,25.043)
\qbezier(60.27,25.043)(60.271,25.106)(60.272,25.169)
\qbezier(60.272,25.169)(60.273,25.198)(60.274,25.227)
\qbezier(60.274,25.227)(60.275,25.217)(60.276,25.207)
\qbezier(60.276,25.207)(60.277,25.16)(60.279,25.112)
\qbezier(60.279,25.112)(60.28,25.032)(60.281,24.952)
\qbezier(60.281,24.952)(60.282,24.851)(60.283,24.749)
\qbezier(60.283,24.749)(60.285,24.641)(60.286,24.532)
\qbezier(60.286,24.532)(60.287,24.431)(60.288,24.329)
\qbezier(60.288,24.329)(60.289,24.25)(60.291,24.171)
\qbezier(60.291,24.171)(60.292,24.125)(60.293,24.079)
\qbezier(60.293,24.079)(60.295,24.073)(60.296,24.068)
\qbezier(60.296,24.068)(60.297,24.104)(60.299,24.14)
\qbezier(60.299,24.14)(60.3,24.214)(60.301,24.288)
\qbezier(60.301,24.288)(60.303,24.39)(60.304,24.492)
\qbezier(60.304,24.492)(60.305,24.608)(60.307,24.724)
\qbezier(60.307,24.724)(60.308,24.838)(60.31,24.952)
\qbezier(60.31,24.952)(60.311,25.048)(60.313,25.144)
\qbezier(60.313,25.144)(60.314,25.209)(60.315,25.273)
\qbezier(60.315,25.273)(60.317,25.297)(60.318,25.32)
\qbezier(60.318,25.32)(60.32,25.298)(60.322,25.275)
\qbezier(60.322,25.275)(60.323,25.21)(60.325,25.145)
\qbezier(60.325,25.145)(60.326,25.045)(60.328,24.945)
\qbezier(60.328,24.945)(60.329,24.823)(60.331,24.701)
\qbezier(60.331,24.701)(60.333,24.575)(60.334,24.448)
\qbezier(60.334,24.448)(60.336,24.334)(60.338,24.221)
\qbezier(60.338,24.221)(60.34,24.136)(60.341,24.051)
\qbezier(60.341,24.051)(60.343,24.008)(60.345,23.965)
\qbezier(60.345,23.965)(60.347,23.97)(60.348,23.975)
\qbezier(60.348,23.975)(60.35,24.029)(60.352,24.082)
\qbezier(60.352,24.082)(60.354,24.178)(60.356,24.274)
\qbezier(60.356,24.274)(60.358,24.4)(60.36,24.526)
\qbezier(60.36,24.526)(60.362,24.664)(60.364,24.802)
\qbezier(60.364,24.802)(60.366,24.934)(60.368,25.066)
\qbezier(60.368,25.066)(60.37,25.172)(60.372,25.279)
\qbezier(60.372,25.279)(60.374,25.345)(60.376,25.411)
\qbezier(60.376,25.411)(60.378,25.427)(60.38,25.442)
\qbezier(60.38,25.442)(60.382,25.403)(60.385,25.364)
\qbezier(60.385,25.364)(60.387,25.275)(60.389,25.186)
\qbezier(60.389,25.186)(60.391,25.059)(60.394,24.931)
\qbezier(60.394,24.931)(60.396,24.782)(60.398,24.633)
\qbezier(60.398,24.633)(60.401,24.483)(60.403,24.332)
\qbezier(60.403,24.332)(60.406,24.202)(60.408,24.071)
\qbezier(60.408,24.071)(60.411,23.979)(60.413,23.888)
\qbezier(60.413,23.888)(60.416,23.848)(60.418,23.809)
\qbezier(60.418,23.809)(60.421,23.829)(60.424,23.849)
\qbezier(60.424,23.849)(60.426,23.927)(60.429,24.005)
\qbezier(60.429,24.005)(60.432,24.132)(60.435,24.258)
\qbezier(60.435,24.258)(60.438,24.417)(60.441,24.576)
\qbezier(60.441,24.576)(60.443,24.745)(60.446,24.915)
\qbezier(60.446,24.915)(60.449,25.071)(60.452,25.228)
\qbezier(60.452,25.228)(60.456,25.349)(60.459,25.471)
\qbezier(60.459,25.471)(60.462,25.539)(60.465,25.607)
\qbezier(60.465,25.607)(60.468,25.611)(60.472,25.614)
\qbezier(60.472,25.614)(60.475,25.551)(60.478,25.488)
\qbezier(60.478,25.488)(60.482,25.366)(60.485,25.243)
\qbezier(60.485,25.243)(60.489,25.076)(60.493,24.909)
\qbezier(60.493,24.909)(60.496,24.72)(60.5,24.531)
\qbezier(60.5,24.531)(60.504,24.346)(60.508,24.16)
\qbezier(60.508,24.16)(60.512,24.005)(60.515,23.85)
\qbezier(60.515,23.85)(60.52,23.747)(60.524,23.645)
\qbezier(60.524,23.645)(60.528,23.611)(60.532,23.578)
\qbezier(60.532,23.578)(60.536,23.62)(60.541,23.662)
\qbezier(60.541,23.662)(60.545,23.777)(60.549,23.891)
\qbezier(60.549,23.891)(60.554,24.064)(60.559,24.238)
\qbezier(60.559,24.238)(60.563,24.447)(60.568,24.656)
\qbezier(60.568,24.656)(60.573,24.874)(60.578,25.091)
\qbezier(60.578,25.091)(60.583,25.286)(60.588,25.481)
\qbezier(60.588,25.481)(60.594,25.626)(60.599,25.77)
\qbezier(60.599,25.77)(60.604,25.842)(60.61,25.913)
\qbezier(60.61,25.913)(60.615,25.9)(60.621,25.886)
\qbezier(60.621,25.886)(60.627,25.785)(60.633,25.684)
\qbezier(60.633,25.684)(60.639,25.507)(60.645,25.33)
\qbezier(60.645,25.33)(60.652,25.1)(60.658,24.869)
\qbezier(60.658,24.869)(60.665,24.615)(60.671,24.36)
\qbezier(60.671,24.36)(60.678,24.117)(60.685,23.874)
\qbezier(60.685,23.874)(60.692,23.677)(60.699,23.48)
\qbezier(60.699,23.48)(60.707,23.358)(60.714,23.237)
\qbezier(60.714,23.237)(60.722,23.212)(60.73,23.187)
\qbezier(60.73,23.187)(60.738,23.266)(60.746,23.345)
\qbezier(60.746,23.345)(60.755,23.521)(60.763,23.698)
\qbezier(60.763,23.698)(60.772,23.952)(60.781,24.205)
\qbezier(60.781,24.205)(60.791,24.504)(60.8,24.803)
\qbezier(60.8,24.803)(60.81,25.106)(60.82,25.41)
\qbezier(60.82,25.41)(60.83,25.676)(60.84,25.942)
\qbezier(60.84,25.942)(60.851,26.131)(60.862,26.32)
\qbezier(60.862,26.32)(60.874,26.401)(60.885,26.482)
\qbezier(60.885,26.482)(60.897,26.438)(60.909,26.393)
\qbezier(60.909,26.393)(60.922,26.223)(60.935,26.053)
\qbezier(60.935,26.053)(60.948,25.774)(60.962,25.495)
\qbezier(60.962,25.495)(60.976,25.141)(60.99,24.786)
\qbezier(60.99,24.786)(61.005,24.401)(61.02,24.015)
\qbezier(61.02,24.015)(61.037,23.653)(61.053,23.29)
\qbezier(61.053,23.29)(61.07,23.002)(61.087,22.715)
\qbezier(61.087,22.715)(61.105,22.548)(61.124,22.381)
\qbezier(61.124,22.381)(61.143,22.367)(61.163,22.353)
\qbezier(61.163,22.353)(61.184,22.505)(61.205,22.656)
\qbezier(61.205,22.656)(61.227,22.964)(61.25,23.272)
\qbezier(61.25,23.272)(61.25,23.272)(61.25,23.272)
\qbezier(61.25,23.272)(61.438,25.614)(61.625,27.957)
\qbezier(61.625,27.957)(61.812,26.244)(62,24.531)
\qbezier(62,24.531)(62.188,22.523)(62.375,20.515)
\qbezier(62.375,20.515)(62.562,19.81)(62.75,19.104)
\qbezier(62.75,19.104)(62.938,19.46)(63.125,19.815)
\qbezier(63.125,19.815)(63.312,20.745)(63.5,21.675)
\qbezier(63.5,21.675)(63.688,22.836)(63.875,23.998)
\qbezier(63.875,23.998)(64.062,25.2)(64.25,26.402)
\qbezier(64.25,26.402)(64.438,27.55)(64.625,28.698)
\qbezier(64.625,28.698)(64.812,29.751)(65,30.805)
\qbezier(65,30.805)(65.188,31.751)(65.375,32.698)
\qbezier(65.375,32.698)(65.562,33.539)(65.75,34.381)
\qbezier(65.75,34.381)(65.938,35.124)(66.125,35.868)
\qbezier(66.125,35.868)(66.312,36.525)(66.5,37.181)
\qbezier(66.5,37.181)(66.688,37.761)(66.875,38.34)
\qbezier(66.875,38.34)(67.062,38.852)(67.25,39.364)
\qbezier(67.25,39.364)(67.438,39.817)(67.625,40.271)
\qbezier(67.625,40.271)(67.812,40.674)(68,41.076)
\qbezier(68,41.076)(68.188,41.435)(68.375,41.794)
\qbezier(68.375,41.794)(68.562,42.114)(68.75,42.434)
\qbezier(68.75,42.434)(68.938,42.721)(69.125,43.008)
\qbezier(69.125,43.008)(69.312,43.266)(69.5,43.524)
\qbezier(69.5,43.524)(69.688,43.756)(69.875,43.988)
\qbezier(69.875,43.988)(70.062,44.198)(70.25,44.408)
\qbezier(70.25,44.408)(70.438,44.599)(70.625,44.789)
\qbezier(70.625,44.789)(70.812,44.962)(71,45.135)
\qbezier(71,45.135)(71.188,45.292)(71.375,45.45)
\qbezier(71.375,45.45)(71.562,45.594)(71.75,45.738)
\qbezier(71.75,45.738)(71.938,45.869)(72.125,46.001)
\qbezier(72.125,46.001)(72.312,46.122)(72.5,46.243)
\qbezier(72.5,46.243)(72.688,46.354)(72.875,46.465)
\qbezier(72.875,46.465)(73.062,46.568)(73.25,46.67)
\qbezier(73.25,46.67)(73.438,46.765)(73.625,46.859)
\qbezier(73.625,46.859)(73.812,46.947)(74,47.034)
\qbezier(74,47.034)(74.188,47.115)(74.375,47.196)
\qbezier(74.375,47.196)(74.562,47.272)(74.75,47.347)
\qbezier(74.75,47.347)(74.938,47.417)(75.125,47.487)
\qbezier(75.125,47.487)(75.312,47.552)(75.5,47.617)
\qbezier(75.5,47.617)(75.688,47.678)(75.875,47.739)
\qbezier(75.875,47.739)(76.062,47.796)(76.25,47.853)
\qbezier(76.25,47.853)(76.438,47.906)(76.625,47.959)
\qbezier(76.625,47.959)(76.812,48.009)(77,48.059)
\qbezier(77,48.059)(77.188,48.106)(77.375,48.152)
\qbezier(77.375,48.152)(77.562,48.196)(77.75,48.24)
\qbezier(77.75,48.24)(77.938,48.282)(78.125,48.323)
\qbezier(78.125,48.323)(78.312,48.362)(78.5,48.401)
\qbezier(78.5,48.401)(78.688,48.437)(78.875,48.474)
\qbezier(78.875,48.474)(79.062,48.509)(79.25,48.543)
\qbezier(79.25,48.543)(79.438,48.576)(79.625,48.609)
\qbezier(79.625,48.609)(79.812,48.64)(80,48.671)
\qbezier(80,48.671)(80.188,48.7)(80.375,48.729)
\qbezier(80.375,48.729)(80.562,48.757)(80.75,48.785)
\qbezier(80.75,48.785)(80.938,48.811)(81.125,48.837)
\qbezier(81.125,48.837)(81.312,48.862)(81.5,48.887)
\qbezier(81.5,48.887)(81.688,48.911)(81.875,48.935)
\qbezier(81.875,48.935)(82.062,48.957)(82.25,48.98)
\qbezier(82.25,48.98)(82.438,49.002)(82.625,49.023)
\qbezier(82.625,49.023)(82.812,49.043)(83,49.064)
\qbezier(83,49.064)(83.188,49.083)(83.375,49.103)
\qbezier(83.375,49.103)(83.562,49.122)(83.75,49.14)
\qbezier(83.75,49.14)(83.938,49.158)(84.125,49.176)
\qbezier(84.125,49.176)(84.312,49.193)(84.5,49.21)
\qbezier(84.5,49.21)(84.688,49.226)(84.875,49.242)
\qbezier(84.875,49.242)(85.062,49.258)(85.25,49.273)
\qbezier(85.25,49.273)(85.438,49.288)(85.625,49.303)
\qbezier(85.625,49.303)(85.812,49.317)(86,49.331)
\qbezier(86,49.331)(86.188,49.345)(86.375,49.359)
\qbezier(86.375,49.359)(86.562,49.372)(86.75,49.385)
\qbezier(86.75,49.385)(86.938,49.397)(87.125,49.41)
\qbezier(87.125,49.41)(87.312,49.422)(87.5,49.434)
\qbezier(87.5,49.434)(87.688,49.445)(87.875,49.457)
\qbezier(87.875,49.457)(88.062,49.468)(88.25,49.479)
\qbezier(88.25,49.479)(88.438,49.49)(88.625,49.501)
\qbezier(88.625,49.501)(88.812,49.511)(89,49.521)
\qbezier(89,49.521)(89.188,49.531)(89.375,49.541)
\qbezier(89.375,49.541)(89.562,49.55)(89.75,49.56)
\qbezier(89.75,49.56)(89.938,49.569)(90.125,49.578)
\qbezier(90.125,49.578)(90.312,49.587)(90.5,49.596)
\qbezier(90.5,49.596)(90.688,49.604)(90.875,49.613)
\qbezier(90.875,49.613)(91.062,49.621)(91.25,49.629)
\qbezier(91.25,49.629)(91.438,49.637)(91.625,49.645)
\qbezier(91.625,49.645)(91.812,49.653)(92,49.661)
\qbezier(92,49.661)(92.188,49.668)(92.375,49.675)
\qbezier(92.375,49.675)(92.562,49.683)(92.75,49.69)
\qbezier(92.75,49.69)(92.938,49.697)(93.125,49.703)
\qbezier(93.125,49.703)(93.312,49.71)(93.5,49.717)
\qbezier(93.5,49.717)(93.688,49.723)(93.875,49.73)
\qbezier(93.875,49.73)(94.062,49.736)(94.25,49.742)
\qbezier(94.25,49.742)(94.438,49.748)(94.625,49.754)
\qbezier(94.625,49.754)(94.812,49.76)(95,49.766)
\qbezier(95,49.766)(95.188,49.772)(95.375,49.777)
\qbezier(95.375,49.777)(95.562,49.783)(95.75,49.788)
\qbezier(95.75,49.788)(95.938,49.794)(96.125,49.799)
\qbezier(96.125,49.799)(96.312,49.804)(96.5,49.809)
\qbezier(96.5,49.809)(96.688,49.814)(96.875,49.819)
\qbezier(96.875,49.819)(97.062,49.824)(97.25,49.829)
\qbezier(97.25,49.829)(97.438,49.834)(97.625,49.839)
\qbezier(97.625,49.839)(97.812,49.843)(98,49.848)
\qbezier(98,49.848)(98.188,49.852)(98.375,49.857)
\qbezier(98.375,49.857)(98.562,49.861)(98.75,49.865)
\qbezier(98.75,49.865)(98.938,49.87)(99.125,49.874)
\qbezier(99.125,49.874)(99.312,49.878)(99.5,49.882)
\qbezier(99.5,49.882)(99.688,49.886)(99.875,49.89)
\qbezier(99.875,49.89)(100.062,49.894)(100.25,49.898)
\qbezier(100.25,49.898)(100.438,49.901)(100.625,49.905)
\qbezier(100.625,49.905)(100.812,49.909)(101,49.912)
\qbezier(101,49.912)(101.188,49.916)(101.375,49.919)
\qbezier(101.375,49.919)(101.562,49.923)(101.75,49.926)
\qbezier(101.75,49.926)(101.938,49.93)(102.125,49.933)
\qbezier(102.125,49.933)(102.312,49.936)(102.5,49.94)
\qbezier(102.5,49.94)(102.688,49.943)(102.875,49.946)
\qbezier(102.875,49.946)(103.062,49.949)(103.25,49.952)
\qbezier(103.25,49.952)(103.438,49.955)(103.625,49.958)
\qbezier(103.625,49.958)(103.812,49.961)(104,49.964)
\qbezier(104,49.964)(104.188,49.967)(104.375,49.97)
\qbezier(104.375,49.97)(104.562,49.973)(104.75,49.976)
\qbezier(104.75,49.976)(104.938,49.978)(105.125,49.981)
\qbezier(105.125,49.981)(105.312,49.984)(105.5,49.986)
\qbezier(105.5,49.986)(105.688,49.989)(105.875,49.992)
\qbezier(105.875,49.992)(106.062,49.994)(106.25,49.997)
\qbezier(106.25,49.997)(106.438,49.999)(106.625,50.002)
\qbezier(106.625,50.002)(106.812,50.004)(107,50.006)
\qbezier(107,50.006)(107.188,50.009)(107.375,50.011)
\qbezier(107.375,50.011)(107.562,50.014)(107.75,50.016)
\qbezier(107.75,50.016)(107.938,50.018)(108.125,50.02)
\qbezier(108.125,50.02)(108.312,50.023)(108.5,50.025)
\qbezier(108.5,50.025)(108.688,50.027)(108.875,50.029)
\qbezier(108.875,50.029)(109.062,50.031)(109.25,50.033)
\qbezier(109.25,50.033)(109.438,50.035)(109.625,50.037)
\qbezier(109.625,50.037)(109.812,50.039)(110,50.041)
\color{black}
\put(60,24.667){{\color{white}\circle*{1.8}}}
\put(60,24.667){\circle{1.8}}
\end{picture}
\endgroup}
\newcommand{\rossSketchReciprocalMinusX}{%
\begingroup
\setlength{\unitlength}{1mm}
\begin{picture}(120,65)
\linethickness{0.3pt}
\put(7,33){\vector(1,0){107}}
\put(60,5){\vector(0,1){57}}
\put(115,32){\makebox(0,0)[l]{$x$}}
\put(61.5,62){\makebox(0,0)[l]{$y$}}
\put(43.333,32.3){\line(0,1){1.4}}
\put(43.333,30.3){\makebox(0,0)[t]{\scriptsize $-1$}}
\put(76.667,32.3){\line(0,1){1.4}}
\put(76.667,30.3){\makebox(0,0)[t]{\scriptsize $1$}}
\put(59.3,14.25){\line(1,0){1.4}}
\put(58.3,14.25){\makebox(0,0)[r]{\scriptsize $-3$}}
\put(59.3,51.75){\line(1,0){1.4}}
\put(58.3,51.75){\makebox(0,0)[r]{\scriptsize $3$}}
\qbezier[35](10,51.75)(60,33)(110,14.25)
\color{mathjockgreen}
\linethickness{0.6pt}
\qbezier(10,49.667)(10.138,49.609)(10.276,49.552)
\qbezier(10.276,49.552)(10.414,49.494)(10.552,49.436)
\qbezier(10.552,49.436)(10.69,49.379)(10.828,49.321)
\qbezier(10.828,49.321)(10.966,49.263)(11.104,49.206)
\qbezier(11.104,49.206)(11.242,49.148)(11.38,49.09)
\qbezier(11.38,49.09)(11.518,49.032)(11.656,48.974)
\qbezier(11.656,48.974)(11.794,48.917)(11.931,48.859)
\qbezier(11.931,48.859)(12.069,48.801)(12.207,48.743)
\qbezier(12.207,48.743)(12.345,48.685)(12.483,48.627)
\qbezier(12.483,48.627)(12.621,48.568)(12.759,48.51)
\qbezier(12.759,48.51)(12.897,48.452)(13.035,48.394)
\qbezier(13.035,48.394)(13.173,48.336)(13.311,48.277)
\qbezier(13.311,48.277)(13.449,48.219)(13.587,48.161)
\qbezier(13.587,48.161)(13.725,48.102)(13.863,48.044)
\qbezier(13.863,48.044)(14.001,47.985)(14.139,47.927)
\qbezier(14.139,47.927)(14.277,47.868)(14.415,47.809)
\qbezier(14.415,47.809)(14.553,47.751)(14.691,47.692)
\qbezier(14.691,47.692)(14.829,47.633)(14.967,47.574)
\qbezier(14.967,47.574)(15.105,47.516)(15.243,47.457)
\qbezier(15.243,47.457)(15.381,47.398)(15.519,47.339)
\qbezier(15.519,47.339)(15.656,47.28)(15.794,47.221)
\qbezier(15.794,47.221)(15.932,47.162)(16.07,47.102)
\qbezier(16.07,47.102)(16.208,47.043)(16.346,46.984)
\qbezier(16.346,46.984)(16.484,46.925)(16.622,46.865)
\qbezier(16.622,46.865)(16.76,46.806)(16.898,46.746)
\qbezier(16.898,46.746)(17.036,46.687)(17.174,46.627)
\qbezier(17.174,46.627)(17.312,46.568)(17.45,46.508)
\qbezier(17.45,46.508)(17.588,46.448)(17.726,46.389)
\qbezier(17.726,46.389)(17.864,46.329)(18.002,46.269)
\qbezier(18.002,46.269)(18.14,46.209)(18.278,46.149)
\qbezier(18.278,46.149)(18.416,46.089)(18.554,46.029)
\qbezier(18.554,46.029)(18.692,45.969)(18.83,45.909)
\qbezier(18.83,45.909)(18.968,45.848)(19.106,45.788)
\qbezier(19.106,45.788)(19.244,45.728)(19.381,45.667)
\qbezier(19.381,45.667)(19.519,45.607)(19.657,45.546)
\qbezier(19.657,45.546)(19.795,45.486)(19.933,45.425)
\qbezier(19.933,45.425)(20.071,45.364)(20.209,45.304)
\qbezier(20.209,45.304)(20.347,45.243)(20.485,45.182)
\qbezier(20.485,45.182)(20.623,45.121)(20.761,45.06)
\qbezier(20.761,45.06)(20.899,44.999)(21.037,44.938)
\qbezier(21.037,44.938)(21.175,44.876)(21.313,44.815)
\qbezier(21.313,44.815)(21.451,44.754)(21.589,44.692)
\qbezier(21.589,44.692)(21.727,44.631)(21.865,44.569)
\qbezier(21.865,44.569)(22.003,44.507)(22.141,44.446)
\qbezier(22.141,44.446)(22.279,44.384)(22.417,44.322)
\qbezier(22.417,44.322)(22.555,44.26)(22.693,44.198)
\qbezier(22.693,44.198)(22.831,44.136)(22.969,44.074)
\qbezier(22.969,44.074)(23.106,44.012)(23.244,43.949)
\qbezier(23.244,43.949)(23.382,43.887)(23.52,43.824)
\qbezier(23.52,43.824)(23.658,43.762)(23.796,43.699)
\qbezier(23.796,43.699)(23.934,43.636)(24.072,43.574)
\qbezier(24.072,43.574)(24.21,43.511)(24.348,43.448)
\qbezier(24.348,43.448)(24.486,43.385)(24.624,43.321)
\qbezier(24.624,43.321)(24.762,43.258)(24.9,43.195)
\qbezier(24.9,43.195)(25.038,43.131)(25.176,43.068)
\qbezier(25.176,43.068)(25.314,43.004)(25.452,42.94)
\qbezier(25.452,42.94)(25.59,42.877)(25.728,42.813)
\qbezier(25.728,42.813)(25.866,42.749)(26.004,42.685)
\qbezier(26.004,42.685)(26.142,42.62)(26.28,42.556)
\qbezier(26.28,42.556)(26.418,42.492)(26.556,42.427)
\qbezier(26.556,42.427)(26.694,42.362)(26.831,42.298)
\qbezier(26.831,42.298)(26.969,42.233)(27.107,42.168)
\qbezier(27.107,42.168)(27.245,42.103)(27.383,42.038)
\qbezier(27.383,42.038)(27.521,41.972)(27.659,41.907)
\qbezier(27.659,41.907)(27.797,41.841)(27.935,41.776)
\qbezier(27.935,41.776)(28.073,41.71)(28.211,41.644)
\qbezier(28.211,41.644)(28.349,41.578)(28.487,41.512)
\qbezier(28.487,41.512)(28.625,41.446)(28.763,41.379)
\qbezier(28.763,41.379)(28.901,41.313)(29.039,41.246)
\qbezier(29.039,41.246)(29.177,41.179)(29.315,41.112)
\qbezier(29.315,41.112)(29.453,41.045)(29.591,40.978)
\qbezier(29.591,40.978)(29.729,40.911)(29.867,40.843)
\qbezier(29.867,40.843)(30.005,40.775)(30.143,40.708)
\qbezier(30.143,40.708)(30.281,40.64)(30.419,40.572)
\qbezier(30.419,40.572)(30.556,40.503)(30.694,40.435)
\qbezier(30.694,40.435)(30.832,40.366)(30.97,40.298)
\qbezier(30.97,40.298)(31.108,40.229)(31.246,40.16)
\qbezier(31.246,40.16)(31.384,40.091)(31.522,40.021)
\qbezier(31.522,40.021)(31.66,39.952)(31.798,39.882)
\qbezier(31.798,39.882)(31.936,39.812)(32.074,39.742)
\qbezier(32.074,39.742)(32.212,39.672)(32.35,39.601)
\qbezier(32.35,39.601)(32.488,39.531)(32.626,39.46)
\qbezier(32.626,39.46)(32.764,39.389)(32.902,39.318)
\qbezier(32.902,39.318)(33.04,39.246)(33.178,39.175)
\qbezier(33.178,39.175)(33.316,39.103)(33.454,39.031)
\qbezier(33.454,39.031)(33.592,38.959)(33.73,38.886)
\qbezier(33.73,38.886)(33.868,38.813)(34.006,38.741)
\qbezier(34.006,38.741)(34.144,38.667)(34.281,38.594)
\qbezier(34.281,38.594)(34.419,38.52)(34.557,38.447)
\qbezier(34.557,38.447)(34.695,38.373)(34.833,38.298)
\qbezier(34.833,38.298)(34.971,38.224)(35.109,38.149)
\qbezier(35.109,38.149)(35.247,38.074)(35.385,37.999)
\qbezier(35.385,37.999)(35.523,37.923)(35.661,37.847)
\qbezier(35.661,37.847)(35.799,37.771)(35.937,37.695)
\qbezier(35.937,37.695)(36.075,37.618)(36.213,37.541)
\qbezier(36.213,37.541)(36.351,37.464)(36.489,37.386)
\qbezier(36.489,37.386)(36.627,37.308)(36.765,37.23)
\qbezier(36.765,37.23)(36.903,37.151)(37.041,37.073)
\qbezier(37.041,37.073)(37.179,36.993)(37.317,36.914)
\qbezier(37.317,36.914)(37.455,36.834)(37.593,36.754)
\qbezier(37.593,36.754)(37.731,36.673)(37.869,36.593)
\qbezier(37.869,36.593)(38.006,36.511)(38.144,36.43)
\qbezier(38.144,36.43)(38.282,36.347)(38.42,36.265)
\qbezier(38.42,36.265)(38.558,36.182)(38.696,36.099)
\qbezier(38.696,36.099)(38.834,36.015)(38.972,35.932)
\qbezier(38.972,35.932)(39.11,35.847)(39.248,35.762)
\qbezier(39.248,35.762)(39.386,35.677)(39.524,35.591)
\qbezier(39.524,35.591)(39.662,35.505)(39.8,35.418)
\qbezier(39.8,35.418)(39.938,35.331)(40.076,35.243)
\qbezier(40.076,35.243)(40.214,35.155)(40.352,35.066)
\qbezier(40.352,35.066)(40.49,34.977)(40.628,34.887)
\qbezier(40.628,34.887)(40.766,34.797)(40.904,34.706)
\qbezier(40.904,34.706)(41.042,34.615)(41.18,34.523)
\qbezier(41.18,34.523)(41.318,34.43)(41.456,34.337)
\qbezier(41.456,34.337)(41.594,34.243)(41.731,34.149)
\qbezier(41.731,34.149)(41.869,34.053)(42.007,33.958)
\qbezier(42.007,33.958)(42.145,33.861)(42.283,33.764)
\qbezier(42.283,33.764)(42.421,33.666)(42.559,33.568)
\qbezier(42.559,33.568)(42.697,33.468)(42.835,33.368)
\qbezier(42.835,33.368)(42.973,33.267)(43.111,33.166)
\qbezier(43.111,33.166)(43.249,33.063)(43.387,32.96)
\qbezier(43.387,32.96)(43.525,32.855)(43.663,32.75)
\qbezier(43.663,32.75)(43.801,32.644)(43.939,32.537)
\qbezier(43.939,32.537)(44.077,32.429)(44.215,32.32)
\qbezier(44.215,32.32)(44.353,32.21)(44.491,32.1)
\qbezier(44.491,32.1)(44.629,31.987)(44.767,31.874)
\qbezier(44.767,31.874)(44.905,31.76)(45.043,31.645)
\qbezier(45.043,31.645)(45.181,31.528)(45.319,31.41)
\qbezier(45.319,31.41)(45.456,31.291)(45.594,31.171)
\qbezier(45.594,31.171)(45.732,31.049)(45.87,30.926)
\qbezier(45.87,30.926)(46.008,30.801)(46.146,30.676)
\qbezier(46.146,30.676)(46.284,30.548)(46.422,30.42)
\qbezier(46.422,30.42)(46.56,30.289)(46.698,30.157)
\qbezier(46.698,30.157)(46.836,30.023)(46.974,29.888)
\qbezier(46.974,29.888)(47.112,29.75)(47.25,29.611)
\qbezier(47.25,29.611)(47.388,29.469)(47.526,29.327)
\qbezier(47.526,29.327)(47.664,29.181)(47.802,29.035)
\qbezier(47.802,29.035)(47.94,28.884)(48.078,28.734)
\qbezier(48.078,28.734)(48.216,28.578)(48.354,28.423)
\qbezier(48.354,28.423)(48.492,28.263)(48.63,28.103)
\qbezier(48.63,28.103)(48.768,27.937)(48.906,27.771)
\qbezier(48.906,27.771)(49.044,27.6)(49.181,27.428)
\qbezier(49.181,27.428)(49.319,27.251)(49.457,27.073)
\qbezier(49.457,27.073)(49.595,26.888)(49.733,26.704)
\qbezier(49.733,26.704)(49.871,26.512)(50.009,26.32)
\qbezier(50.009,26.32)(50.147,26.12)(50.285,25.921)
\qbezier(50.285,25.921)(50.423,25.712)(50.561,25.504)
\qbezier(50.561,25.504)(50.699,25.286)(50.837,25.068)
\qbezier(50.837,25.068)(50.975,24.84)(51.113,24.611)
\qbezier(51.113,24.611)(51.251,24.372)(51.389,24.132)
\qbezier(51.389,24.132)(51.527,23.88)(51.665,23.628)
\qbezier(51.665,23.628)(51.803,23.363)(51.941,23.097)
\qbezier(51.941,23.097)(52.079,22.816)(52.217,22.535)
\qbezier(52.217,22.535)(52.355,22.238)(52.493,21.94)
\qbezier(52.493,21.94)(52.631,21.624)(52.769,21.307)
\qbezier(52.769,21.307)(52.906,20.97)(53.044,20.632)
\qbezier(53.044,20.632)(53.182,20.271)(53.32,19.91)
\qbezier(53.32,19.91)(53.458,19.522)(53.596,19.135)
\qbezier(53.596,19.135)(53.734,18.717)(53.872,18.299)
\qbezier(53.872,18.299)(54.01,17.846)(54.148,17.394)
\qbezier(54.148,17.394)(54.286,16.902)(54.424,16.409)
\qbezier(54.424,16.409)(54.562,15.871)(54.7,15.333)
\qbezier(54.7,15.333)(54.838,14.742)(54.976,14.151)
\qbezier(54.976,14.151)(55.114,13.496)(55.252,12.842)
\qbezier(55.252,12.842)(55.39,12.114)(55.528,11.385)
\qbezier(55.528,11.385)(55.666,10.568)(55.804,9.75)
\qbezier(64.196,56.25)(64.334,55.432)(64.472,54.615)
\qbezier(64.472,54.615)(64.61,53.886)(64.748,53.158)
\qbezier(64.748,53.158)(64.886,52.504)(65.024,51.849)
\qbezier(65.024,51.849)(65.162,51.258)(65.3,50.667)
\qbezier(65.3,50.667)(65.438,50.129)(65.576,49.591)
\qbezier(65.576,49.591)(65.714,49.098)(65.852,48.606)
\qbezier(65.852,48.606)(65.99,48.154)(66.128,47.701)
\qbezier(66.128,47.701)(66.266,47.283)(66.404,46.865)
\qbezier(66.404,46.865)(66.542,46.478)(66.68,46.09)
\qbezier(66.68,46.09)(66.818,45.729)(66.956,45.368)
\qbezier(66.956,45.368)(67.094,45.03)(67.231,44.693)
\qbezier(67.231,44.693)(67.369,44.376)(67.507,44.06)
\qbezier(67.507,44.06)(67.645,43.762)(67.783,43.465)
\qbezier(67.783,43.465)(67.921,43.184)(68.059,42.903)
\qbezier(68.059,42.903)(68.197,42.637)(68.335,42.372)
\qbezier(68.335,42.372)(68.473,42.12)(68.611,41.868)
\qbezier(68.611,41.868)(68.749,41.628)(68.887,41.389)
\qbezier(68.887,41.389)(69.025,41.16)(69.163,40.932)
\qbezier(69.163,40.932)(69.301,40.714)(69.439,40.496)
\qbezier(69.439,40.496)(69.577,40.288)(69.715,40.079)
\qbezier(69.715,40.079)(69.853,39.88)(69.991,39.68)
\qbezier(69.991,39.68)(70.129,39.488)(70.267,39.296)
\qbezier(70.267,39.296)(70.405,39.112)(70.543,38.927)
\qbezier(70.543,38.927)(70.681,38.749)(70.819,38.572)
\qbezier(70.819,38.572)(70.956,38.4)(71.094,38.229)
\qbezier(71.094,38.229)(71.232,38.063)(71.37,37.897)
\qbezier(71.37,37.897)(71.508,37.737)(71.646,37.577)
\qbezier(71.646,37.577)(71.784,37.422)(71.922,37.266)
\qbezier(71.922,37.266)(72.06,37.116)(72.198,36.965)
\qbezier(72.198,36.965)(72.336,36.819)(72.474,36.673)
\qbezier(72.474,36.673)(72.612,36.531)(72.75,36.389)
\qbezier(72.75,36.389)(72.888,36.25)(73.026,36.112)
\qbezier(73.026,36.112)(73.164,35.977)(73.302,35.843)
\qbezier(73.302,35.843)(73.44,35.711)(73.578,35.58)
\qbezier(73.578,35.58)(73.716,35.452)(73.854,35.324)
\qbezier(73.854,35.324)(73.992,35.199)(74.13,35.074)
\qbezier(74.13,35.074)(74.268,34.951)(74.406,34.829)
\qbezier(74.406,34.829)(74.544,34.709)(74.681,34.59)
\qbezier(74.681,34.59)(74.819,34.472)(74.957,34.355)
\qbezier(74.957,34.355)(75.095,34.24)(75.233,34.126)
\qbezier(75.233,34.126)(75.371,34.013)(75.509,33.9)
\qbezier(75.509,33.9)(75.647,33.79)(75.785,33.68)
\qbezier(75.785,33.68)(75.923,33.571)(76.061,33.463)
\qbezier(76.061,33.463)(76.199,33.356)(76.337,33.25)
\qbezier(76.337,33.25)(76.475,33.145)(76.613,33.04)
\qbezier(76.613,33.04)(76.751,32.937)(76.889,32.834)
\qbezier(76.889,32.834)(77.027,32.733)(77.165,32.632)
\qbezier(77.165,32.632)(77.303,32.532)(77.441,32.432)
\qbezier(77.441,32.432)(77.579,32.334)(77.717,32.236)
\qbezier(77.717,32.236)(77.855,32.139)(77.993,32.042)
\qbezier(77.993,32.042)(78.131,31.947)(78.269,31.851)
\qbezier(78.269,31.851)(78.406,31.757)(78.544,31.663)
\qbezier(78.544,31.663)(78.682,31.57)(78.82,31.477)
\qbezier(78.82,31.477)(78.958,31.385)(79.096,31.294)
\qbezier(79.096,31.294)(79.234,31.203)(79.372,31.113)
\qbezier(79.372,31.113)(79.51,31.023)(79.648,30.934)
\qbezier(79.648,30.934)(79.786,30.845)(79.924,30.757)
\qbezier(79.924,30.757)(80.062,30.669)(80.2,30.582)
\qbezier(80.2,30.582)(80.338,30.495)(80.476,30.409)
\qbezier(80.476,30.409)(80.614,30.323)(80.752,30.238)
\qbezier(80.752,30.238)(80.89,30.153)(81.028,30.068)
\qbezier(81.028,30.068)(81.166,29.985)(81.304,29.901)
\qbezier(81.304,29.901)(81.442,29.818)(81.58,29.735)
\qbezier(81.58,29.735)(81.718,29.653)(81.856,29.57)
\qbezier(81.856,29.57)(81.994,29.489)(82.131,29.407)
\qbezier(82.131,29.407)(82.269,29.327)(82.407,29.246)
\qbezier(82.407,29.246)(82.545,29.166)(82.683,29.086)
\qbezier(82.683,29.086)(82.821,29.007)(82.959,28.927)
\qbezier(82.959,28.927)(83.097,28.849)(83.235,28.77)
\qbezier(83.235,28.77)(83.373,28.692)(83.511,28.614)
\qbezier(83.511,28.614)(83.649,28.536)(83.787,28.459)
\qbezier(83.787,28.459)(83.925,28.382)(84.063,28.305)
\qbezier(84.063,28.305)(84.201,28.229)(84.339,28.153)
\qbezier(84.339,28.153)(84.477,28.077)(84.615,28.001)
\qbezier(84.615,28.001)(84.753,27.926)(84.891,27.851)
\qbezier(84.891,27.851)(85.029,27.776)(85.167,27.702)
\qbezier(85.167,27.702)(85.305,27.627)(85.443,27.553)
\qbezier(85.443,27.553)(85.581,27.48)(85.719,27.406)
\qbezier(85.719,27.406)(85.856,27.333)(85.994,27.259)
\qbezier(85.994,27.259)(86.132,27.187)(86.27,27.114)
\qbezier(86.27,27.114)(86.408,27.041)(86.546,26.969)
\qbezier(86.546,26.969)(86.684,26.897)(86.822,26.825)
\qbezier(86.822,26.825)(86.96,26.754)(87.098,26.682)
\qbezier(87.098,26.682)(87.236,26.611)(87.374,26.54)
\qbezier(87.374,26.54)(87.512,26.469)(87.65,26.399)
\qbezier(87.65,26.399)(87.788,26.328)(87.926,26.258)
\qbezier(87.926,26.258)(88.064,26.188)(88.202,26.118)
\qbezier(88.202,26.118)(88.34,26.048)(88.478,25.979)
\qbezier(88.478,25.979)(88.616,25.909)(88.754,25.84)
\qbezier(88.754,25.84)(88.892,25.771)(89.03,25.702)
\qbezier(89.03,25.702)(89.168,25.634)(89.306,25.565)
\qbezier(89.306,25.565)(89.444,25.497)(89.581,25.428)
\qbezier(89.581,25.428)(89.719,25.36)(89.857,25.292)
\qbezier(89.857,25.292)(89.995,25.225)(90.133,25.157)
\qbezier(90.133,25.157)(90.271,25.089)(90.409,25.022)
\qbezier(90.409,25.022)(90.547,24.955)(90.685,24.888)
\qbezier(90.685,24.888)(90.823,24.821)(90.961,24.754)
\qbezier(90.961,24.754)(91.099,24.687)(91.237,24.621)
\qbezier(91.237,24.621)(91.375,24.554)(91.513,24.488)
\qbezier(91.513,24.488)(91.651,24.422)(91.789,24.356)
\qbezier(91.789,24.356)(91.927,24.29)(92.065,24.224)
\qbezier(92.065,24.224)(92.203,24.159)(92.341,24.093)
\qbezier(92.341,24.093)(92.479,24.028)(92.617,23.962)
\qbezier(92.617,23.962)(92.755,23.897)(92.893,23.832)
\qbezier(92.893,23.832)(93.031,23.767)(93.169,23.702)
\qbezier(93.169,23.702)(93.306,23.638)(93.444,23.573)
\qbezier(93.444,23.573)(93.582,23.508)(93.72,23.444)
\qbezier(93.72,23.444)(93.858,23.38)(93.996,23.315)
\qbezier(93.996,23.315)(94.134,23.251)(94.272,23.187)
\qbezier(94.272,23.187)(94.41,23.123)(94.548,23.06)
\qbezier(94.548,23.06)(94.686,22.996)(94.824,22.932)
\qbezier(94.824,22.932)(94.962,22.869)(95.1,22.805)
\qbezier(95.1,22.805)(95.238,22.742)(95.376,22.679)
\qbezier(95.376,22.679)(95.514,22.615)(95.652,22.552)
\qbezier(95.652,22.552)(95.79,22.489)(95.928,22.426)
\qbezier(95.928,22.426)(96.066,22.364)(96.204,22.301)
\qbezier(96.204,22.301)(96.342,22.238)(96.48,22.176)
\qbezier(96.48,22.176)(96.618,22.113)(96.756,22.051)
\qbezier(96.756,22.051)(96.894,21.988)(97.031,21.926)
\qbezier(97.031,21.926)(97.169,21.864)(97.307,21.802)
\qbezier(97.307,21.802)(97.445,21.74)(97.583,21.678)
\qbezier(97.583,21.678)(97.721,21.616)(97.859,21.554)
\qbezier(97.859,21.554)(97.997,21.493)(98.135,21.431)
\qbezier(98.135,21.431)(98.273,21.369)(98.411,21.308)
\qbezier(98.411,21.308)(98.549,21.246)(98.687,21.185)
\qbezier(98.687,21.185)(98.825,21.124)(98.963,21.062)
\qbezier(98.963,21.062)(99.101,21.001)(99.239,20.94)
\qbezier(99.239,20.94)(99.377,20.879)(99.515,20.818)
\qbezier(99.515,20.818)(99.653,20.757)(99.791,20.696)
\qbezier(99.791,20.696)(99.929,20.636)(100.067,20.575)
\qbezier(100.067,20.575)(100.205,20.514)(100.343,20.454)
\qbezier(100.343,20.454)(100.481,20.393)(100.619,20.333)
\qbezier(100.619,20.333)(100.756,20.272)(100.894,20.212)
\qbezier(100.894,20.212)(101.032,20.152)(101.17,20.091)
\qbezier(101.17,20.091)(101.308,20.031)(101.446,19.971)
\qbezier(101.446,19.971)(101.584,19.911)(101.722,19.851)
\qbezier(101.722,19.851)(101.86,19.791)(101.998,19.731)
\qbezier(101.998,19.731)(102.136,19.671)(102.274,19.611)
\qbezier(102.274,19.611)(102.412,19.552)(102.55,19.492)
\qbezier(102.55,19.492)(102.688,19.432)(102.826,19.373)
\qbezier(102.826,19.373)(102.964,19.313)(103.102,19.254)
\qbezier(103.102,19.254)(103.24,19.194)(103.378,19.135)
\qbezier(103.378,19.135)(103.516,19.075)(103.654,19.016)
\qbezier(103.654,19.016)(103.792,18.957)(103.93,18.898)
\qbezier(103.93,18.898)(104.068,18.838)(104.206,18.779)
\qbezier(104.206,18.779)(104.344,18.72)(104.481,18.661)
\qbezier(104.481,18.661)(104.619,18.602)(104.757,18.543)
\qbezier(104.757,18.543)(104.895,18.484)(105.033,18.426)
\qbezier(105.033,18.426)(105.171,18.367)(105.309,18.308)
\qbezier(105.309,18.308)(105.447,18.249)(105.585,18.191)
\qbezier(105.585,18.191)(105.723,18.132)(105.861,18.073)
\qbezier(105.861,18.073)(105.999,18.015)(106.137,17.956)
\qbezier(106.137,17.956)(106.275,17.898)(106.413,17.839)
\qbezier(106.413,17.839)(106.551,17.781)(106.689,17.723)
\qbezier(106.689,17.723)(106.827,17.664)(106.965,17.606)
\qbezier(106.965,17.606)(107.103,17.548)(107.241,17.49)
\qbezier(107.241,17.49)(107.379,17.432)(107.517,17.373)
\qbezier(107.517,17.373)(107.655,17.315)(107.793,17.257)
\qbezier(107.793,17.257)(107.931,17.199)(108.069,17.141)
\qbezier(108.069,17.141)(108.206,17.083)(108.344,17.026)
\qbezier(108.344,17.026)(108.482,16.968)(108.62,16.91)
\qbezier(108.62,16.91)(108.758,16.852)(108.896,16.794)
\qbezier(108.896,16.794)(109.034,16.737)(109.172,16.679)
\qbezier(109.172,16.679)(109.31,16.621)(109.448,16.564)
\qbezier(109.448,16.564)(109.586,16.506)(109.724,16.448)
\qbezier(109.724,16.448)(109.862,16.391)(110,16.333)
\color{black}
\end{picture}
\endgroup}
\newcommand{\rossSketchTwoPoles}{%
\begingroup
\setlength{\unitlength}{1mm}
\begin{picture}(120,65)
\linethickness{0.3pt}
\put(7,23.385){\vector(1,0){107}}
\put(28.182,5){\vector(0,1){57}}
\put(115,22.385){\makebox(0,0)[l]{$x$}}
\put(29.682,62){\makebox(0,0)[l]{$y$}}
\put(46.364,22.685){\line(0,1){1.4}}
\put(46.364,20.685){\makebox(0,0)[t]{\scriptsize $1$}}
\put(64.545,22.685){\line(0,1){1.4}}
\put(64.545,20.685){\makebox(0,0)[t]{\scriptsize $2$}}
\put(100.909,22.685){\line(0,1){1.4}}
\put(100.909,20.685){\makebox(0,0)[t]{\scriptsize $4$}}
\put(27.482,13.769){\line(1,0){1.4}}
\put(26.482,13.769){\makebox(0,0)[r]{\scriptsize $-5$}}
\put(27.482,42.615){\line(1,0){1.4}}
\put(26.482,42.615){\makebox(0,0)[r]{\scriptsize $10$}}
\multiput(46.364,8)(0,4){13}{\line(0,1){2}}
\multiput(64.545,8)(0,4){13}{\line(0,1){2}}
\color{mathjockgreen}
\linethickness{0.6pt}
\qbezier(10,23.278)(10.101,23.277)(10.202,23.276)
\qbezier(10.202,23.276)(10.302,23.276)(10.403,23.275)
\qbezier(10.403,23.275)(10.504,23.274)(10.605,23.274)
\qbezier(10.605,23.274)(10.705,23.273)(10.806,23.272)
\qbezier(10.806,23.272)(10.907,23.271)(11.008,23.271)
\qbezier(11.008,23.271)(11.108,23.27)(11.209,23.269)
\qbezier(11.209,23.269)(11.31,23.268)(11.411,23.267)
\qbezier(11.411,23.267)(11.511,23.267)(11.612,23.266)
\qbezier(11.612,23.266)(11.713,23.265)(11.814,23.264)
\qbezier(11.814,23.264)(11.914,23.263)(12.015,23.263)
\qbezier(12.015,23.263)(12.116,23.262)(12.217,23.261)
\qbezier(12.217,23.261)(12.317,23.26)(12.418,23.259)
\qbezier(12.418,23.259)(12.519,23.258)(12.62,23.258)
\qbezier(12.62,23.258)(12.72,23.257)(12.821,23.256)
\qbezier(12.821,23.256)(12.922,23.255)(13.023,23.254)
\qbezier(13.023,23.254)(13.123,23.253)(13.224,23.252)
\qbezier(13.224,23.252)(13.325,23.251)(13.426,23.25)
\qbezier(13.426,23.25)(13.527,23.249)(13.627,23.248)
\qbezier(13.627,23.248)(13.728,23.247)(13.829,23.246)
\qbezier(13.829,23.246)(13.93,23.246)(14.03,23.245)
\qbezier(14.03,23.245)(14.131,23.244)(14.232,23.243)
\qbezier(14.232,23.243)(14.333,23.241)(14.433,23.24)
\qbezier(14.433,23.24)(14.534,23.239)(14.635,23.238)
\qbezier(14.635,23.238)(14.736,23.237)(14.836,23.236)
\qbezier(14.836,23.236)(14.937,23.235)(15.038,23.234)
\qbezier(15.038,23.234)(15.139,23.233)(15.239,23.232)
\qbezier(15.239,23.232)(15.34,23.231)(15.441,23.23)
\qbezier(15.441,23.23)(15.542,23.228)(15.642,23.227)
\qbezier(15.642,23.227)(15.743,23.226)(15.844,23.225)
\qbezier(15.844,23.225)(15.945,23.224)(16.045,23.223)
\qbezier(16.045,23.223)(16.146,23.221)(16.247,23.22)
\qbezier(16.247,23.22)(16.348,23.219)(16.448,23.218)
\qbezier(16.448,23.218)(16.549,23.216)(16.65,23.215)
\qbezier(16.65,23.215)(16.751,23.214)(16.852,23.212)
\qbezier(16.852,23.212)(16.952,23.211)(17.053,23.21)
\qbezier(17.053,23.21)(17.154,23.208)(17.255,23.207)
\qbezier(17.255,23.207)(17.355,23.206)(17.456,23.204)
\qbezier(17.456,23.204)(17.557,23.203)(17.658,23.201)
\qbezier(17.658,23.201)(17.758,23.2)(17.859,23.199)
\qbezier(17.859,23.199)(17.96,23.197)(18.061,23.196)
\qbezier(18.061,23.196)(18.161,23.194)(18.262,23.193)
\qbezier(18.262,23.193)(18.363,23.191)(18.464,23.19)
\qbezier(18.464,23.19)(18.564,23.188)(18.665,23.186)
\qbezier(18.665,23.186)(18.766,23.185)(18.867,23.183)
\qbezier(18.867,23.183)(18.967,23.182)(19.068,23.18)
\qbezier(19.068,23.18)(19.169,23.178)(19.27,23.176)
\qbezier(19.27,23.176)(19.37,23.175)(19.471,23.173)
\qbezier(19.471,23.173)(19.572,23.171)(19.673,23.17)
\qbezier(19.673,23.17)(19.773,23.168)(19.874,23.166)
\qbezier(19.874,23.166)(19.975,23.164)(20.076,23.162)
\qbezier(20.076,23.162)(20.177,23.16)(20.277,23.159)
\qbezier(20.277,23.159)(20.378,23.157)(20.479,23.155)
\qbezier(20.479,23.155)(20.58,23.153)(20.68,23.151)
\qbezier(20.68,23.151)(20.781,23.149)(20.882,23.147)
\qbezier(20.882,23.147)(20.983,23.145)(21.083,23.143)
\qbezier(21.083,23.143)(21.184,23.14)(21.285,23.138)
\qbezier(21.285,23.138)(21.386,23.136)(21.486,23.134)
\qbezier(21.486,23.134)(21.587,23.132)(21.688,23.13)
\qbezier(21.688,23.13)(21.789,23.127)(21.889,23.125)
\qbezier(21.889,23.125)(21.99,23.123)(22.091,23.12)
\qbezier(22.091,23.12)(22.192,23.118)(22.292,23.116)
\qbezier(22.292,23.116)(22.393,23.113)(22.494,23.111)
\qbezier(22.494,23.111)(22.595,23.108)(22.695,23.106)
\qbezier(22.695,23.106)(22.796,23.103)(22.897,23.101)
\qbezier(22.897,23.101)(22.998,23.098)(23.098,23.095)
\qbezier(23.098,23.095)(23.199,23.093)(23.3,23.09)
\qbezier(23.3,23.09)(23.401,23.087)(23.502,23.084)
\qbezier(23.502,23.084)(23.602,23.082)(23.703,23.079)
\qbezier(23.703,23.079)(23.804,23.076)(23.905,23.073)
\qbezier(23.905,23.073)(24.005,23.07)(24.106,23.067)
\qbezier(24.106,23.067)(24.207,23.064)(24.308,23.061)
\qbezier(24.308,23.061)(24.408,23.058)(24.509,23.055)
\qbezier(24.509,23.055)(24.61,23.051)(24.711,23.048)
\qbezier(24.711,23.048)(24.811,23.045)(24.912,23.042)
\qbezier(24.912,23.042)(25.013,23.038)(25.114,23.035)
\qbezier(25.114,23.035)(25.214,23.031)(25.315,23.028)
\qbezier(25.315,23.028)(25.416,23.024)(25.517,23.021)
\qbezier(25.517,23.021)(25.617,23.017)(25.718,23.013)
\qbezier(25.718,23.013)(25.819,23.009)(25.92,23.006)
\qbezier(25.92,23.006)(26.02,23.002)(26.121,22.998)
\qbezier(26.121,22.998)(26.222,22.994)(26.323,22.99)
\qbezier(26.323,22.99)(26.423,22.986)(26.524,22.982)
\qbezier(26.524,22.982)(26.625,22.977)(26.726,22.973)
\qbezier(26.726,22.973)(26.827,22.969)(26.927,22.964)
\qbezier(26.927,22.964)(27.028,22.96)(27.129,22.955)
\qbezier(27.129,22.955)(27.23,22.951)(27.33,22.946)
\qbezier(27.33,22.946)(27.431,22.941)(27.532,22.937)
\qbezier(27.532,22.937)(27.633,22.932)(27.733,22.927)
\qbezier(27.733,22.927)(27.834,22.922)(27.935,22.917)
\qbezier(27.935,22.917)(28.036,22.911)(28.136,22.906)
\qbezier(28.136,22.906)(28.237,22.901)(28.338,22.895)
\qbezier(28.338,22.895)(28.439,22.89)(28.539,22.884)
\qbezier(28.539,22.884)(28.64,22.879)(28.741,22.873)
\qbezier(28.741,22.873)(28.842,22.867)(28.942,22.861)
\qbezier(28.942,22.861)(29.043,22.855)(29.144,22.849)
\qbezier(29.144,22.849)(29.245,22.843)(29.345,22.836)
\qbezier(29.345,22.836)(29.446,22.83)(29.547,22.823)
\qbezier(29.547,22.823)(29.648,22.817)(29.748,22.81)
\qbezier(29.748,22.81)(29.849,22.803)(29.95,22.796)
\qbezier(29.95,22.796)(30.051,22.789)(30.152,22.782)
\qbezier(30.152,22.782)(30.252,22.775)(30.353,22.767)
\qbezier(30.353,22.767)(30.454,22.759)(30.555,22.752)
\qbezier(30.555,22.752)(30.655,22.744)(30.756,22.736)
\qbezier(30.756,22.736)(30.857,22.728)(30.958,22.72)
\qbezier(30.958,22.72)(31.058,22.711)(31.159,22.703)
\qbezier(31.159,22.703)(31.26,22.694)(31.361,22.685)
\qbezier(31.361,22.685)(31.461,22.676)(31.562,22.667)
\qbezier(31.562,22.667)(31.663,22.657)(31.764,22.648)
\qbezier(31.764,22.648)(31.864,22.638)(31.965,22.628)
\qbezier(31.965,22.628)(32.066,22.618)(32.167,22.608)
\qbezier(32.167,22.608)(32.267,22.597)(32.368,22.587)
\qbezier(32.368,22.587)(32.469,22.576)(32.57,22.565)
\qbezier(32.57,22.565)(32.67,22.554)(32.771,22.542)
\qbezier(32.771,22.542)(32.872,22.531)(32.973,22.519)
\qbezier(32.973,22.519)(33.073,22.506)(33.174,22.494)
\qbezier(33.174,22.494)(33.275,22.481)(33.376,22.469)
\qbezier(33.376,22.469)(33.477,22.455)(33.577,22.442)
\qbezier(33.577,22.442)(33.678,22.428)(33.779,22.414)
\qbezier(33.779,22.414)(33.88,22.4)(33.98,22.385)
\qbezier(33.98,22.385)(34.081,22.37)(34.182,22.355)
\qbezier(34.182,22.355)(34.283,22.34)(34.383,22.324)
\qbezier(34.383,22.324)(34.484,22.308)(34.585,22.291)
\qbezier(34.585,22.291)(34.686,22.274)(34.786,22.257)
\qbezier(34.786,22.257)(34.887,22.239)(34.988,22.222)
\qbezier(34.988,22.222)(35.089,22.203)(35.189,22.184)
\qbezier(35.189,22.184)(35.29,22.165)(35.391,22.145)
\qbezier(35.391,22.145)(35.492,22.125)(35.592,22.104)
\qbezier(35.592,22.104)(35.693,22.083)(35.794,22.062)
\qbezier(35.794,22.062)(35.895,22.039)(35.995,22.017)
\qbezier(35.995,22.017)(36.096,21.993)(36.197,21.97)
\qbezier(36.197,21.97)(36.298,21.945)(36.398,21.921)
\qbezier(36.398,21.921)(36.499,21.895)(36.6,21.869)
\qbezier(36.6,21.869)(36.701,21.841)(36.802,21.814)
\qbezier(36.802,21.814)(36.902,21.786)(37.003,21.757)
\qbezier(37.003,21.757)(37.104,21.727)(37.205,21.696)
\qbezier(37.205,21.696)(37.305,21.665)(37.406,21.633)
\qbezier(37.406,21.633)(37.507,21.599)(37.608,21.565)
\qbezier(37.608,21.565)(37.708,21.53)(37.809,21.494)
\qbezier(37.809,21.494)(37.91,21.457)(38.011,21.419)
\qbezier(38.011,21.419)(38.111,21.38)(38.212,21.34)
\qbezier(38.212,21.34)(38.313,21.298)(38.414,21.255)
\qbezier(38.414,21.255)(38.514,21.211)(38.615,21.166)
\qbezier(38.615,21.166)(38.716,21.118)(38.817,21.071)
\qbezier(38.817,21.071)(38.917,21.02)(39.018,20.97)
\qbezier(39.018,20.97)(39.119,20.916)(39.22,20.862)
\qbezier(39.22,20.862)(39.32,20.805)(39.421,20.747)
\qbezier(39.421,20.747)(39.522,20.686)(39.623,20.624)
\qbezier(39.623,20.624)(39.723,20.558)(39.824,20.492)
\qbezier(39.824,20.492)(39.925,20.422)(40.026,20.351)
\qbezier(40.026,20.351)(40.127,20.275)(40.227,20.199)
\qbezier(40.227,20.199)(40.328,20.118)(40.429,20.036)
\qbezier(40.429,20.036)(40.53,19.948)(40.63,19.86)
\qbezier(40.63,19.86)(40.731,19.764)(40.832,19.669)
\qbezier(40.832,19.669)(40.933,19.565)(41.033,19.462)
\qbezier(41.033,19.462)(41.134,19.35)(41.235,19.237)
\qbezier(41.235,19.237)(41.336,19.114)(41.436,18.992)
\qbezier(41.436,18.992)(41.537,18.858)(41.638,18.724)
\qbezier(41.638,18.724)(41.739,18.576)(41.839,18.429)
\qbezier(41.839,18.429)(41.94,18.267)(42.041,18.105)
\qbezier(42.041,18.105)(42.142,17.926)(42.242,17.746)
\qbezier(42.242,17.746)(42.343,17.547)(42.444,17.348)
\qbezier(42.444,17.348)(42.545,17.125)(42.645,16.903)
\qbezier(42.645,16.903)(42.746,16.653)(42.847,16.404)
\qbezier(42.847,16.404)(42.948,16.122)(43.048,15.84)
\qbezier(43.048,15.84)(43.149,15.519)(43.25,15.199)
\qbezier(43.25,15.199)(43.351,14.831)(43.452,14.464)
\qbezier(43.452,14.464)(43.552,14.04)(43.653,13.615)
\qbezier(43.653,13.615)(43.754,13.12)(43.855,12.624)
\qbezier(43.855,12.624)(43.955,12.039)(44.056,11.453)
\qbezier(44.056,11.453)(44.157,10.752)(44.258,10.05)
\qbezier(44.258,10.05)(44.358,9.196)(44.459,8.342)
\qbezier(47.535,57.498)(47.58,56.45)(47.625,55.401)
\qbezier(47.625,55.401)(47.67,54.495)(47.715,53.589)
\qbezier(47.715,53.589)(47.76,52.798)(47.805,52.008)
\qbezier(47.805,52.008)(47.85,51.312)(47.895,50.617)
\qbezier(47.895,50.617)(47.94,50.001)(47.985,49.385)
\qbezier(47.985,49.385)(48.03,48.836)(48.075,48.287)
\qbezier(48.075,48.287)(48.12,47.795)(48.165,47.303)
\qbezier(48.165,47.303)(48.21,46.86)(48.255,46.417)
\qbezier(48.255,46.417)(48.3,46.017)(48.345,45.616)
\qbezier(48.345,45.616)(48.39,45.252)(48.435,44.888)
\qbezier(48.435,44.888)(48.48,44.556)(48.525,44.225)
\qbezier(48.525,44.225)(48.57,43.921)(48.615,43.618)
\qbezier(48.615,43.618)(48.66,43.34)(48.705,43.062)
\qbezier(48.705,43.062)(48.75,42.806)(48.795,42.55)
\qbezier(48.795,42.55)(48.84,42.315)(48.885,42.079)
\qbezier(48.885,42.079)(48.93,41.862)(48.975,41.644)
\qbezier(48.975,41.644)(49.02,41.443)(49.065,41.242)
\qbezier(49.065,41.242)(49.11,41.055)(49.155,40.868)
\qbezier(49.155,40.868)(49.2,40.695)(49.245,40.522)
\qbezier(49.245,40.522)(49.29,40.361)(49.335,40.2)
\qbezier(49.335,40.2)(49.38,40.05)(49.425,39.901)
\qbezier(49.425,39.901)(49.47,39.761)(49.515,39.621)
\qbezier(49.515,39.621)(49.56,39.491)(49.605,39.361)
\qbezier(49.605,39.361)(49.65,39.24)(49.695,39.119)
\qbezier(49.695,39.119)(49.74,39.005)(49.785,38.892)
\qbezier(49.785,38.892)(49.83,38.786)(49.875,38.681)
\qbezier(49.875,38.681)(49.92,38.582)(49.965,38.483)
\qbezier(49.965,38.483)(50.01,38.391)(50.055,38.298)
\qbezier(50.055,38.298)(50.1,38.212)(50.145,38.126)
\qbezier(50.145,38.126)(50.19,38.045)(50.235,37.965)
\qbezier(50.235,37.965)(50.28,37.89)(50.325,37.814)
\qbezier(50.325,37.814)(50.37,37.744)(50.415,37.674)
\qbezier(50.415,37.674)(50.46,37.609)(50.505,37.543)
\qbezier(50.505,37.543)(50.55,37.482)(50.595,37.421)
\qbezier(50.595,37.421)(50.64,37.365)(50.685,37.308)
\qbezier(50.685,37.308)(50.73,37.256)(50.775,37.203)
\qbezier(50.775,37.203)(50.82,37.154)(50.865,37.105)
\qbezier(50.865,37.105)(50.91,37.06)(50.955,37.015)
\qbezier(50.955,37.015)(51,36.974)(51.045,36.932)
\qbezier(51.045,36.932)(51.09,36.894)(51.135,36.855)
\qbezier(51.135,36.855)(51.18,36.82)(51.225,36.785)
\qbezier(51.225,36.785)(51.27,36.753)(51.315,36.721)
\qbezier(51.315,36.721)(51.36,36.692)(51.405,36.663)
\qbezier(51.405,36.663)(51.45,36.637)(51.495,36.611)
\qbezier(51.495,36.611)(51.54,36.587)(51.585,36.564)
\qbezier(51.585,36.564)(51.63,36.543)(51.675,36.522)
\qbezier(51.675,36.522)(51.72,36.504)(51.765,36.486)
\qbezier(51.765,36.486)(51.81,36.47)(51.855,36.455)
\qbezier(51.855,36.455)(51.9,36.442)(51.945,36.428)
\qbezier(51.945,36.428)(51.99,36.417)(52.035,36.407)
\qbezier(52.035,36.407)(52.08,36.398)(52.125,36.39)
\qbezier(52.125,36.39)(52.17,36.383)(52.215,36.377)
\qbezier(52.215,36.377)(52.26,36.373)(52.305,36.369)
\qbezier(52.305,36.369)(52.35,36.367)(52.395,36.366)
\qbezier(52.395,36.366)(52.44,36.366)(52.485,36.366)
\qbezier(52.485,36.366)(52.53,36.369)(52.575,36.371)
\qbezier(52.575,36.371)(52.62,36.376)(52.665,36.381)
\qbezier(52.665,36.381)(52.71,36.387)(52.755,36.394)
\qbezier(52.755,36.394)(52.8,36.403)(52.845,36.411)
\qbezier(52.845,36.411)(52.89,36.422)(52.935,36.433)
\qbezier(52.935,36.433)(52.98,36.446)(53.025,36.458)
\qbezier(53.025,36.458)(53.07,36.473)(53.115,36.488)
\qbezier(53.115,36.488)(53.16,36.505)(53.205,36.522)
\qbezier(53.205,36.522)(53.25,36.541)(53.295,36.559)
\qbezier(53.295,36.559)(53.34,36.58)(53.385,36.601)
\qbezier(53.385,36.601)(53.43,36.624)(53.475,36.647)
\qbezier(53.475,36.647)(53.52,36.672)(53.565,36.697)
\qbezier(53.565,36.697)(53.61,36.724)(53.655,36.751)
\qbezier(53.655,36.751)(53.7,36.78)(53.745,36.808)
\qbezier(53.745,36.808)(53.79,36.84)(53.835,36.871)
\qbezier(53.835,36.871)(53.88,36.904)(53.925,36.937)
\qbezier(53.925,36.937)(53.97,36.972)(54.015,37.007)
\qbezier(54.015,37.007)(54.06,37.045)(54.105,37.082)
\qbezier(54.105,37.082)(54.15,37.122)(54.195,37.161)
\qbezier(54.195,37.161)(54.24,37.203)(54.285,37.245)
\qbezier(54.285,37.245)(54.33,37.289)(54.375,37.332)
\qbezier(54.375,37.332)(54.42,37.379)(54.465,37.425)
\qbezier(54.465,37.425)(54.51,37.473)(54.555,37.522)
\qbezier(54.555,37.522)(54.6,37.573)(54.645,37.624)
\qbezier(54.645,37.624)(54.69,37.677)(54.735,37.73)
\qbezier(54.735,37.73)(54.78,37.786)(54.825,37.842)
\qbezier(54.825,37.842)(54.87,37.9)(54.915,37.958)
\qbezier(54.915,37.958)(54.96,38.019)(55.005,38.08)
\qbezier(55.005,38.08)(55.05,38.143)(55.095,38.206)
\qbezier(55.095,38.206)(55.14,38.273)(55.185,38.339)
\qbezier(55.185,38.339)(55.23,38.408)(55.275,38.476)
\qbezier(55.275,38.476)(55.32,38.548)(55.365,38.62)
\qbezier(55.365,38.62)(55.41,38.695)(55.455,38.769)
\qbezier(55.455,38.769)(55.5,38.847)(55.545,38.925)
\qbezier(55.545,38.925)(55.59,39.005)(55.635,39.086)
\qbezier(55.635,39.086)(55.68,39.17)(55.725,39.254)
\qbezier(55.725,39.254)(55.77,39.341)(55.815,39.429)
\qbezier(55.815,39.429)(55.86,39.519)(55.905,39.61)
\qbezier(55.905,39.61)(55.95,39.704)(55.995,39.799)
\qbezier(55.995,39.799)(56.04,39.897)(56.085,39.995)
\qbezier(56.085,39.995)(56.13,40.096)(56.175,40.198)
\qbezier(56.175,40.198)(56.22,40.304)(56.265,40.409)
\qbezier(56.265,40.409)(56.31,40.519)(56.355,40.629)
\qbezier(56.355,40.629)(56.4,40.743)(56.445,40.857)
\qbezier(56.445,40.857)(56.49,40.975)(56.535,41.093)
\qbezier(56.535,41.093)(56.58,41.216)(56.625,41.339)
\qbezier(56.625,41.339)(56.67,41.467)(56.715,41.594)
\qbezier(56.715,41.594)(56.76,41.727)(56.805,41.86)
\qbezier(56.805,41.86)(56.85,41.997)(56.895,42.135)
\qbezier(56.895,42.135)(56.94,42.278)(56.985,42.422)
\qbezier(56.985,42.422)(57.03,42.57)(57.075,42.719)
\qbezier(57.075,42.719)(57.12,42.874)(57.165,43.029)
\qbezier(57.165,43.029)(57.21,43.189)(57.255,43.35)
\qbezier(57.255,43.35)(57.3,43.517)(57.345,43.685)
\qbezier(57.345,43.685)(57.39,43.859)(57.435,44.032)
\qbezier(57.435,44.032)(57.48,44.213)(57.525,44.394)
\qbezier(57.525,44.394)(57.57,44.583)(57.615,44.771)
\qbezier(57.615,44.771)(57.66,44.967)(57.705,45.163)
\qbezier(57.705,45.163)(57.75,45.368)(57.795,45.572)
\qbezier(57.795,45.572)(57.84,45.785)(57.885,45.997)
\qbezier(57.885,45.997)(57.93,46.219)(57.975,46.441)
\qbezier(57.975,46.441)(58.02,46.672)(58.065,46.904)
\qbezier(58.065,46.904)(58.11,47.145)(58.155,47.386)
\qbezier(58.155,47.386)(58.2,47.638)(58.245,47.889)
\qbezier(58.245,47.889)(58.29,48.152)(58.335,48.415)
\qbezier(58.335,48.415)(58.38,48.69)(58.425,48.964)
\qbezier(58.425,48.964)(58.47,49.251)(58.515,49.538)
\qbezier(58.515,49.538)(58.56,49.839)(58.605,50.139)
\qbezier(58.605,50.139)(58.65,50.453)(58.695,50.767)
\qbezier(58.695,50.767)(58.74,51.096)(58.785,51.424)
\qbezier(58.785,51.424)(58.83,51.769)(58.875,52.113)
\qbezier(58.875,52.113)(58.92,52.474)(58.965,52.835)
\qbezier(58.965,52.835)(59.01,53.214)(59.055,53.593)
\qbezier(59.055,53.593)(59.1,53.991)(59.145,54.388)
\qbezier(59.145,54.388)(59.19,54.806)(59.235,55.224)
\qbezier(59.235,55.224)(59.28,55.663)(59.325,56.102)
\qbezier(59.325,56.102)(59.37,56.564)(59.415,57.026)
\qbezier(59.415,57.026)(59.46,57.513)(59.505,58)
\qbezier(68.669,53.866)(68.795,52.009)(68.921,50.153)
\qbezier(68.921,50.153)(69.047,48.602)(69.173,47.052)
\qbezier(69.173,47.052)(69.299,45.745)(69.425,44.438)
\qbezier(69.425,44.438)(69.551,43.326)(69.677,42.215)
\qbezier(69.677,42.215)(69.803,41.262)(69.929,40.31)
\qbezier(69.929,40.31)(70.055,39.488)(70.181,38.666)
\qbezier(70.181,38.666)(70.307,37.953)(70.433,37.24)
\qbezier(70.433,37.24)(70.559,36.616)(70.685,35.993)
\qbezier(70.685,35.993)(70.811,35.446)(70.937,34.899)
\qbezier(70.937,34.899)(71.063,34.417)(71.189,33.934)
\qbezier(71.189,33.934)(71.315,33.506)(71.441,33.079)
\qbezier(71.441,33.079)(71.567,32.698)(71.693,32.317)
\qbezier(71.693,32.317)(71.819,31.977)(71.945,31.637)
\qbezier(71.945,31.637)(72.071,31.332)(72.197,31.027)
\qbezier(72.197,31.027)(72.323,30.753)(72.449,30.478)
\qbezier(72.449,30.478)(72.575,30.231)(72.701,29.983)
\qbezier(72.701,29.983)(72.827,29.759)(72.953,29.534)
\qbezier(72.953,29.534)(73.079,29.331)(73.205,29.127)
\qbezier(73.205,29.127)(73.331,28.942)(73.457,28.757)
\qbezier(73.457,28.757)(73.583,28.587)(73.709,28.418)
\qbezier(73.709,28.418)(73.835,28.264)(73.961,28.109)
\qbezier(73.961,28.109)(74.087,27.967)(74.213,27.825)
\qbezier(74.213,27.825)(74.339,27.695)(74.465,27.565)
\qbezier(74.465,27.565)(74.591,27.445)(74.717,27.325)
\qbezier(74.717,27.325)(74.843,27.214)(74.969,27.103)
\qbezier(74.969,27.103)(75.095,27.001)(75.221,26.899)
\qbezier(75.221,26.899)(75.347,26.804)(75.473,26.71)
\qbezier(75.473,26.71)(75.599,26.622)(75.725,26.534)
\qbezier(75.725,26.534)(75.851,26.453)(75.977,26.371)
\qbezier(75.977,26.371)(76.103,26.295)(76.229,26.22)
\qbezier(76.229,26.22)(76.355,26.149)(76.481,26.079)
\qbezier(76.481,26.079)(76.607,26.013)(76.733,25.947)
\qbezier(76.733,25.947)(76.859,25.885)(76.985,25.824)
\qbezier(76.985,25.824)(77.111,25.766)(77.237,25.709)
\qbezier(77.237,25.709)(77.363,25.655)(77.489,25.601)
\qbezier(77.489,25.601)(77.615,25.551)(77.741,25.5)
\qbezier(77.741,25.5)(77.867,25.453)(77.993,25.405)
\qbezier(77.993,25.405)(78.119,25.361)(78.245,25.316)
\qbezier(78.245,25.316)(78.371,25.274)(78.497,25.232)
\qbezier(78.497,25.232)(78.623,25.193)(78.749,25.154)
\qbezier(78.749,25.154)(78.876,25.116)(79.002,25.079)
\qbezier(79.002,25.079)(79.128,25.044)(79.254,25.009)
\qbezier(79.254,25.009)(79.38,24.976)(79.506,24.943)
\qbezier(79.506,24.943)(79.632,24.912)(79.758,24.88)
\qbezier(79.758,24.88)(79.884,24.851)(80.01,24.821)
\qbezier(80.01,24.821)(80.136,24.793)(80.262,24.765)
\qbezier(80.262,24.765)(80.388,24.739)(80.514,24.712)
\qbezier(80.514,24.712)(80.64,24.687)(80.766,24.662)
\qbezier(80.766,24.662)(80.892,24.638)(81.018,24.614)
\qbezier(81.018,24.614)(81.144,24.591)(81.27,24.569)
\qbezier(81.27,24.569)(81.396,24.547)(81.522,24.525)
\qbezier(81.522,24.525)(81.648,24.505)(81.774,24.484)
\qbezier(81.774,24.484)(81.9,24.465)(82.026,24.445)
\qbezier(82.026,24.445)(82.152,24.427)(82.278,24.408)
\qbezier(82.278,24.408)(82.404,24.39)(82.53,24.373)
\qbezier(82.53,24.373)(82.656,24.356)(82.782,24.339)
\qbezier(82.782,24.339)(82.908,24.323)(83.034,24.307)
\qbezier(83.034,24.307)(83.16,24.291)(83.286,24.276)
\qbezier(83.286,24.276)(83.412,24.261)(83.538,24.247)
\qbezier(83.538,24.247)(83.664,24.233)(83.79,24.219)
\qbezier(83.79,24.219)(83.916,24.205)(84.042,24.192)
\qbezier(84.042,24.192)(84.168,24.179)(84.294,24.166)
\qbezier(84.294,24.166)(84.42,24.154)(84.546,24.141)
\qbezier(84.546,24.141)(84.672,24.13)(84.798,24.118)
\qbezier(84.798,24.118)(84.924,24.107)(85.05,24.095)
\qbezier(85.05,24.095)(85.176,24.084)(85.302,24.074)
\qbezier(85.302,24.074)(85.428,24.063)(85.554,24.053)
\qbezier(85.554,24.053)(85.68,24.043)(85.806,24.033)
\qbezier(85.806,24.033)(85.932,24.023)(86.058,24.014)
\qbezier(86.058,24.014)(86.184,24.005)(86.31,23.995)
\qbezier(86.31,23.995)(86.436,23.987)(86.562,23.978)
\qbezier(86.562,23.978)(86.688,23.969)(86.814,23.961)
\qbezier(86.814,23.961)(86.94,23.953)(87.066,23.945)
\qbezier(87.066,23.945)(87.192,23.937)(87.318,23.929)
\qbezier(87.318,23.929)(87.444,23.921)(87.57,23.914)
\qbezier(87.57,23.914)(87.696,23.906)(87.822,23.899)
\qbezier(87.822,23.899)(87.948,23.892)(88.074,23.885)
\qbezier(88.074,23.885)(88.2,23.878)(88.326,23.872)
\qbezier(88.326,23.872)(88.452,23.865)(88.578,23.859)
\qbezier(88.578,23.859)(88.704,23.852)(88.83,23.846)
\qbezier(88.83,23.846)(88.956,23.84)(89.082,23.834)
\qbezier(89.082,23.834)(89.208,23.828)(89.334,23.822)
\qbezier(89.334,23.822)(89.46,23.817)(89.586,23.811)
\qbezier(89.586,23.811)(89.712,23.806)(89.838,23.8)
\qbezier(89.838,23.8)(89.964,23.795)(90.09,23.79)
\qbezier(90.09,23.79)(90.216,23.785)(90.342,23.78)
\qbezier(90.342,23.78)(90.468,23.775)(90.594,23.77)
\qbezier(90.594,23.77)(90.72,23.765)(90.846,23.76)
\qbezier(90.846,23.76)(90.972,23.756)(91.098,23.751)
\qbezier(91.098,23.751)(91.224,23.747)(91.351,23.742)
\qbezier(91.351,23.742)(91.477,23.738)(91.603,23.734)
\qbezier(91.603,23.734)(91.729,23.729)(91.855,23.725)
\qbezier(91.855,23.725)(91.981,23.721)(92.107,23.717)
\qbezier(92.107,23.717)(92.233,23.713)(92.359,23.709)
\qbezier(92.359,23.709)(92.485,23.706)(92.611,23.702)
\qbezier(92.611,23.702)(92.737,23.698)(92.863,23.695)
\qbezier(92.863,23.695)(92.989,23.691)(93.115,23.688)
\qbezier(93.115,23.688)(93.241,23.684)(93.367,23.681)
\qbezier(93.367,23.681)(93.493,23.677)(93.619,23.674)
\qbezier(93.619,23.674)(93.745,23.671)(93.871,23.668)
\qbezier(93.871,23.668)(93.997,23.664)(94.123,23.661)
\qbezier(94.123,23.661)(94.249,23.658)(94.375,23.655)
\qbezier(94.375,23.655)(94.501,23.652)(94.627,23.649)
\qbezier(94.627,23.649)(94.753,23.646)(94.879,23.644)
\qbezier(94.879,23.644)(95.005,23.641)(95.131,23.638)
\qbezier(95.131,23.638)(95.257,23.635)(95.383,23.633)
\qbezier(95.383,23.633)(95.509,23.63)(95.635,23.627)
\qbezier(95.635,23.627)(95.761,23.625)(95.887,23.622)
\qbezier(95.887,23.622)(96.013,23.62)(96.139,23.617)
\qbezier(96.139,23.617)(96.265,23.615)(96.391,23.612)
\qbezier(96.391,23.612)(96.517,23.61)(96.643,23.608)
\qbezier(96.643,23.608)(96.769,23.605)(96.895,23.603)
\qbezier(96.895,23.603)(97.021,23.601)(97.147,23.599)
\qbezier(97.147,23.599)(97.273,23.597)(97.399,23.594)
\qbezier(97.399,23.594)(97.525,23.592)(97.651,23.59)
\qbezier(97.651,23.59)(97.777,23.588)(97.903,23.586)
\qbezier(97.903,23.586)(98.029,23.584)(98.155,23.582)
\qbezier(98.155,23.582)(98.281,23.58)(98.407,23.578)
\qbezier(98.407,23.578)(98.533,23.576)(98.659,23.575)
\qbezier(98.659,23.575)(98.785,23.573)(98.911,23.571)
\qbezier(98.911,23.571)(99.037,23.569)(99.163,23.567)
\qbezier(99.163,23.567)(99.289,23.566)(99.415,23.564)
\qbezier(99.415,23.564)(99.541,23.562)(99.667,23.56)
\qbezier(99.667,23.56)(99.793,23.559)(99.919,23.557)
\qbezier(99.919,23.557)(100.045,23.555)(100.171,23.554)
\qbezier(100.171,23.554)(100.297,23.552)(100.423,23.551)
\qbezier(100.423,23.551)(100.549,23.549)(100.675,23.548)
\qbezier(100.675,23.548)(100.801,23.546)(100.927,23.545)
\qbezier(100.927,23.545)(101.053,23.543)(101.179,23.542)
\qbezier(101.179,23.542)(101.305,23.54)(101.431,23.539)
\qbezier(101.431,23.539)(101.557,23.537)(101.683,23.536)
\qbezier(101.683,23.536)(101.809,23.535)(101.935,23.533)
\qbezier(101.935,23.533)(102.061,23.532)(102.187,23.531)
\qbezier(102.187,23.531)(102.313,23.529)(102.439,23.528)
\qbezier(102.439,23.528)(102.565,23.527)(102.691,23.526)
\qbezier(102.691,23.526)(102.817,23.524)(102.943,23.523)
\qbezier(102.943,23.523)(103.069,23.522)(103.195,23.521)
\qbezier(103.195,23.521)(103.321,23.52)(103.447,23.518)
\qbezier(103.447,23.518)(103.573,23.517)(103.699,23.516)
\qbezier(103.699,23.516)(103.826,23.515)(103.952,23.514)
\qbezier(103.952,23.514)(104.078,23.513)(104.204,23.512)
\qbezier(104.204,23.512)(104.33,23.511)(104.456,23.51)
\qbezier(104.456,23.51)(104.582,23.508)(104.708,23.507)
\qbezier(104.708,23.507)(104.834,23.506)(104.96,23.505)
\qbezier(104.96,23.505)(105.086,23.504)(105.212,23.503)
\qbezier(105.212,23.503)(105.338,23.502)(105.464,23.501)
\qbezier(105.464,23.501)(105.59,23.5)(105.716,23.5)
\qbezier(105.716,23.5)(105.842,23.499)(105.968,23.498)
\qbezier(105.968,23.498)(106.094,23.497)(106.22,23.496)
\qbezier(106.22,23.496)(106.346,23.495)(106.472,23.494)
\qbezier(106.472,23.494)(106.598,23.493)(106.724,23.492)
\qbezier(106.724,23.492)(106.85,23.491)(106.976,23.491)
\qbezier(106.976,23.491)(107.102,23.49)(107.228,23.489)
\qbezier(107.228,23.489)(107.354,23.488)(107.48,23.487)
\qbezier(107.48,23.487)(107.606,23.486)(107.732,23.486)
\qbezier(107.732,23.486)(107.858,23.485)(107.984,23.484)
\qbezier(107.984,23.484)(108.11,23.483)(108.236,23.482)
\qbezier(108.236,23.482)(108.362,23.482)(108.488,23.481)
\qbezier(108.488,23.481)(108.614,23.48)(108.74,23.479)
\qbezier(108.74,23.479)(108.866,23.479)(108.992,23.478)
\qbezier(108.992,23.478)(109.118,23.477)(109.244,23.477)
\qbezier(109.244,23.477)(109.37,23.476)(109.496,23.475)
\qbezier(109.496,23.475)(109.622,23.475)(109.748,23.474)
\qbezier(109.748,23.474)(109.874,23.473)(110,23.473)
\color{black}
\end{picture}
\endgroup}

The approach set is part of a limit's meaning. A two-sided limit at
$a$ uses nearby points different from $a$; a right-hand limit uses
points strictly greater than $a$. The value at $a$ is irrelevant to
these punctured limits. Limits may be finite or may equal one of
the two infinite symbols. ``Does not exist'' means that none of
these possibilities is a limit.

The sketches below are original schematic plots in finite windows.
Open circles mark missing points. Dashed lines mark asymptotes;
the axes themselves also mark an asymptote when appropriate.
No finite drawing can show all the oscillations accumulating at zero.
The formulas and estimates specify those oscillations exactly.

\exercise{20.1}{A sign graph and its five limiting behaviors}
\goal Sketch $f(x)=x/\abs{x}$ and inspect the five limits requested
in the exercise, in their stated order.

\step{Separate the positive and negative halves of the domain.}
The denominator is zero at $x=0$, so the domain is $\R\setminus\{0\}$.
For $x>0$, $\abs{x}=x$, whereas for $x<0$, $\abs{x}=-x$.
Consequently
\[
 f(x)=\begin{cases}1,&x>0,\\-1,&x<0.\end{cases}
\]
The graph consists of two horizontal half-lines, with the endpoints
at $(0,1)$ and $(0,-1)$ omitted.
\begin{center}\rossSketchSign\end{center}

\step{Read the limits, keeping the two sides separate.}
\[
\begin{array}{c|ccccc}
 \text{approach}&+\infty&0^+&0^-&-\infty&0\\ \hline
 \text{limit}&1&1&-1&-1&\text{does not exist}
\end{array}
\]
The right-hand values are constantly $1$ and the left-hand values
are constantly $-1$. Since these two one-sided limits differ,
Theorem~20.10 excludes a two-sided limit. Detailed verification
of the assertions is given in Exercise~20.5.
\answer{The five answers are $1,1,-1,-1$, and nonexistence.}

\exercise{20.2}{Two halves of a quadratic graph}
\goal Repeat the five-limit inspection for $f(x)=x^3/\abs{x}$.

\step{Cancel only where the denominator is nonzero.}
For $x>0$, $x^3/\abs{x}=x^2$. For $x<0$, division by $-x$
gives $x^3/\abs{x}=-x^2$. Thus
\[
 f(x)=\begin{cases}x^2,&x>0,\\-x^2,&x<0.\end{cases}
\]
The right branch rises to $+\infty$, the left branch falls to
$-\infty$ as $x$ runs toward the corresponding infinity, and both
branches approach the missing point $(0,0)$.
\begin{center}\rossSketchSignedSquare\end{center}
\[
\begin{array}{c|ccccc}
 \text{approach}&+\infty&0^+&0^-&-\infty&0\\ \hline
 \text{limit}&+\infty&0&0&-\infty&0
\end{array}
\]
The two-sided limit is $0$ even though $f(0)$ is undefined.
\answer{The five answers are $+\infty,0,0,-\infty,0$.
Exercise~20.6 proves them.}

\exercise{20.3}{The sine quotient and its removable hole}
\goal Sketch $f(x)=\sin x/x$, $x\ne0$, and inspect the five limits.

\step{Identify the hole, symmetry, and distant oscillations.}
Since $\sin(-x)=-\sin x$, division by $-x$ shows that
$f(-x)=f(x)$. The graph is even. Its zeros are the nonzero
integer multiples of $\pi$, and
\[
 \abs{f(x)}\le\frac1{\abs{x}}\qquad(x\ne0).
\]
The oscillations therefore shrink at either infinity. Near zero,
the curve approaches height $1$, with $(0,1)$ omitted.
\begin{center}\rossSketchSinc\end{center}
\[
\begin{array}{c|ccccc}
 \text{approach}&+\infty&0^+&0^-&-\infty&0\\ \hline
 \text{limit}&0&1&1&0&1
\end{array}
\]
\answer{The five limits are $0,1,1,0,1$.
Exercise~20.7 supplies the estimates and the proof near zero.}

\exercise{20.4}{Oscillations at zero and a horizontal asymptote}
\goal Repeat the inspection for $f(x)=x\sin(1/x)$, $x\ne0$.

\step{Describe both scales of the graph.}
The function is even, since
$(-x)\sin(-1/x)=x\sin(1/x)$. Its zeros are
$x=1/(k\pi)$ for nonzero integers $k$. Near zero they accumulate,
but
\[
 \abs{x\sin(1/x)}\le\abs{x},
\]
so the entire oscillating graph is squeezed into a narrowing pair
of straight lines. The origin is a hole. At either infinity put
$u=1/x$; then $x\sin(1/x)=\sin u/u$ approaches $1$.
\begin{center}\rossSketchOscillatingProduct\end{center}
The dashed horizontal line is $y=1$. Only finitely many of the
oscillations near the missing origin can be displayed.
\[
\begin{array}{c|ccccc}
 \text{approach}&+\infty&0^+&0^-&-\infty&0\\ \hline
 \text{limit}&1&0&0&1&0
\end{array}
\]
\answer{The five answers are $1,0,0,1,0$.
Exercise~20.8 proves both the zero and infinity assertions.}

\exercise{20.5}{Proof of the sign-function limit assertions}
\goal Prove every limit and nonexistence assertion in Exercise~20.1.

\step{The two positive-side limits.}
For $x>0$, $f(x)=1$, so $\abs{f(x)-1}=0<\varepsilon$
for every $\varepsilon>0$. For the right-hand limit at zero,
any $\delta>0$ works, because $0<x<\delta$ implies $x>0$.
For the limit at $+\infty$, the threshold $x>0$ works for every
$\varepsilon$. Thus these limits both equal $1$.

\step{The two negative-side limits.}
For $x<0$, $f(x)=-1$, so $\abs{f(x)-(-1)}=0$.
For the left-hand limit at zero, any $\delta>0$ works whenever
$-\delta<x<0$. For the limit at $-\infty$, all $x<0$ already
have the required value. Thus these two limits equal $-1$.

\step{Why zero has no two-sided limit of any kind.}
The one-sided limits are distinct finite numbers. By
Theorem~20.10 a two-sided limit, including an infinite limit,
would have to agree with both. It cannot. Equivalently, the
sequences $x_n=1/n$ and $y_n=-1/n$ both approach zero, but their
function values are constantly $1$ and $-1$, respectively.
\answer{All four one-sided or infinity limits hold, and the
two-sided limit at zero does not exist.}

\exercise{20.6}{Proof of the signed-quadratic limit assertions}
\goal Prove all five assertions in Exercise~20.2.

\step{The absolute value controls every approach to zero.}
For $x\ne0$, the piecewise formula gives $\abs{f(x)}=x^2$.
Given $\varepsilon>0$, take $\delta=\sqrt\varepsilon>0$.
Then
\[
 0<\abs{x}<\delta\quad\Longrightarrow\quad
 \abs{f(x)-0}=x^2<\delta^2=\varepsilon.
\]
This is the two-sided limit $0$. The same implication restricted
to $x>0$ or $x<0$ proves the corresponding one-sided limits.

\step{The two infinities have different signs.}
Fix $M>0$. If $x>\sqrt M$, then $x>0$ and
$f(x)=x^2>M$. Hence $f(x)\to+\infty$ as $x\to+\infty$.
If $x<-\sqrt M$, then $x<0$, $x^2>M$, and
$f(x)=-x^2<-M$. This is $f(x)\to-\infty$ as $x\to-\infty$.
\answer{The three limits at zero are $0$; the limits at positive
and negative infinity are $+\infty$ and $-\infty$, respectively.}

\exercise{20.7}{Proof of the sine-quotient limit assertions}
\goal Justify all five limits in Exercise~20.3.

\step{At either infinity, use a bound that does not depend on phase.}
For $x\ne0$, $\abs{\sin x/x}\le1/\abs{x}$.
Given $\varepsilon>0$, if $\abs{x}>1/\varepsilon$, then
\[
 \abs{\frac{\sin x}{x}-0}\le\frac1{\abs{x}}<\varepsilon.
\]
Taking either $x>1/\varepsilon$ or $x<-1/\varepsilon$ proves
the limits $0$ at $+\infty$ and $-\infty$.

\step{At zero, justify the familiar sine limit.}
We use the Mean Value Theorem, already permitted in Section~19,
and the familiar identities $(\sin t)'=\cos t$ and $\sin0=0$.
For every nonzero $x$, the theorem on the interval with endpoints
$0$ and $x$ gives a point $c_x$ strictly between these endpoints
such that
\[
 \frac{\sin x-\sin0}{x-0}=\cos c_x,
 \qquad \abs{c_x}<\abs{x}.
\]
The hypotheses hold because sine is continuous on that closed
interval and differentiable on its interior.

Let $\varepsilon>0$. Continuity of cosine at zero, and $\cos0=1$,
give $\eta>0$ such that $\abs{t}<\eta$ implies
$\abs{\cos t-1}<\varepsilon$. If $0<\abs{x}<\eta$, then
$\abs{c_x}<\eta$, so
\[
 \abs{\frac{\sin x}{x}-1}
 =\abs{\cos c_x-1}<\varepsilon.
\]
This proves the two-sided limit $1$. Restricting the same
estimate to positive or negative $x$ proves both one-sided limits.
It also verifies directly the continuous extension discussed in
Section~19, Example~9.
\answer{The two infinity limits are $0$ and all three limits
at zero equal $1$.}

\exercise{20.8}{Proof for the oscillating product}
\goal Prove the five assertions in Exercise~20.4.

\step{At zero, the small prefactor controls every oscillation.}
The bound $\abs{\sin t}\le1$ gives
\[
 \abs{x\sin(1/x)-0}\le\abs{x}\qquad(x\ne0).
\]
Given $\varepsilon>0$, choose $\delta=\varepsilon$.
Then $0<\abs{x}<\delta$ forces the displayed error to be less
than $\varepsilon$. Hence the two-sided limit is $0$, and so
are its two one-sided restrictions.

\step{At infinity, transform the expression rather than its factors separately.}
For $u=1/x$, one has $u\ne0$ and
\[
 x\sin(1/x)=\frac{\sin u}{u}.
\]
Exercise~20.7 proved that for every $\varepsilon>0$ there is
$\eta>0$ with
\[
 0<\abs{u}<\eta\quad\Longrightarrow\quad
 \abs{\frac{\sin u}{u}-1}<\varepsilon.
\]
If $\abs{x}>1/\eta$, then $0<\abs{1/x}<\eta$, so this bound
applies. It works for positive and negative $x$, proving a limit
of $1$ at both infinities. Multiplying the separate informal
limits $x\to\infty$ and $\sin(1/x)\to0$ would not be a valid
finite product-limit calculation; the quotient identity supplies
the needed control.
\answer{All three zero limits equal $0$ and both infinity
limits equal $1$.}

\exercise{20.9}{A reciprocal term and a slant asymptote}
\goal Sketch $f(x)=(1-x^2)/x$ and inspect the same five limits.

\step{Rewrite the function and identify the branches.}
The domain excludes zero, and for $x\ne0$,
\[
 f(x)=\frac1x-x.
\]
The graph is odd, since $f(-x)=-f(x)$. It crosses the $x$-axis
at $x=\pm1$. The line $x=0$ is a vertical asymptote, while
$y=-x$ is a slant asymptote because $f(x)-(-x)=1/x\to0$
at either infinity.
\begin{center}\rossSketchReciprocalMinusX\end{center}
The dotted diagonal is $y=-x$. The vertical axis is the other
asymptote; neither branch contains a point at $x=0$.
\[
\begin{array}{c|ccccc}
 \text{approach}&+\infty&0^+&0^-&-\infty&0\\ \hline
 \text{limit}&-\infty&+\infty&-\infty&+\infty&\text{does not exist}
\end{array}
\]
\answer{The answers are $-\infty,+\infty,-\infty,+\infty$,
and nonexistence. Exercise~20.10 proves them.}

\exercise{20.10}{Proof for the reciprocal-minus-linear function}
\goal Prove every limit assertion in Exercise~20.9.

\step{At positive infinity, the negative linear term dominates.}
Fix $M>0$, and choose $R=\max\{1,M+1\}$.
If $x>R$, then $1/x<1$ and $x>M+1$. Consequently
\[
 f(x)=\frac1x-x<1-x<-M.
\]
This proves the limit $-\infty$ at positive infinity.
For $x<-R$, put $t=-x>R$. Oddness gives
$f(x)=f(-t)=-f(t)>M$. Hence the limit at negative infinity
is $+\infty$.

\step{At zero, the reciprocal term dominates.}
For the same arbitrary $M>0$, choose
$\delta=\min\{1,1/(M+1)\}>0$.
If $0<x<\delta$, then $1/x>M+1$ and $x<1$, giving
\[
 f(x)=\frac1x-x>(M+1)-1=M.
\]
This proves the right-hand limit $+\infty$ at zero.
If $-\delta<x<0$, put $t=-x\in(0,\delta)$. Then
$f(x)=-f(t)<-M$, proving the left-hand limit $-\infty$.
The one-sided limits differ, so Theorem~20.10 shows that the
two-sided limit does not exist, even as an infinite limit.
\answer{All five assertions follow from explicit threshold
estimates, with the signs preserved on each side of zero.}

\exercise{20.11}{Three removable algebraic quotients}
\goal Find the three limits by simplifying only at points where
the original quotients are defined.

\step{(a) Factor the difference of squares.}
For $x\ne a$,
\[
 \frac{x^2-a^2}{x-a}
 =\frac{(x-a)(x+a)}{x-a}=x+a.
\]
The expression on the right is continuous at $a$. Since the
punctured expressions agree, their limits agree, and the answer
is $a+a=2a$. No value of the original quotient at $a$ is needed.

\step{(b) Multiply by the conjugate, with the domain checked.}
Here $b>0$. For example, $\abs{x-b}<b/2$ ensures $x>b/2>0$,
so both square roots exist. If also $x\ne b$, then
\[
 \frac{\sqrt x-\sqrt b}{x-b}
 =\frac{(\sqrt x-\sqrt b)(\sqrt x+\sqrt b)}
        {(x-b)(\sqrt x+\sqrt b)}
 =\frac1{\sqrt x+\sqrt b}.
\]
The conjugate is positive, and the cancellation uses $x-b\ne0$.
As $x\to b$, continuity of the square-root function gives a
denominator tending to $2\sqrt b>0$. The reciprocal limit is
therefore $1/(2\sqrt b)$.

\step{(c) Factor the difference of cubes.}
The identity
$x^3-a^3=(x-a)(x^2+ax+a^2)$ follows by multiplying out; the
two mixed terms cancel. For $x\ne a$ it gives
\[
 \frac{x^3-a^3}{x-a}=x^2+ax+a^2.
\]
This polynomial is continuous at $a$, so the limit is
$a^2+a^2+a^2=3a^2$.
\answer{The limits in (a), (b), and (c) are
$2a$, $1/(2\sqrt b)$, and $3a^2$, respectively.}

\exercise{20.12}{One simple pole and one squared pole}
\goal Sketch $f(x)=1/[(x-1)(x-2)^2]$, determine all four
requested one-sided limits, and decide the two two-sided limits.

\step{(a) Use signs and poles to organize the sketch.}
The domain is $\R\setminus\{1,2\}$.
For $x\ne2$, $(x-2)^2>0$, so the sign of $f(x)$ is the sign
of $x-1$. The graph is negative to the left of $1$ and positive
on both $(1,2)$ and $(2,\infty)$. There are vertical asymptotes
at $1$ and $2$ and a horizontal asymptote $y=0$ at both infinities.
For additional shape information, differentiation on each domain
interval gives
\[
 f'(x)=-\frac{3x-4}{(x-1)^2(x-2)^3}.
\]
Thus the middle positive branch decreases on $(1,4/3)$ and
increases on $(4/3,2)$, reaching $f(4/3)=27/4$.
The left branch decreases, and the branch to the right of $2$
decreases to zero. These derivative facts concern only the sketch;
the limit estimates below do not require them.
\begin{center}\rossSketchTwoPoles\end{center}
The dashed vertical lines are $x=1$ and $x=2$. The portions
outside the finite plotting window continue toward the indicated
infinite limits.

\step{(b), near $2$: the squared factor is positive on both sides.}
If $0<\abs{x-2}<1/2$, then $1/2<x-1<3/2$. Hence
\[
 f(x)=\frac1{(x-1)(x-2)^2}
 >\frac{2}{3\abs{x-2}^2}>0.
\]
Given $M>0$, choose
$\delta=\min\{1/2,\sqrt{2/(3M)}\}$.
Then $0<\abs{x-2}<\delta$ makes this last lower bound exceed
$M$. Thus both $\lim_{x\to2^+}f(x)$ and
$\lim_{x\to2^-}f(x)$ equal $+\infty$.

\step{(b), near $1$: the unsquared factor changes sign.}
If $0<\abs{x-1}<1/2$, then $1/2<\abs{x-2}<3/2$, so
\[
 \abs{f(x)}=\frac1{\abs{x-1}(x-2)^2}
 >\frac4{9\abs{x-1}}.
\]
For $M>0$, taking
$\delta=\min\{1/2,4/(9M)\}$ makes the magnitude exceed $M$.
For $x>1$ the function is positive, so the right-hand limit is
$+\infty$. For $x<1$ it is negative, so the left-hand limit is
$-\infty$.

\step{(c) Compare the two sides at each pole.}
At $2$ the two sides agree, giving
$\lim_{x\to2}f(x)=+\infty$ by Theorem~20.10.
At $1$ they disagree, so $\lim_{x\to1}f(x)$ does not exist.
\answer{In the requested order the one-sided limits are
$+\infty,+\infty,+\infty,-\infty$.
The two-sided limit at $2$ is $+\infty$; the one at $1$ does
not exist.}

\exercise{20.13}{Finite limit laws with reciprocal and square root}
\goal From $\lim_{x\to a}f(x)=3$ and
$\lim_{x\to a}g(x)=2$, prove all three asserted limits.

\step{(a) Apply the finite sum and product laws.}
The limits are finite, so Theorem~20.4 gives
\[
 \lim_{x\to a}g(x)^2=2\cdot2=4,
 \qquad\lim_{x\to a}3f(x)=3\cdot3=9.
\]
Adding these limits gives
\[
 \lim_{x\to a}[3f(x)+g(x)^2]=9+4=13.
\]

\step{(b) Ensure the reciprocal is defined before using its law.}
Since $g(x)\to2$, there is a punctured neighborhood of $a$
on which $\abs{g(x)-2}<1$. There $g(x)>1$, so $g(x)\ne0$.
On that neighborhood the quotient law applies, with the constant
numerator $1$ and denominator limit $2\ne0$. Therefore
\[
 \lim_{x\to a}\frac1{g(x)}=\frac12.
\]

\step{(c) Check the square-root domain near the limiting point.}
The sum and scalar laws first give
\[
 h(x)=3f(x)+8g(x)\longrightarrow3\cdot3+8\cdot2=25.
\]
For all sufficiently close $x\ne a$, one has
$\abs{h(x)-25}<1$, hence $h(x)>24>0$.
Thus the square root is defined throughout a punctured
neighborhood. The square-root function is continuous at $25$;
Theorem~20.5, applied on this neighborhood, gives
\[
 \lim_{x\to a}\sqrt{3f(x)+8g(x)}=\sqrt{25}=5.
\]
\answer{The three limits are $13$, $1/2$, and $5$.
Both nonzero-denominator and nonnegative-radicand conditions hold
on a sufficiently small punctured neighborhood.}

\exercise{20.14}{The two one-sided infinite reciprocal limits}
\goal Prove $\lim_{x\to0^+}1/x=+\infty$ and
$\lim_{x\to0^-}1/x=-\infty$ directly.

\step{Right side: choose the distance from the target height.}
Fix $M>0$ and put $\delta=1/M>0$.
If $0<x<\delta$, reciprocation of positive numbers reverses
the strict inequality and yields $1/x>1/\delta=M$.
This is precisely the threshold definition of the first limit.

\step{Left side: keep track of the negative sign.}
For the same $M$ and $\delta$, if $-\delta<x<0$, then
$0<-x<\delta$. The previous positive estimate gives
$1/(-x)>M$. Multiplying by $-1$ reverses the inequality:
$1/x<-M$. Since every negative threshold is $-M$ for some
$M>0$, this proves the second limit.
\answer{The reciprocal tends to $+\infty$ from the right and
to $-\infty$ from the left.}

\exercise{20.15}{Two limits of the shifted reciprocal cube}
\goal For the function in Example~4,
$f(x)=1/(x-2)^3$ with $x\ne2$, prove the limits at $-\infty$
and $2^+$ requested in the exercise.

\step{At negative infinity, bound the absolute error from zero.}
Let $\varepsilon>0$ and set $A=-\varepsilon^{-1/3}<0$.
If $x<A$, then $x<0$ and
$\abs{x-2}=2-x>-x>\varepsilon^{-1/3}$.
Cubing positive quantities and then reciprocating gives
\[
 \abs{f(x)-0}=\frac1{\abs{x-2}^3}<\varepsilon.
\]
This proves $\lim_{x\to-\infty}f(x)=0$.

\step{At the pole from the right, choose a small positive denominator.}
Given $M>0$, put $\delta=M^{-1/3}>0$.
If $2<x<2+\delta$, then $0<x-2<\delta$. Hence
$(x-2)^3<\delta^3=1/M$ and
$f(x)=1/(x-2)^3>M$.
Thus $\lim_{x\to2^+}f(x)=+\infty$.
\answer{The two limits are $0$ and $+\infty$, respectively.}

\exercise{20.16}{Order of right-hand limits and the loss of strictness}
\goal Prove the order assertion in part (a), and determine whether
strict pointwise order implies strict order of limits in part (b).
The existing limits may be finite or infinite.

\step{(a) Approach through one common sequence.}
Shrink $b>a$ if necessary so that the interval $(a,b)$ lies
in the neighborhoods used to define both right-hand limits.
Let
\[
 x_n=a+\frac{b-a}{n+1}.
\]
Then $a<x_n<b$ and $x_n\to a$. The hypotheses therefore give
\[
 f_1(x_n)\longrightarrow L_1,\qquad
 f_2(x_n)\longrightarrow L_2,\qquad
 f_1(x_n)\le f_2(x_n)\quad(n\in\N).
\]
The order rule for existing sequence limits, Exercise~9.9(c),
now gives $L_1\le L_2$, including the possible infinite values.
To recall why that order rule includes these values, if
$L_1>L_2$, choose a finite real $c$ strictly between them.
Convergence would eventually give both $f_1(x_n)>c$ and
$f_2(x_n)<c$, contrary to the pointwise inequality.

\step{(b) Give a counterexample with strictly ordered values.}
On $(a,a+1)$ put $f_1(x)=0$ and $f_2(x)=x-a$.
Every point satisfies $f_1(x)<f_2(x)$, since $x>a$.
Nevertheless both right-hand limits at $a$ equal $0$.
The positive gap $x-a$ shrinks to zero, so no strict gap remains
between the limits.
\answer{Part (a) gives $L_1\le L_2$. The answer to part (b)
is no: equality of limits is possible despite strict inequality
at every point.}

\exercise{20.17}{The squeeze lemma proves existence of the middle limit}
\goal Suppose $f_1(x)\le f_2(x)\le f_3(x)$ on $(a,b)$, and
both outside functions have right-hand limit $L$ at $a$.
Prove that $f_2$ has that same limit, without assuming it exists.

\step{Finite $L$: enclose all middle values in a prescribed error interval.}
Let $\varepsilon>0$. The two given limits supply
$\delta_1,\delta_3>0$ such that
\[
 \begin{split}
 a<x<a+\delta_1&\quad\Longrightarrow\quad
             L-\varepsilon<f_1(x)<L+\varepsilon,\\
 a<x<a+\delta_3&\quad\Longrightarrow\quad
             L-\varepsilon<f_3(x)<L+\varepsilon.
 \end{split}
\]
Take $\delta=\min\{\delta_1,\delta_3,b-a\}>0$.
Whenever $a<x<a+\delta$, the point lies in $(a,b)$ and
\[
 L-\varepsilon<f_1(x)\le f_2(x)\le f_3(x)<L+\varepsilon.
\]
Thus $\abs{f_2(x)-L}<\varepsilon$. This verifies the
definition of its limit and establishes existence as part
of the conclusion.

\step{Infinite $L$: one of the outside functions already forces the threshold.}
If $L=+\infty$, let $M>0$. The limit of $f_1$ gives
$\delta_1>0$ such that $a<x<a+\delta_1$ implies $f_1(x)>M$.
Restrict also to $x<b$. Then $f_2(x)\ge f_1(x)>M$, proving
$f_2(x)\to+\infty$.
If $L=-\infty$, let $M>0$ and use the upper function instead:
eventually $f_3(x)<-M$, so $f_2(x)\le f_3(x)<-M$.
This proves $f_2(x)\to-\infty$.
\answer{The middle limit exists and equals $L$ in every finite
or infinite case. No order theorem requiring its prior existence
was used.}

\exercise{20.18}{Rationalizing a square-root difference}
\goal Show the limit at zero exists for
\[
 f(x)=\frac{\sqrt{1+3x^2}-1}{x^2}\qquad(x\ne0),
\]
and determine its value.

\step{Check every factor used in the rationalization.}
For every real $x$, $1+3x^2\ge1$, so the square root exists
and $\sqrt{1+3x^2}+1\ge2>0$. If $x\ne0$, then $x^2>0$.
We may therefore multiply the numerator and denominator by the
conjugate and cancel the nonzero factor $x^2$:
\[
 \begin{aligned}
 f(x)&=\frac{(\sqrt{1+3x^2}-1)(\sqrt{1+3x^2}+1)}
             {x^2(\sqrt{1+3x^2}+1)}\\
 &=\frac{(1+3x^2)-1}{x^2(\sqrt{1+3x^2}+1)}
 =\frac3{\sqrt{1+3x^2}+1}.
 \end{aligned}
\]

\step{Use a continuous expression that also makes sense at zero.}
The polynomial $1+3x^2$ tends to $1$ as $x\to0$.
Continuity of square root gives a denominator tending to
$\sqrt1+1=2\ne0$. The reciprocal and scalar limit laws yield
\[
 \lim_{x\to0}f(x)=\frac32.
\]
This establishes existence, not just a possible value: the
original function agrees on every nonzero point with an
expression continuous at zero.
\answer{The limit exists and equals $3/2$.}

\exercise{20.19}{Changing the right-hand approach interval}
\goal Let $a<b_1<b_2$ and let $f$ be defined on $(a,b_2)$.
Show that existence and value of the limit along $(a,b_i)$
are independent of choosing $i=1$ or $i=2$.

\step{A limit along the larger interval restricts to the smaller one.}
Suppose the limit along $S_2=(a,b_2)$ exists and equals $L$.
Every sequence in $S_1=(a,b_1)$ is also a sequence in $S_2$,
since $S_1\subseteq S_2$. Thus whenever a sequence in $S_1$
tends to $a$, its function values tend to $L$ by the given
limit. This is the sequential definition of the limit along
$S_1$, with the same value.

\step{A sequence in the larger interval eventually enters the smaller one.}
Conversely, suppose the limit along $S_1$ exists and equals $L$.
Take an arbitrary sequence $(x_n)$ in $S_2$ tending to $a$.
The positive number $b_1-a$ is an allowable tolerance for
convergence. Therefore for some integer $N$,
\[
 n>N\quad\Longrightarrow\quad
 \abs{x_n-a}<b_1-a.
\]
Since all these $x_n>a$, the inequality says $a<x_n<b_1$.
Hence the tail $x_{N+1},x_{N+2},\ldots$ is a sequence in
$S_1$ and still tends to $a$. By the assumed smaller-interval
limit, the corresponding function values tend to $L$.
Adding back finitely many initial terms does not change a
finite or infinite sequence limit. Thus $f(x_n)\to L$ for the
whole arbitrary sequence, proving the limit along $S_2$.

Both approach sets admit sequences tending to $a$, for example
$a+(b_1-a)/(n+1)$, so their limits are meaningful under
Definition~20.1 and are unique. The proof did not require $L$
to be finite.
\answer{Either interval defines the same right-hand limit,
finite or infinite, which is denoted $\lim_{x\to a^+}f(x)$.}

\exercise{20.20}{Infinite sum and product limit rules}
\goal Assume $\lim_{x\to a^S}f_1(x)=+\infty$ and that
$L_2=\lim_{x\to a^S}f_2(x)$ exists. Prove all three asserted
rules and analyze the zero case in part (d).
Here $a$ may be finite or infinite, with the approach set $S$
as in Definition~20.1.

\step{Use an arbitrary admissible sequence, so all approach types are covered.}
Take any sequence $x_n\in S$ with $x_n\to a$, and write
$u_n=f_1(x_n)$, $v_n=f_2(x_n)$. The hypotheses give
$u_n\to+\infty$ and $v_n\to L_2$. We prove each conclusion
for this arbitrary sequence. Definition~20.1 will then give
the corresponding function limit.

\step{(a) A finite or positive-infinite second limit cannot cancel the first.}
Suppose $L_2\ne-\infty$. There is a finite real lower bound
$c$ for all sufficiently late $v_n$: if $L_2$ is finite, take
$c=L_2-1$; if $L_2=+\infty$, take $c=0$.
For any $M>0$, the limit $u_n\to+\infty$ eventually gives
$u_n>M-c$ (or a larger positive threshold if needed).
On the common tail,
\[
 u_n+v_n\ge u_n+c>M.
\]
Hence $u_n+v_n\to+\infty$, so
$\lim_{x\to a^S}(f_1+f_2)(x)=+\infty$.

\step{(b) A positive second limit gives a fixed positive factor.}
Suppose $0<L_2\le+\infty$.
If $L_2$ is finite, take $c=L_2/2>0$.
If $L_2=+\infty$, take $c=1$.
In either case $v_n\ge c$ eventually.
For any $M>0$, eventually $u_n>M/c>0$ as well. Multiplication
by this positive $u_n$ preserves the inequality, giving
\[
 u_nv_n\ge cu_n>M.
\]
Thus $u_nv_n\to+\infty$, proving the asserted product rule.

\step{(c) Reflect a negative second limit before applying the positive rule.}
Suppose $-\infty\le L_2<0$. Then $-v_n\to-L_2>0$,
where $-L_2=+\infty$ if $L_2=-\infty$.
Part (b)'s argument applied to $u_n$ and $-v_n$ gives
$u_n(-v_n)\to+\infty$.
Negating reverses every threshold inequality, so
$u_nv_n\to-\infty$. Since the admissible sequence was
arbitrary, this proves
$\lim_{x\to a^S}(f_1f_2)(x)=-\infty$.

\step{(d) A zero second limit does not determine the product.}
Use the common approach $x\to0^+$ on $S=(0,1)$, and fix
$f_1(x)=1/x^2$, which tends to $+\infty$.
For any prescribed real $c$, choose $f_2(x)=cx^2$.
Then $f_2(x)\to0$ but $f_1(x)f_2(x)=c$, so any finite real
limit can occur. If $f_2(x)=x$, its limit is again $0$ but
the product is $1/x\to+\infty$.
If $f_2(x)=-x$, the product tends to $-\infty$.

There need not be any product limit: choose
$f_2(x)=x^2\sin(1/x)$. The bound $\abs{f_2(x)}\le x^2$
gives $f_2(x)\to0$, while the product is $\sin(1/x)$.
Along
\[
 x_n=\frac1{2\pi n+\pi/2},\qquad
 y_n=\frac1{2\pi n+3\pi/2},
\]
both of which belong to $(0,1)$ and tend to zero, its values
are $1$ and $-1$, respectively. These incompatible sequence
limits exclude a function limit.
\answer{Parts (a), (b), and (c) give $+\infty$, $+\infty$,
and $-\infty$, respectively. When $L_2=0$, the product may
have any finite limit, either infinite limit, or no limit.}
